% =====================================================================
% 中文对照翻译版说明：
% 本文件为 openai/math 仓库 preprints/The-Exponential-Time-Hypothesis-is-False-October-25-2026
% 中 main_full_version.tex 的非官方中文翻译。公式、标签（\label）、
% 引用（\cite/\ref）与编号均与英文原文保持一致，以英文原文为准。
% =====================================================================
\documentclass[11pt,reqno]{amsart}
\usepackage[T1]{fontenc}
\PassOptionsToPackage{no-math}{fontspec}
\usepackage{xeCJK}
\setCJKmainfont[BoldFont=SimHei,ItalicFont=KaiTi]{SimSun}
\setCJKsansfont[BoldFont=SimHei]{SimHei}
\setCJKmonofont{FangSong}
\usepackage{amsmath,amssymb,amsthm,mathtools}
\usepackage{dsfont}
\usepackage[margin=1in]{geometry}
\usepackage{microtype}
\usepackage[colorlinks=true,allcolors=blue,backref=page]{hyperref}
\numberwithin{equation}{section}
\newtheorem{theorem}{定理}[section]
\newtheorem{lemma}[theorem]{引理}
\newtheorem{proposition}[theorem]{命题}
\newtheorem{corollary}[theorem]{推论}
\theoremstyle{definition}\newtheorem{definition}[theorem]{定义}
\theoremstyle{remark}\newtheorem{remark}[theorem]{注}
\newcommand{\poly}{\operatorname{poly}}
\newcommand{\Prob}{\mathbb P}
\DeclareMathOperator{\Bin}{Bin}
\allowdisplaybreaks[2]
\renewcommand{\le}{\leqslant}
\renewcommand{\ge}{\geqslant}
\setlength{\emergencystretch}{1.5em}
\setcounter{tocdepth}{1}
\title{时间指数猜想是错误的\\[0.6em]
 {\large\textmd{The Exponential Time Hypothesis is False}}\\[0.2em]
 {\small\textmd{（非官方中文对照翻译，以英文原文为准）}}}
\author{OpenAI}
\date{2026 年 10 月 17 日，2026 年 10 月 25 日修订}
\hypersetup{pdftitle={ETH is False},
 pdfauthor={OpenAI},
 pdfsubject={Exact counting, sparse parity tests, and polynomial interpolation}}
\begin{document}
\begin{abstract}
我们给出一个确定性算法，用于精确计数任意 $k$-CNF 公式的可满足赋值个数。
对每个固定的 $k\ge2$，在含 $n$ 个声明变量、$m$ 个子句的公式上，
我们的算法使用 $\poly_k(n,m)$ 空间，运行时间为
$(n+m)^{O_k(1)}2^{\,n/(\log n)^{\frac12+o(1)}}$。
我们的算法证伪了时间指数猜想（Exponential Time Hypothesis, ETH），
该猜想断言：判定 $n$ 个变量的 $3$-SAT 在最坏情形下需要
$2^{\Omega(n)}$ 时间。
\end{abstract}
\maketitle

\tableofcontents
\vspace{-\baselineskip}
\section{引言}
对每个固定的 $k\ge3$，$k$-SAT 是典型的 NP 完全问题~\cite{Cook71}，
其指数时间复杂度构成了时间指数猜想（ETH）与强时间指数猜想
（SETH）的基础~\cite{IP01,IPZ}。
本文证伪 ETH：对每个固定的 $k\ge2$，我们给出一个确定性算法，
在亚指数时间、多项式空间内计数 $k$-CNF 的可满足赋值。

将输入的、含 $n$ 个声明变量与 $m$ 个输入子句的 $k$-CNF 写作
\[
 F(x_1,\ldots,x_n)=\bigwedge_{e=1}^m
 \bigl[(x_{i_{e,1}}\oplus b_{e,1})\vee\cdots\vee
       (x_{i_{e,k}}\oplus b_{e,k})\bigr],
\]
其中 $i_{e,j}\in\{1,\ldots,n\}$ 与 $b_{e,j}\in\{0,1\}$ 分别是第
$e$ 个子句中第 $j$ 个文字的变量下标与取反位，$\oplus$ 表示模二加。
重复某个文字即允许子句的文字数少于 $k$。
下标 $i_{e,j}$ 与位 $b_{e,j}$ 由 $F$ 固定；
只有 $x_1,\ldots,x_n$ 在赋值范围内变化。
我们的目标是计算量
\begin{equation}\label{eq:solution-count}
 \#F=\bigl|\{x\in\{0,1\}^n:F(x)=1\}\bigr|.
\end{equation}
判定 $\#F>0$ 也就判定了可满足性。
当 $n$ 增大时我们保持 $k$ 固定。
全文中 $\log$ 表示自然对数。
\begin{theorem}\label{thm:main}
对每个固定的 $k\ge2$，存在确定性算法：给定含 $n\ge20$ 个声明变量、
$m$ 个子句的 $k$-CNF $F$，按 \eqref{eq:solution-count} 计算 $\#F$，
使用
\[
 (n+m)^{O_k(1)}2^{O_k(n\log\log n/\sqrt{\log n})}
\]
次位运算与 $(n+m)^{O_k(1)}$ 比特空间。
对每个固定的 $D>0$，在满足 $m\le Dn$ 的输入上，其运行时间为
\[
 (n+m)^{O_{k,D}(1)}2^{O_{k,D}(n/\sqrt{\log n})}.
\]
\end{theorem}

第~\ref{subsec:general-reduction} 节的有理插值论证给出较弱的指数
$O_k(\rho(n))$，其中
\[
 \rho(n)=\frac{n(\log\log n)^2}{\sqrt{\log n}}
          =\frac{n}{(\log n)^{\frac12+o(1)}}.
\]
附录~\ref{app:quantitative} 用张量收缩证明定理~\ref{thm:main}
的两个界。

当 $\#F>0$ 时，标准的自归约方法通过每次固定一个变量并检查剩余
计数是否为正来找出一个可满足赋值。若改为按两个条件计数来选取
每一位，则可得到完全均匀分布的随机可满足赋值，期望额外开销为
$\poly_k(n,m)$~\cite[Theorem~3.3(1)]{JVV86}。

该定理还给出任意精度的、亚指数时间的 Max-$k$-SAT 近似算法。
唯一游戏定理~\cite[Section~8.1]{OpenAIUGC26} 蕴含 Max-$k$-SAT
的 NP 难度界，其在逼近比上任意小误差范围内匹配~\cite{Rag08}
的半定规划算法给出的保证。
记 $\operatorname{OPT}(F)$ 为一个赋值同时满足的 $F$ 子句数的最大值。
\begin{corollary}\label{cor:max-sat-approximation}
对每个固定的 $k\ge2$ 与 $\varepsilon>0$，存在随机化算法：给定含
$n\ge20$ 个变量、$m$ 个子句的 $k$-CNF $F$，以概率
$1-2^{-\Omega(n)}$ 找到一个至少满足
$\operatorname{OPT}(F)-\varepsilon m$ 个子句的赋值，
时间为
\[
 (n+m)^{O_{k,\varepsilon}(1)}2^{O_{k,\varepsilon}(\rho(n))}
\]
，空间为 $(n+m)^{O_{k,\varepsilon}(1)}$。
\end{corollary}
\begin{proof}
可设 $m>0$。
设 $f(x)$ 为 $x$ 满足的子句比例，$\widehat f(x)$ 为从 $m$ 个子句中
独立均匀有放回抽取 $O(n/\varepsilon^2)$ 个组成的样本上的该比例。
由 Hoeffding 不等式与对所有 $2^n$ 个赋值的联合界，以概率
$1-2^{-\Omega(n)}$ 有
\[
 |\widehat f(x)-f(x)|\le\varepsilon/2
 \qquad\text{对每个 }x\in\{0,1\}^n.
\]
于是，若 $\widehat x$ 最大化 $\widehat f$，则
$f(\widehat x)\ge\max_x f(x)-\varepsilon$。
给定整数阈值 $b$，一个每门至多两个输入、规模
$O_{k,\varepsilon}(n)$ 的布尔电路可检验某赋值是否至少满足 $b$
个抽样子句。
标准的电路到 $3$-CNF 归约~\cite[Lemmas~15.9--15.10]{Bar23}
产出一个含 $O_{k,\varepsilon}(n)$ 个变量与子句的公式，其可满足
当且仅当该阈值可以达到。
定理~\ref{thm:main}、对阈值的二分搜索与自归约随即在所述界内找到
精确的样本最优者。
对每个样本，在 $O(n+\log(n/\varepsilon^2))$ 次尝试后截断拒绝采样，
即可保证最坏情形 $\poly(n,m,\varepsilon^{-1})$ 的采样时间，
且失败概率至多增加 $2^{-\Omega(n)}$。
\end{proof}

注意，形如 $2^{o(n)}$ 的运行时间界在一般多项式时间归约下不可传递
（传递需要 $2^{n^{o(1)}}$），因此我们的结果并不蕴含一般 NP 问题
都有亚指数时间算法。不过我们确实得到了计数图着色与完美匹配的
亚指数时间算法。

\begin{corollary}\label{cor:graph-counting}
在含 $v$ 个顶点、$e$ 条边的简单图上，独立集个数，以及每个固定的
$q\ge2$ 的正常 $q$-着色个数，均可确定性地在 $2^{o(v)}$ 时间、
$\poly_q(v,e)$ 空间内计算。完美匹配个数可确定性地在 $2^{o(v+e)}$
时间、$\poly(v,e)$ 空间内计算。
\end{corollary}
\begin{proof}
对 $G=(V,E)$，为每个 $u\in V$ 引入布尔变量 $x_u$。
\[
 \bigwedge_{\{u,w\}\in E}(\neg x_u\vee\neg x_w)
\]
的可满足赋值与独立集 $\{u\in V:x_u=1\}$ 一一对应。
对正常 $q$-着色，为每个（顶点，颜色）对使用一个布尔变量，
要求每个顶点恰取一种颜色，并禁止相邻顶点取相同颜色。
这给出一个保持计数的 $q$-CNF 编码，含 $qv$ 个变量与
$O_q(v+e)$ 个子句。

对完美匹配，为每条边使用一个布尔变量，并用一个线性规模电路
检验每个顶点恰与一条被选边关联。标准的电路到 $3$-CNF
归约~\cite[Lemmas~15.9--15.10]{Bar23} 给出 $O(v+e)$ 个变量与
子句。一旦边位固定，门变量取值唯一，因此该编码同样保持计数。
定理~\ref{thm:main} 给出所述时间与空间界。
\end{proof}

定理~\ref{thm:main} 的固定密度情形给出
$\mathrm E^{\mathrm{NP}}$（可由确定性时间 $2^{O(n)}$ 加 NP 谕示
判定的语言类）的电路下界。

\begin{corollary}\label{cor:quantitative-circuits}
设 $C_L(n)$ 是在 $\{0,1\}^n$ 上计算语言 $L$ 的电路所需非输入门的
最小数目，允许任意扇入至多二的布尔门与无限制扇出。存在语言
$L\in\mathrm E^{\mathrm{NP}}$ 与常数 $c>0$，使得
\begin{equation}\label{eq:quantitative-circuit-lower-bound}
 C_L(n)>c\,n\sqrt{\log n}\qquad\text{对无穷多个 }n.
\end{equation}
\end{corollary}

附录~\ref{subsec:circuit-lower-bounds} 证明该推论。
附录~\ref{subsec:local-inversion} 给出常数输出局部性的布尔映射的
求逆与精确原像计数算法。

\subsection{相关工作}\label{sec:related}
我们先讨论此前的 SAT 与计数算法，以及被定理~\ref{thm:main}
否定的若干复杂度假设。
随后将我们的方法与既有的代数算法比较，并指出证明中用到的既有构造，
它们的作用将在第~\ref{sec:technical-overview} 节的路线图中展开。

\subsubsection{SAT 算法与指数时间假设}\label{subsec:related-sat}
Sch\"oning 的局部搜索算法与 PPSZ 是 $k$-SAT 指数时间算法发展史上
的里程碑~\cite{Schon99,PPSZ05}。
Sch\"oning 的随机化算法从随机赋值出发，反复在某个未被满足的子句中
均匀选一个变量翻转。对每个固定的 $k\ge3$，其有界错误判定时间为
\[
 (n+m)^{O(1)}(2-2/k)^n
\]
~\cite[Section~2]{Schon99}。
特别地，对 $3$-SAT 给出 $(n+m)^{O(1)}(4/3)^n$ 时间。
PPSZ 算法~\cite{PPSZ05} 采用不同机制：先用有界消解导出更多子句，
再按随机顺序给变量赋值，能取被迫值时取被迫值，否则猜测。
对 $k=3,4$，其原始分析对唯一可满足公式给出比一般公式更强的界。
唯一 SAT 的保证后来被推广到一般 $3$-SAT 与 $4$-SAT；对 $3$-SAT
由此得到随机化 $O(1.308^n)$ 界~\cite{Hertli14}。
Hansen、Kaplan、Zamir 与 Zwick 通过给算法的猜测加偏置，改进了每个
固定 $k\ge3$ 的 PPSZ 界~\cite{HKZZ19}。
Scheder 通过锐化原始 PPSZ 算法的分析，在同一 $k$ 范围得到改进
~\cite[Theorem~1.5]{Scheder24}。
对一般 $3$-SAT，Jiang 与 Cai 的预印本给出随机化
$O^*(1.307031578^n)$ 界~\cite[Corollary~1.2]{JC26}，其中 $O^*$
略去输入规模的多项式因子。Liu 的确定性算法结合分支与去随机化
局部搜索，以 $O(1.32793^n)$ 时间求解 $3$-SAT~\cite[Theorem~21]{Liu18}。

随机化局部搜索也产生了有效的启发式算法。
WalkSAT 启发式选取一个未被满足的子句并翻转其一个变量，
把随机选择与贪心移动相结合~\cite{SKC94}。
在电路综合与诊断产生的公式上的实验表明了该方法的实用价值
~\cite[pp.~340--342, Tables~4--5]{SKC94}。

在均匀随机的 $k$-CNF 中，可满足性阈值处的实例曾被猜测是判定
可满足性最难的情形~\cite[Section~1]{VW21}；
对固定的 $k\ge3$ 与较低的子句密度 $m/n\le 2^k/k^C$（$C$
为普适常数），已知在公式上以高概率存在多项式时间的近均匀
可满足赋值采样~\cite[Theorem~1.2]{CLWYY25}。
对子句数为多项式的均匀随机公式，在以可满足为条件的情形下，
Vyas 与 Williams~\cite[Theorem~3]{VW21} 给出基于 PPZ 的算法，
在公式与算法随机选取二者之上以高概率在
$2^{n(1-\Omega(\log k/k))}$ 时间内找到一个解。

时间指数猜想（ETH）断言 $3$-SAT 最坏情形需要 $2^{\Omega(n)}$
确定性时间，从而对每个固定的 $k\ge3$，$k$-SAT 亦然
~\cite[Theorem~1]{IP01}。
强时间指数猜想（SETH）给出更尖锐的预言：随 $k$ 增大，$k$-SAT
需要 $2^{n(1-o_k(1))}$ 时间。
确切地说，对每个 $\varepsilon\in(0,1)$，存在固定的 $k$，
使任何确定性算法都不能在 $2^{(1-\varepsilon)n}(n+m)^{O(1)}$
时间内求解 $k$-SAT。
Sch\"oning 算法与 PPSZ 的运行时间界助长了这一预言：
它们相对穷举的改进随 $k$ 增大而衰减~\cite[Open Problems, item~2]{IP01}。

间隙时间指数猜想（Gap-ETH）强化 ETH：断言存在常数
$\delta,\gamma>0$，使得任何 $O(2^{\gamma n})$ 时间算法都无法区分
一个可满足的 $3$-CNF 与一个每个赋值都至少违反 $\delta$
比例子句的 $3$-CNF，即使 $m=O(n)$~\cite{Din16,MR17}。
其随机化版本排除以相同运行时间在每个被承诺输入上以至少 $2/3$
概率区分两种情形的算法~\cite[Conjecture~2.5]{CCK20}。
它被用于证明逼近 Label Cover~\cite{Din16} 与稠密约束满足问题
~\cite{MR17} 的难度，也排除了 Clique 与 Dominating Set 的非平凡
固定参数近似。对 Clique，被排除的比为 $o(\mathrm{OPT})$，
其中 $\mathrm{OPT}$ 为最大团规模~\cite[Corollaries~5.3 and~5.6]{CCK20}。

非确定性强时间指数猜想（NSETH）断言：对每个 $\varepsilon>0$，
某个固定的 $k$ 不存在非确定性 $2^{(1-\varepsilon)n}(n+m)^{O(1)}$
时间算法判定 $k$-CNF 是否\emph{不可满足}，等价地判定 $k$-DNF
是否为重言式~\cite{CGIMPS16}。
NSETH 被用于建立从 SAT 到 3SUM、全对最短路径等问题的确定性
细粒度归约的障碍~\cite{CGIMPS16}，并用于划定基于 SETH 的有向图
直径近似下界~\cite{Li21}。

定理~\ref{thm:main} 证伪 ETH、SETH、Gap-ETH 与 NSETH。
事实上，判定精确计数为正或为零即可对每个固定的 $k$ 以确定性
$2^{o(n)}(n+m)^{O_k(1)}$ 时间同时判定可满足与不可满足；
判定可满足性也解决了 Gap-ETH 的承诺问题。
该确定性算法还证伪了 ETH~\cite[Section~1]{DHMTW}、
SETH~\cite[Definition~2.1]{CDL16} 与 Gap-ETH 的有界错误随机化版本。
它与受限证明系统（如正则消解）的无条件指数下界不冲突
~\cite{BeckImp13}。
2016 年，Ryan Williams 通过给出一个以 $2^{n/2}\poly(n,m)$ 时间
证明 $k$-CNF 不可满足性的协议，证伪了 Merlin--Arthur 版 SETH
（MA-SETH）~\cite[Theorem~3.3]{Wil16}。

\subsubsection{精确计数与计数复杂度}\label{subsec:related-counting}
在 Valiant 的框架下，$\#\mathrm{SAT}$ 与 $\#3$-SAT 在吝啬归约下
是 $\#\mathrm P$ 完全的~\cite[Section~3, Facts~1--2]{Val79}。
对固定的 $k\ge3$，此前的精确计数算法达到
$\poly_k(n,m)2^{n-\Omega(n/k)}$ 时间，其中 Las Vegas 期望时间界见
~\cite[Lemma~3.5]{IMP12}，确定性界见~\cite[Theorem~1.6]{CW21}。
对 $\#2$-SAT，Wahlstr\"om 的算法以 $O^*(1.2377^n)$ 时间计数可满足
赋值~\cite{Wahlstrom08}。对 $\#3$-SAT，Kutzkov 用互补分支以
$O^*(1.637^n)$ 时间计数可满足赋值~\cite{Kutzkov26}。
两个算法都是确定性的。

计数时间指数猜想（$\#\mathrm{ETH}$）断言：精确计数 $3$-CNF 的
可满足赋值需要 $2^{\Omega(n)}$ 确定性时间~\cite[Section~1]{DHMTW}。
在 $\#\mathrm{ETH}$ 之下，Curticapean~\cite{Cur18} 排除了计数完美
匹配的确定性 $2^{o(v)}$ 时间算法，即使在含 $v$ 个顶点、$O(v)$
条边的无权二部图上亦然。
推论~\ref{cor:graph-counting} 的完美匹配算法特别适用于这些图。
我们的算法（定理~\ref{thm:main}）以 $2^{o(n)}\poly(n,m)$ 时间
确定性求解 $\#3$-SAT，因此证伪 $\#\mathrm{ETH}$。

\subsubsection{对细粒度复杂度的影响}\label{subsec:related-consequences}
ETH 与 SETH 成为参数化复杂度与细粒度复杂度的标准假设
~\cite{IPZ,VW19}。
基于 ETH 的归约给出了 Clique 对解规模的依赖~\cite{CHKX06}，
以及约束满足与分片子图同构对树宽的依赖~\cite{Marx10} 的下界。
SETH 支持 Hitting Set 与 Set Splitting~\cite{CDL16}，以及按树宽
参数化的 Independent Set 与 Dominating Set~\cite{LMS18} 的尖锐
指数界。
它还给出 Subset Sum 与 Bicriteria Path 的伪多项式算法可能加速的
界~\cite{ABHS19}。

对 $\mathrm P$ 中的问题，SETH 被用于排除编辑距离
~\cite{BackursIndyk18}、最长公共子序列与动态时间规整~\cite{BK15}、
Fr\'echet 距离~\cite{Bri14}、正则表达式匹配~\cite{BI16}
相对二次运行时间的多项式因子加速。
其他应用涉及图直径近似~\cite{RV13}，以及在图更新下维护可达计数
或最大强连通分量规模~\cite{AW14}。
这些结果共同使 SETH 成为解释众多迥异问题上算法表观最优性的
常用基础。

我们对 ETH 与 SETH 的证伪移除了上述难度断言的一个基础。
其中一些下界可由更弱的、可能仍然成立的假设推出。
例如，最长公共子序列与正则表达式匹配的下界可由正交向量猜想
（OVC）推出~\cite{BK15,BI16}。OVC 断言：对每个固定的
$\varepsilon>0$，不存在 $O(N^{2-\varepsilon}\poly(d))$ 时间算法
判定两个由 $N$ 个 $d$ 维布尔向量组成的集合是否含有正交对（各取
一个向量）。从 SAT 到正交向量的归约表明 SETH 蕴含
OVC~\cite{Wil05,VW19}。
Fr\'echet 距离、图直径近似、以及 Independent Set 与 Dominating
Set 对树宽的依赖的下界，也可由「无限制 CNF-SAT 对任何常数
$\delta>0$ 都不容许 $2^{(1-\delta)n}\poly(n,m)$ 时间算法」这一
假设推出~\cite{Bri14,RV13,LMS18}。定理~\ref{thm:main} 对每个固定
子句宽度 $k$ 给出运行时间保证，而该假设允许宽度随 $n$ 增长。

要用我们的 SAT 算法为目标问题得到更快算法，需要反方向的归约：
从该问题归约到固定宽度的 SAT，并控制变量数。
集合规模至多 $k$ 的 Set Splitting 具有此类归约
~\cite[Observation~3.6]{CDL16}：为每个元素使用一个布尔变量指定其
分侧。对每个集合，一个子句要求其变量至少一个为真，另一个子句
要求至少一个为假。所得 $k$-CNF 每个全集元素一个变量，因此对固定
$k$，定理~\ref{thm:main} 给出以全集大小计的亚指数时间（相差一个
输入规模的多项式因子）。

\subsubsection{组合问题的代数算法先行工作}\label{subsec:related-algebraic}
我们的算法是代数型的，基于对精心选取的辅助变量的多项式插值。
下面讨论用代数方法为组合问题设计算法的先行工作。
更一般的背景见 Koutis 与 Williams 关于代数指纹的综述
~\cite{KW16}、Williams 关于算法设计中多项式方法的综述
~\cite{Wil14}，以及 Husfeldt 关于算法容斥的综述~\cite{Hus11}。

Williams 的 split-and-list 算法~\cite{Wil05} 将 $2$-CSP 的变量
分为三组，枚举部分赋值，并用快速矩阵乘法合并。
Koivisto~\cite{Koi06} 给出直接的「和--积」表述。
第~\ref{sec:counting}--\ref{sec:hashes} 节的算法同样在组内枚举
赋值。固定共享测量 $\vec y$ 后，各组的贡献可以独立相乘。
Scott 与 Sorkin~\cite[Section~6]{SS09} 将 split-and-list 推广到
多项式值权与任意环。
我们的求和例程同样使用截断多项式环，在一次计算中保留多个系数。

Bj\"orklund~\cite[Sections~2--3]{Bjo16} 用容斥与多项式插值将
积和式与有向哈密顿圈计数归约为更小实例上的加权和。
他在表中把相同实例的权相加，然后对总权非零的每个不同实例求解一次。
Bj\"orklund、Kaski 与 Williams~\cite{BKW19} 用低次曲线与广义挂谷集
从小规模结构化点集上的取值恢复多项式的值，得到关于费米子式
（fermionant，包括元素绝对值被固定常数界住的整数矩阵的积和式）
的更快算法。
我们的多项式替代形式起类似作用：它们把目标计数与结构化取值
连接起来，从而可以重构。这里有效的结构是固定少数共享测量后
按组分离。

Bhargava、Ghosh、Kumar 与
Mohapatra~\cite[Sections~3.1--3.2]{BGKM23} 沿低次曲线在有限域上
求值多元多项式，在仅取值不够时使用 Hasse 导数与 Hermite 插值。
对一元限制 $g$，$t$ 处的 $j$ 阶 Hasse 导数是 $g(t+z)$ 中 $z^j$
的系数。我们使用相同的恢复原理：在每个允许快速求和的点上计算
若干 Taylor 系数，再用 Hermite 插值恢复原计数。

Woodruff 与 Yekhanin~\cite{WY07} 在私有信息检索中用沿直线的取值
与一阶导数恢复目标值。
Kopparty、Saraf 与 Yekhanin~\cite{KSY14} 用高阶 Hasse 导数做重数码
的局部校正，推广了基于直线的 Reed--Muller 码恢复。
在我们的算法中，矩阵构造提供低次替代形式，使我们能在固定少数
共享测量后，用快速求和计算插值点处的 Taylor 系数。

\subsubsection{矩阵构造的渊源}\label{subsec:related-spectral}
我们的计数算法在布尔赋值上求多项式乘积的和。这些多项式共享
变量，因此这些和不能分解为独立的小和之和。
我们引入（对插值参数多项式依赖的）线性替换，在固定少数整数线性
测量后使这种分解成为可能。一个辅助的有理方阵 $M_0$ 支撑这些
替换：$\operatorname{rank}(I-M_0)$ 的界限制测量个数，而 $M_0$
非零元上的限制保持替换次数很小。
第~\ref{subsec:preview-auxiliary} 节解释这一联系。

相关插值方法已被用于构造结构化矩阵的低秩近似。
在~\cite[Section~3]{AW17} 中，一个低次多项式在整数内积 $j$ 的
区间上匹配 $(-1)^j$，用的是~\cite[Lemma~3.1]{AW15} 的插值引理。
这给出在大多数位置与 Walsh--Hadamard 矩阵一致的低秩矩阵，
是矩阵刚性研究中的一个结果~\cite{Val77}。
相关插值在~\cite[Theorem~6.1]{AGP23} 中给出 Kronecker 幂的刚性界。
Dvir 与 Liu~\cite[Claim~2.20 and Lemmas~3.3--3.4]{DL20} 用谱方法：
对坐标置换不变函数的稀疏改动使许多 Fourier 系数为零，
降低其群循环矩阵的秩。

我们的构造同样在一个短区间上插值 $(-1)^j$，但此刻是为了得到
交换置换矩阵的多项式。同一个多项式既给出 $I-M_0$ 的秩界，
也给出 $M_0$ 非零元上的限制。我们在第~\ref{sec:construction} 节
直接证明这些性质及其定量界。

\subsubsection{稀疏化与奇偶检验}\label{subsec:related-reductions}
我们使用稀疏化引理~\cite{IPZ} 的保持计数版本~\cite[Appendix~A]{DHMTW}，
定量界取自~\cite[Section~6]{CIP} 与
\cite[Lemma~5.14 and equation~(5.4)]{CalabroThesis}。
这把可满足赋值分拆到同一组布尔变量上的稀疏公式中。
附录~\ref{app:sparsify} 记录了不交分支论证，以及其参数随输入
规模变化时所需的界。

初始的奇偶检验是 Tanner 码~\cite[Section~II]{Tanner81} 的约束，
作用在记录赋值违反了哪些子句的二进制向量上。
我们使用~\cite[Sections~IV and~VI]{SS96} 的基于边的展开码构造：
每条图边承载一位，每个顶点要求其关联的位属于某个固定的二元线性码。
局部码的最小距离与图展开共同迫使满足所有这些约束的每个非零
向量具有大支撑。该构造与距离论证是标准的。
引理~\ref{lem:sparse-checks} 给出此处所需的参数选取，包括检验的
个数与支撑大小的界，并用确定性穷举搜索构造图与局部码。
图是一个 Cayley 多重图的二部双覆盖。定义该多重图的移位几乎
平衡基群上的每个非常数奇偶性：这就是小偏置条件
~\cite{NN93,AR94}，在 \eqref{eq:hash-walsh} 中显式表述。
将未用的边位固定为零，得到恰在 $m$ 个坐标上的码。
另外的奇偶检验排除剩余的非零违反向量。我们用条件期望法背后的
平均原理~\cite{Raghavan88} 选取这些检验，用精确计数评估每个
候选。第~\ref{sec:hashes} 节给出此处所需的构造与定量界。

定理~\ref{thm:main} 提示人们重新寻找那些已知难度结果依赖于
ETH 或 SETH 的问题的算法。
第~\ref{subsec:related-consequences} 节的例子既显示了直接推论，
也显示了一些需要新算法思想的问题。

\subsection{记号}\label{subsec:notation}
对整数 $s\ge0$，记 $[s]=\{1,\ldots,s\}$，$[0]=\varnothing$。
对条件 $E$，标量 $\mathds{1}_{\{E\}}$ 在 $E$ 成立时为一，否则为零。
我们使用自然对数，底为二时写作 $\log_2$。
对复杂度界，$\poly(t_1,\ldots,t_s)$ 表示形如
$(1+t_1+\cdots+t_s)^{O(1)}$ 的界，其中隐含常数为绝对常数。
下标（如 $\poly_{k,c}(t_1,\ldots,t_s)$）允许这些常数依赖所标出的
固定参数。整数以带符号二进制记法存储，故整数 $z$ 使用
$O(\log(|z|+2))$ 比特。
有理数存储为整数分子与正整数分母；所述位长界对二者分别成立。
空间以比特计。算术运算次数与位运算次数区分使用。

对有限集 $S$，$\mathbb Q^S$ 中的向量以 $S$ 为坐标下标，
除另有说明外均为列向量。
对 $U\subseteq S$，指示向量 $\mathbf1_U\in\{0,1\}^S$ 的分量
$(\mathbf1_U)_v=\mathds{1}_{\{v\in U\}}$（$v\in S$）。
零向量记作 $\vec{0}$。对 $v\in S$，行 $\delta_v\in\mathbb Q^{1\times S}$
选取第 $v$ 个坐标，即 $\delta_vx=x_v$。
以 $U\subseteq S$ 为下标的向量可延拓到 $S$：在 $S\setminus U$ 的
每个坐标处补零。

我们写 $\vec{X}=(X_1,\ldots,X_n)$ 表示形式变量向量，用箭头表示
多项式系数向量，如 $\vec{P}_{f,v}(T)$。
对系数行 $\vec{p}=(p_i)_{i\in S}$ 与向量 $z=(z_i)_{i\in S}$，
二者的标量配对记作 $\langle\vec{p},z\rangle=\sum_{i\in S}p_i z_i$。
特别地，$\langle\vec{P}_{f,v}(T),\vec{X}\rangle$ 是变量 $X_i$
的标量线性型。

对变量 $\vec{X}$ 与变量 $z$ 的多项式 $P(\vec{X},z)$，
$[z^j]P$ 表示 $z^j$ 的系数（视为 $\vec{X}$ 的多项式），
$\deg_{\vec{X}} P$ 表示对 $\vec{X}$ 的总次数。
多项式行的次数是其各元素次数的最大值。
除非另有说明，多项式系数为有理数，矩阵核在 $\mathbb Q$ 上取。
\section{技术概览}\label{sec:technical-overview}
我们先对每个变量至多出现在三个子句中的 $3$-CNF 计数可满足赋值。
其子句多项式的乘法次数至多为 $3n$，因此可以用插值从小变量组上的
和恢复精确计数。
随后第~\ref{subsec:preview-general-sat} 节用不交稀疏化与变量复制，
对每个固定的 $k\ge2$，在 $(n+m)^{O_k(1)}2^{o(n)}$ 时间内计数任意
$k$-CNF 的可满足赋值。取 $k=3$ 即已证伪 ETH。
第~\ref{subsec:preview-full-case} 节解释在原变量上计数如何给出
命题~\ref{prop:interpolation-bound} 的指数 $O_k(\rho(n))$。

\subsection{受限出现次数的情形}\label{subsec:preview-setting}
设 $F$ 含 $n$ 个声明布尔变量。每个子句至多含三个文字，
每个变量（正出现与负出现合计）至多出现在三个子句中。
我们的任务是精确计数 $F$ 的可满足赋值。

我们删除重言子句与子句内重复的文字。空子句直接使计数为零。
这些操作保持出现次数界。对剩下的每个子句 $e$，令
$I_e\subseteq[n]$ 为其变量集。
对不出现在任何子句中的变量，添加一个 scope 为 $\{v\}$、取值恒为
一的因子。这样的因子不改变计数，并保证每个变量在构造中都有
 assigned 的位置。我们称这些恒一项为\emph{哑子句}（dummy clauses）。
令 $\mathcal F$ 为全部子句（含哑子句）的下标集。则
\begin{equation}\label{eq:overview-incidences}
 |I_e|\le3,\qquad
 \#\{e\in\mathcal F:v\in I_e\}\le3,\qquad
 n\le M:=\sum_{e\in\mathcal F}|I_e|\le3n.
\end{equation}
下文所有求和仍遍历原来的 $n$ 个布尔变量。

\subsection{计数的直接多项式}\label{subsec:preview-encoding}
对正文字写 $L_{e,v}(X_v)=X_v$，对负文字写 $L_{e,v}(X_v)=1-X_v$。
子句 $e$ 的满足多项式为
\[
 \Phi_e((X_v)_{v\in I_e})
 =1-\prod_{v\in I_e}(1-L_{e,v}(X_v)).
\]
例如，$x_a\vee\neg x_b\vee x_c$ 给出
\[
 \Phi_e(X_a,X_b,X_c)=1-(1-X_a)X_b(1-X_c).
\]
在布尔输入上，子句被满足时该多项式取一，否则取零。
对哑子句 $e$，置 $\Phi_e=1$。因此
\begin{equation}\label{eq:overview-direct-count}
 \#F=\sum_{x\in\{0,1\}^n}\prod_{e\in\mathcal F}
                   \Phi_e((x_v)_{v\in I_e}).
\end{equation}
该乘积的总次数至多为 $M\le3n$。
这个线性次数界正是三次出现假设带来的有用结论。

式~\eqref{eq:overview-direct-count} 是计数的精确表达式，
但对其 $2^n$ 个被加项求值就相当于穷举搜索。
为引出算法，假设给定变量到 $r$ 个组 $S_1,\ldots,S_r$
的一个划分，同时满足以下两个条件：
\begin{enumerate}
\item 每组至多含 $2n/r$ 个变量。
\item 每个子句的全部变量都属于同一组。
\end{enumerate}
这是一个假想的、限制性极强的假设；我们并不假定输入公式容许
这样的划分。

在该假设下，我们可以在每个组内独立计数可满足赋值。
设 $N_g$ 为 $S_g$ 中变量上满足该组全部子句的赋值个数。
这样的赋值可以与其他组的任意可满足赋值组合，因为没有子句
连接不同的组。每个全局可满足赋值恰以一种方式如此产生。
于是
\[
 \#F=\prod_{g=1}^r N_g.
\]
在组内穷举计算每个 $N_g$ 总共至多检查
\[
 \sum_{g=1}^r 2^{|S_g|}\le r\,2^{2n/r}
\]
个局部赋值，每个赋值上的工作量为多项式。
对大的 $r$，这远快于枚举全部 $2^n$ 个全局赋值。

我们将替换子句的输入，使不同组只通过一小批共享线性摘要
相互作用。固定这些摘要的取值后，便又可对每组分别求和。

\subsection{更容易求值点上的一族多项式}
\label{subsec:preview-deformation}
我们引入参数 $T$。在 $T=0$ 处，替换必须恢复原来的子句输入；
在选定的非零点 $T=1,\ldots,Q$ 处，替换应使我们能通过固定共享
线性摘要来分离各组。这里 $Q$ 是选定点的个数，
$R$ 将界定替代系数的次数。两个参数稍后选取。

写 $\vec{X}=(X_1,\ldots,X_n)$ 表示形式变量向量。
对每个出现 $v\in I_e$，用多项式系数向量
$\vec{P}_{e,v}(T)=(p_{e,v,1}(T),\ldots,p_{e,v,n}(T))$
把子句 $e$ 中的 $X_v$ 替换为标量内积
\[
 \langle \vec{P}_{e,v}(T),\vec{X}\rangle
 =\sum_{i=1}^n p_{e,v,i}(T)X_i.
\]
我们称此表达式为变量 $v$ 在子句 $e$ 中出现的
\emph{替代线性型}（replacement linear form）。它替换子句多项式
$\Phi_e$ 中的输入 $X_v$。
我们要求
\[
 \langle \vec{P}_{e,v}(0),\vec{X}\rangle=X_v,
 \qquad \deg \vec{P}_{e,v}\le R.
\]
不同子句可对同一变量使用不同的替代线性型。
定义
\begin{equation}\label{eq:simple-counting-polynomial}
 H(T)=\sum_{x\in\{0,1\}^n}\prod_{e\in\mathcal F}
       \Phi_e\bigl((\langle \vec{P}_{e,v}(T), x\rangle)_{v\in I_e}\bigr)
       \in\mathbb Q[T].
\end{equation}
在零点每个替换都恢复原变量，而每个子句多项式的次数至多增长
$R$ 倍。因此
\begin{equation}\label{eq:overview-value-degree}
 H(0)=\#F,\qquad \deg H\le MR\le3nR.
\end{equation}
我们将从 $t=1,\ldots,Q$ 处的 Taylor 系数恢复 $H(0)$。
次数界意味着每点只需 $O(1+nR/Q)$ 个系数，
第~\ref{subsec:preview-interpolation} 节将作解释。
因此我们的目标是让 $R$ 相对 $Q$ 较小：点多而替代次数低，
意味着每点需要的系数更少。

在非零的 $t$ 处，替代线性型的取值
$\langle\vec{P}_{e,v}(t),x\rangle$ 未必落在 $\{0,1\}$ 中，
即使 $x$ 是布尔取值。因此变换后的子句多项式
不必再充当满足指示子。

\subsection{通过固定测量分离求和}
\label{subsec:preview-separation}
固定 $t\in[Q]$。此时给定系数向量 $\vec{P}_{e,v}(t)$，
其分量已是常数。我们说明什么结构使我们可以快速计算 $H(t)$，
以及如何使用它。本小节中记 $S_g=S_{t,g}$、$C_g=C_{t,g}$。

\begin{definition}[$(r,\tau,D)$-design]
\label{def:overview-measurement-structure}
对给定的系数向量 $\vec{P}_{e,v}(t)$，
\emph{$t$ 处的 $(r,\tau,D)$-design}
由组 $S_g$、子句指派 $\sigma_e(t)$ 与测量矩阵 $C_g$
组成，满足以下四个条件：
\begin{enumerate}
\item \emph{小组。}将变量划分为两两不交的组
$[n]=S_1\sqcup\cdots\sqcup S_r$，且 $|S_g|\le2n/r$。
\item \emph{每子句一个指派组。}把每个子句 $e$ 指派到一个组
$\sigma_e(t)\in[r]$。条件~4 必须对子句 $e$ 的每个替代线性型
都在这同一组上成立。
\item \emph{少量有界测量。}每组 $g$ 有一个整数矩阵 $C_g$，
列以 $S_g$ 为下标，共 $c_g$ 行。每一行指定对 $S_g$
中变量的一个线性测量。总行数至多 $2\tau n$，且每行系数的
绝对值之和至多为 $D$。
对赋值 $x$，写
\[
 x_g=x|_{S_g},\qquad y_g=C_gx_g,\qquad
 \vec{y}=(y_1,\ldots,y_r).
\]
\item \emph{对其他组的依赖只通过测量。}设子句 $e$ 被指派到组
$g=\sigma_e(t)$。该子句的每个替代线性型可以自由依赖 $S_g$
中的变量；它对任何其他组 $S_h$ 的依赖必须是该组测量 $y_h$
的线性组合。确切地说，对每个 $v\in I_e$，必须存在系数向量
$\vec{\lambda}_{e,v,h}\in\mathbb Q^{c_h}$（$h\ne g$），使得
对每个 $x\in\mathbb Q^n$，
\begin{equation}\label{eq:outline-separation}
 \langle\vec{P}_{e,v}(t),x\rangle
 =\langle\vec{P}_{e,v}(t)|_{S_g},x_g\rangle
   +\sum_{h\ne g}\langle\vec{\lambda}_{e,v,h},y_h\rangle.
\end{equation}
等价地，$\vec{P}_{e,v}(t)|_{S_h}$ 对每个 $h\ne g$ 都落在 $C_h$
的行空间中。
\end{enumerate}
\end{definition}
矩阵 $C_h$ 为所有子句共享。子句的指派组不必包含其原变量。
不过一旦 $\vec{y}$ 固定，整个变换后的子句多项式就只依赖
其指派组中的变量。

\smallskip\noindent
\textbf{算法。}
给定 $t$ 处的 $(r,\tau,D)$-design，
先通过求解条件~4 保证的线性方程组找到系数向量
$\vec{\lambda}_{e,v,h}$。
对固定的 $\vec{y}$，令 $\mathcal X_{t,g}(y_g)$ 为测量值等于
$y_g$ 的 $S_g$ 上变量的赋值集合：
\[
 \mathcal X_{t,g}(y_g)
 =\{u\in\{0,1\}^{S_g}:C_gu=y_g\}.
\]
对这样的赋值 $u$，如下计算组 $g$ 的贡献。对指派到 $g$ 的每个
子句 $e$，先用 \eqref{eq:outline-separation} 从 $u$ 与固定的
测量计算其替代线性型的取值；再在这些取值上求子句多项式
$\Phi_e$ 的值；最后把指派到 $g$ 的全部子句的取值相乘。
记该乘积为 $A_{t,g}(u;\vec{y})$。此计算不需要其他组的赋值。

布尔赋值的每个测量都是 $[-D,D]$ 中的整数。
初始化 $Z=0$，执行以下步骤：
\begin{enumerate}
\item 尝试每个测量元组
$\vec{y}\in\mathcal Y:=\prod_{g=1}^r\{-D,\ldots,D\}^{c_g}$。
对每个元组执行第 2、3 步。
\item 对每个组 $g$，尝试其变量的每个布尔赋值 $u$，
仅当 $C_gu=y_g$ 时保留。把对应的乘积 $A_{t,g}(u;\vec{y})$
相加，得到组和
\[
 B_{t,g}(\vec{y})=
 \sum_{u\in\mathcal X_{t,g}(y_g)}A_{t,g}(u;\vec{y}).
\]
若无赋值满足 $C_gu=y_g$，置 $B_{t,g}(\vec{y})=0$。
\item 把各组之和相乘，并将 $\prod_{g=1}^r B_{t,g}(\vec{y})$
加到 $Z$ 上。
\end{enumerate}
处理完所有测量元组后，返回 $Z$。

为说明它返回 $H(t)$，固定 $\vec{y}$，并从每个集合
$\mathcal X_{t,g}(y_g)$ 中各选一个赋值。由于各组包含互不相交
的变量集，这些选择合成一个完整赋值。每个测量为 $\vec{y}$ 的
完整赋值恰以一种方式如此产生。其子句多项式的乘积就是各组
贡献 $A_{t,g}$ 的乘积。先把组和相乘、再对 $\vec{y}$ 求和，
即得
\begin{equation}\label{eq:outline-separated-sum}
 H(t)=\sum_{\vec{y}\in\mathcal Y}\prod_{g=1}^r
       \left(\sum_{u\in\mathcal X_{t,g}(y_g)}
                         A_{t,g}(u;\vec{y})\right).
\end{equation}
每个布尔赋值恰有一个测量元组，因此该和把每个赋值的贡献
恰好计一次。

测量元组至多有 $(2D+1)^{2\tau n}$ 个。对每个元组，算法对各组
内部至多尝试 $\sum_g2^{|S_g|}\le r2^{2n/r}$ 个赋值。
因此，不计多项式因子，运行时间为
\[
 (2D+1)^{2\tau n}2^{2n/r}.
\]
第一个因子界定测量元组的个数；第二个界定一组内的赋值个数。
算法一次扫描一个组，只保留当前测量元组、当前赋值以及运行中
的和与积；见引理~\ref{lem:measurement-separation}。

对第~\ref{subsec:preview-parameters} 节的参数选取，
我们可以构造次数至多 $R$ 的替代型，使其在零点恢复原变量，
并在每个 $t\in[Q]$ 处容许 $(r,\tau,D)$-design。
第~\ref{subsec:preview-auxiliary} 节给出构造。
我们先解释这些性质如何让我们高效地恢复计数。

\subsection{由 Taylor 系数恢复计数}
\label{subsec:preview-interpolation}
我们的目标是从 $t=1,\ldots,Q$ 处的信息恢复 $H(0)=\#F$。
由于 $\deg H\le3nR$，仅有 $Q$ 个值 $H(t)$ 未必确定 $H$。
我们改为在每个点计算若干 Taylor 系数。取
\[
 q=\left\lfloor\frac{3nR}{Q}\right\rfloor,\qquad J=q+1.
\]
假设对每个 $t\in[Q]$，我们知道 $H(t+z)$ 中 $z^0,\ldots,z^q$
的系数。这些数据唯一确定 $H$：若两个次数至多 $3nR$
的多项式有相同的数据，则其差对每个 $t\in[Q]$ 都被
$(T-t)^J$ 整除，从而被 $\prod_{t=1}^Q(T-t)^J$ 整除。
该乘积的次数为 $QJ>3nR$，故差必为零。
每个 Taylor 系数给出 $H$ 的系数的一个线性方程。
解这个有理线性方程组恢复 $H$，其常数项即所要的计数 $H(0)$。

剩下的是计算这些 Taylor 系数。固定 $t\in[Q]$；
我们要 $H(t+z)$ 中 $z$ 的次数至多 $q$ 的部分。
design 控制替代型在 $t$ 处的形态，但未必控制它们从 $t$
到 $t+z$ 的变化。我们先把这些变化单列为一个输入，
使固定该输入时组的分离得以保持。

\smallskip\noindent
\textbf{保持平移输入的分离。}
写
\[
 \Delta\vec P_{e,v}(z)=\vec P_{e,v}(t+z)-\vec P_{e,v}(t).
\]
保留 $x$ 作为被求和的布尔赋值，并为变化引入单独的形式向量
$\vec X$。定义
\begin{equation}\label{eq:outline-taylor-product}
 \mathcal G_x(\vec X,z)=
 \prod_{e\in\mathcal F}\Phi_e\bigl(
       (\langle\vec P_{e,v}(t),x\rangle+
        \langle\Delta\vec P_{e,v}(z),\vec X\rangle)_{v\in I_e}\bigr).
\end{equation}
这里 $x$ 是多项式的下标，而 $\vec X$ 与 $z$ 是其形式变量。
置 $\vec X=x$ 即恢复每个平移输入
$\langle\vec P_{e,v}(t+z),x\rangle$，因此
\[
 H(t+z)=\sum_{x\in\{0,1\}^n}\mathcal G_x(x,z).
\]
另一方面，对某个 $U\subseteq[n]$ 固定 $\vec X=\mathbf1_U$，
使每个修正项 $\langle\Delta\vec P_{e,v}(z),\mathbf1_U\rangle$
与 $x$ 无关。于是对 $x$ 的全部依赖都经
$\langle\vec P_{e,v}(t),x\rangle$ 传递，design 在此适用。
我们的任务因此是从 $\vec X$ 固定的和恢复 $\vec X=x$ 的和。

次数界使少量固定代入即可做到。
$\Delta\vec P_{e,v}(z)$ 的每个分量都被 $z$ 整除，故展开项中
$X_i$ 的每次出现至少贡献一个 $z$ 的因子。于是
\[
 \deg_{\vec X}([z^j]\mathcal G_x(\vec X,z))\le j.
\]
因此，丢弃 $q$ 以上的 $z$ 的幂后，每个单项至多涉及 $\vec X$
中 $q$ 个不同变量。

\smallskip\noindent
\textbf{用容斥恢复平移后的和。}
在模 $z^J$ 意义下工作，只保留 $z^q$ 以下的系数。
把 $\mathcal G_x$ 中 $X_i$ 的每个正幂替换为 $X_i$。
这保持其在布尔输入上的取值。由上述次数界，对布尔 $\vec X$，
\[
 \mathcal G_x(\vec X,z)\equiv
 \sum_{\substack{E\subseteq[n]\\|E|\le q}}
       a_E(x,z)\prod_{i\in E}X_i\pmod{z^J},
\]
其中 $a_E(x,z)$ 是该替换后所指单项的系数。
在 $\mathbf1_U$（恰在 $U$ 处取一的向量）处求值，即对以 $U$
的子集为下标的系数求和。容斥逐一提取：
\[
 a_E(x,z)\equiv
 \sum_{U\subseteq E}(-1)^{|E|-|U|}
                 \mathcal G_x(\mathbf1_U,z)\pmod{z^J}.
\]
当置 $\vec X=x$ 时，单项 $\prod_{i\in E}x_i$ 恰在 $E$ 中所有位
为一时等于一。它对 $x$ 上求和的贡献因此是
$\sum_{x:\,x_i=1\ (i\in E)}a_E(x,z)$（$x$ 取布尔值）。
代入上面的系数公式，这便化归为形式输入固定的和。

\smallskip\noindent
\textbf{计算受限和。}
对 $U\subseteq E\subseteq[n]$、$|E|\le q$，定义
\begin{equation}\label{eq:outline-restricted-taylor-sum}
 \mathcal Z_{t,E,U}(z)=
 \sum_{\substack{x\in\{0,1\}^n\\x_i=1\ (i\in E)}}
                    \mathcal G_x(\mathbf1_U,z)
       \quad\text{（在 }\mathbb Q[z]/(z^J)\text{ 中）}.
\end{equation}
这里 $U$ 固定形式输入 $\vec X=\mathbf1_U$，而 $E$ 把求和变量
$x$ 的某些位固定为一。
对单项贡献求和给出
\begin{equation}\label{eq:outline-taylor-recovery}
 H(t+z)\equiv
 \sum_{\substack{E\subseteq[n]\\|E|\le q}}
       \sum_{U\subseteq E}(-1)^{|E|-|U|}
                         \mathcal Z_{t,E,U}(z)\pmod{z^J}.
\end{equation}
这同时恢复全部 $J=q+1$ 个 Taylor 系数。
引理~\ref{lem:boolean-differences} 陈述了这里使用的一般恒等式。

每个和 $\mathcal Z_{t,E,U}$ 都可以用
第~\ref{subsec:preview-separation} 节的算法求值。
固定 $\vec X=\mathbf1_U$ 后，每个修正项
$\langle\Delta\vec P_{e,v}(z),\mathbf1_U\rangle$ 与 $x$ 无关。
对 $x$ 的全部依赖都来自型
$\langle\vec P_{e,v}(t),x\rangle$，因此固定 $\vec y$
仍然分离各子句多项式。限制 $x_i=1$ 是各组局部的。
分离恒等式只用加法、乘法与分配律，因此也适用于这些截断
多项式值。每个子句都在代入 $\vec X=\mathbf1_U$ 后求值，
只用次数小于 $J$ 的一元 $z$ 多项式，且每次乘法后丢弃更高次幂。
多元展开及其系数 $a_E(x,z)$ 只用于论证恒等式；算法从不
构造它们。

在 $t=1,\ldots,Q$ 处重复这一计算，即提供上述插值所需的全部
Taylor 系数，从而恢复 $H(0)=\#F$。
\subsection{三项成本与运行时间}
\label{subsec:preview-parameters}
先按下不表：第~\ref{subsec:preview-auxiliary} 节将构造替代型，
使其在零点恢复原变量、次数至多 $R$，并在每个 $t\in[Q]$
处容许 $(r,\tau,D)$-design。对足够大的 $n$，构造耗时多项式，
并达到以下参数界：
\begin{equation}\label{eq:overview-parameter-bounds}
 \begin{aligned}
 r&=\Theta\!\left(\frac{\sqrt{\log n}}{(\log\log n)^2}\right),
 &Q&=\Theta\!\left(\frac{\log n}{\log\log n}\right),\\
 R&=\Theta(\sqrt{\log n}),
 &\tau&=\Theta\!\left(\frac1{r\log n}\right),\qquad
 D=r^{2Q}\le(2n)^{1/128}.
 \end{aligned}
\end{equation}
我们先说明这些界如何给出快速计数算法；
下一小节将论证它们可以同时达成。

回忆 $q=J-1=\lfloor3nR/Q\rfloor$。对每个求值点 $t$，
第~\ref{subsec:preview-interpolation} 节对每一对
$U\subseteq E$（$|E|\le q$）计算一个受限和。
大小为 $j$ 的 $E$ 有 $\binom nj$ 种取法，随后 $U$ 有 $2^j$
种取法。每个受限和枚举测量元组与各组内部的赋值。
因此，不计多项式因子，三项成本相乘给出
\begin{equation}\label{eq:overview-three-costs}
 \underbrace{\sum_{j=0}^{\min(q,n)}2^j\binom nj}_{\text{Taylor 数据}}
 \ \cdot\
 \underbrace{(2D+1)^{2\tau n}}_{\text{测量}}
 \ \cdot\
 \underbrace{2^{2n/r}}_{\text{单组}}.
\end{equation}
为界定第一项成本，注意对足够大的 $n$ 有 $1\le q\le n/2$
（因 $R/Q\to0$）。用 $\binom nj\le(en/j)^j$ 以及
$(2en/j)^j$ 在 $1\le j\le n$ 上递增的事实，得
\[
 \sum_{j=0}^q2^j\binom nj
 \le(q+1)\left(\frac{2en}{q}\right)^q.
\]
由于 $q=\Theta(nR/Q)$，该项成本的对数为
$O(n(R/Q)\log(2Q/R))$。对 \eqref{eq:overview-three-costs}
中三项成本取对数，即得
\[
 O\!\left(n\left[
       \frac RQ\log\frac{2Q}{R}
       +\tau\log(2D+1)+\frac1r\right]\right).
\]
在 \eqref{eq:overview-parameter-bounds} 下每一项都很小：
\[
 \begin{aligned}
 \frac RQ\log\frac{2Q}{R}
   &=O\!\left(\frac{\log\log n}{\sqrt{\log n}}
                      \cdot\log\log n\right),\\
 \tau\log(2D+1)&=O\!\left(\frac1{r\log n}\cdot\log n\right)
                  =O(1/r),\\
 \frac1r&=O\!\left(\frac{(\log\log n)^2}{\sqrt{\log n}}\right).
 \end{aligned}
\]
于是受限出现次数的实例可在
\begin{equation}\label{eq:overview-bounded-time}
 n^{O(1)}2^{O(n(\log\log n)^2/\sqrt{\log n})}
\end{equation}
次位运算与多项式空间内完成计数。有限多个较小的输入用穷举
处理。

点数 $Q$、组数 $r$ 与截断阶 $J=O(1+nR/Q)$ 都是 $n$ 的多项式。
因此对点与组的循环、次数小于 $J$ 的多项式上的算术、以及最终
插值，只在上述三个枚举的每次迭代中增加多项式个运算。

位复杂度与空间界见附录~\ref{subsec:overview-arithmetic}。

\subsection{构造低次替代线性型}
\label{subsec:preview-auxiliary}
考虑子句中 $X_v$ 的一次出现。在 $T=0$ 处其替换必须等于
$X_v$。在每个 $t=1,\ldots,Q$ 处，该子句必须能只用
其指派组与少数共享测量来求值同一替换。
在每个点选取不同的线性型再插值，一般要付出次数 $Q$ 的代价。
为使 Taylor 恢复廉价，我们需要次数 $R$ 远小于 $Q$。

假设替代中的某一项只在 $t=2$ 处使用子句指派组之外的变量。
把该项乘以 $1-T/2$：在零点不变，在 $t=2$ 处消失。
在其他每个求值点，子句都能从自己的组读到该变量。
同样的技巧适用于多个这样的点：每点一个因子，每点付出
一次次数。因此，挑战是把每个输入写成这样的和：共享测量
确定其中一部分，其余各项只在少数几个点用到指派组之外的
变量。我们先在记录组选取的辅助坐标上构造这样的分解，
再代入公式的变量。

\medskip\noindent
\textbf{第 1 步：建立辅助坐标。}\par\nobreak
\smallskip\noindent
在每个求值点，出现 $(e,v)$ 有两个组选取：变量 $v$
所在的组，以及子句 $e$ 被指派的组。跨越 $Q$ 个点，
这些选取构成 $Q$ 个对的序列。因此我们使用坐标集
\[
 \mathcal B=([r]\times[r])^Q,\qquad D=|\mathcal B|=r^{2Q}.
\]
把下标写作
\[
 u=\bigl((a_u(1),d_u(1)),\ldots,(a_u(Q),d_u(Q))\bigr).
\]
我们将为每个这样的下标构造一行替代，因此无论以后对公式
作出何种组选取，它都可用。

设 $\vec\zeta=(\zeta_w)_{w\in\mathcal B}$ 是原变量的形式
占位符，将在第 4 步代入。
在点 $t$，坐标 $\zeta_w$ 属于组 $a_w(t)$，而 $\zeta_u$
的替代被指派到组 $d_u(t)$。于是该替代可自由使用的坐标是
\[
 \mathcal S_{t,d_u(t)},\qquad
 \mathcal S_{t,g}:=\{w\in\mathcal B:a_w(t)=g\}.
\]
特别地，$\zeta_w$ 属于指派组当且仅当 $a_w(t)=d_u(t)$。
集合 $\mathcal S_{t,g}$ 划分 $\mathcal B$，每组含 $D/r$
个坐标。

我们的辅助目标是构造系数行 $\vec P_u(T)\in\mathbb Q[T]^D$
与一个矩阵
\[
 C\in\{0,1,-1\}^{c\times D},\qquad c\le\frac{D\tau}{3r},
\]
使得
\[
 \langle\vec P_u(0),\vec\zeta\rangle=\zeta_u,
 \qquad \deg\vec P_u\le R-1,
\]
且在每个 $t\in[Q]$ 处，替代线性型
$\langle\vec P_u(t),\vec\zeta\rangle$ 是可用数据
\begin{equation}\label{eq:overview-auxiliary-requirement}
 \vec\zeta|_{\mathcal S_{t,d_u(t)}}\quad\text{与}\quad C\vec\zeta
\end{equation}
的线性组合。同一批测量 $C\vec\zeta$ 服务于每一行与每个
求值点。这正是第~\ref{subsec:preview-separation} 节 design
中依赖条件的辅助版本。

\medskip\noindent
\textbf{第 2 步：构造常数矩阵。}\par\nobreak
\smallskip\noindent
从恒等式
\[
 \zeta_u=((I_D-M_0)\vec\zeta)_u+
                   \sum_{w\in\mathcal B}M_0[u,w]\zeta_w.
\]
出发。我们希望第一项由 $C\vec\zeta$ 确定。
在和式中，每个非零系数应当只在少数几个点连接不相容的组，
使得低次多项式可以去除这些依赖。定义这些点为
\[
 E_{u,w}=\{t\in[Q]:a_w(t)\ne d_u(t)\}.
\]
对 $M_0\in\mathbb Q^{D\times D}$ 的两个要求是
\begin{equation}\label{eq:outline-matrix-conditions}
 \begin{gathered}
 M_0\vec\xi=\vec\xi\quad\text{对每个 }\vec\xi\in\ker C,\\
 M_0[u,w]\ne0\quad\Longrightarrow\quad |E_{u,w}|\le R-1.
 \end{gathered}
\end{equation}
第一条说 $I_D-M_0$ 在 $\ker C$ 上为零。等价地，其每行都在
$C$ 的行空间中，故 $(I_D-M_0)\vec\zeta$ 是测量 $C\vec\zeta$
的线性函数。第二条将为第 3 步所用多项式的次数提供界。

\smallskip\noindent
\emph{通过保持 $f$ 低次来限制支撑。}
设 $\Pi_s$ 为交换每个下标中第 $s$ 对的两个分量、其余对
不变的置换矩阵。这些矩阵两两交换，且满足 $\Pi_s^2=I_D$。
令
\[
 V_0=\prod_{s=1}^Q\Pi_s,\qquad
 W_0=\frac12\sum_{s=1}^Q(I_D-\Pi_s),\qquad
 M_0=V_0f(W_0),
\]
其中我们将选取次数至多 $R-1$ 的多项式 $f$。
完全交换 $V_0$ 已经尊重每个组：其在行 $u$ 的非零元位于满足
\[
 (a_w(s),d_w(s))=(d_u(s),a_u(s))\quad\text{对每个 }s
\]
的列 $w$。特别地，$E_{u,w}$ 为空。

乘以 $f(W_0)$ 可以撤销一些交换，但 $f$ 的次数限制了个数。
展开 $f(W_0)$ 给出至多 $R-1$ 个交换的乘积。
由 $\Pi_s^2=I_D$，重复的交换相消。
因此乘以 $V_0$ 给出
\begin{equation}\label{eq:overview-partial-swaps}
 M_0=\sum_{\substack{S\subseteq[Q]\\|S|\le R-1}}
             \gamma_S\prod_{s\notin S}\Pi_s
 \qquad\text{（对某些 }\gamma_S\in\mathbb Q\text{）}.
\end{equation}
以 $S$ 为下标的被加项交换 $S$ 之外的每一对。其非零元满足
$a_w(s)=d_u(s)$ 对每个 $s\notin S$，故 $E_{u,w}\subseteq S$。
只要 $|E_{u,w}|>R-1$，每个被加项都为零，其和亦然。
这证明了对任何次数至多 $R-1$ 的多项式 $f$，
\eqref{eq:outline-matrix-conditions} 的第二个条件成立。

\smallskip\noindent
\emph{选取 $f$ 使其固定大多数基方向。}
现在选取 $f$，使 $M_0$ 固定大多数基方向。
先考虑一个位置，其坐标以 $[r]\times[r]$ 为下标。
设 $\vec e_{ab}$ 为在 $(a,b)$ 取一、其余取零的向量。
交换该对有正交特征基
\[
 \begin{array}{ll}
 \vec e_{aa}\quad(a\in[r]),\quad
 \vec e_{ab}+\vec e_{ba}\quad(a<b),
     &\text{特征值 }+1,\\[2pt]
 \vec e_{ab}-\vec e_{ba}\quad(a<b),
     &\text{特征值 }-1.
 \end{array}
\]
特征值 $+1$ 的向量有 $r(r+1)/2$ 个，特征值 $-1$ 的有
$r(r-1)/2$ 个。所有分量都在 $\{0,1,-1\}$ 中。

对每个位置 $s\in[Q]$，选取这样一个向量 $\vec\psi_s$，
并按分量定义 $\vec\psi\in\mathbb Q^{\mathcal B}$：
\[
 \psi_u=\prod_{s=1}^Q\psi_s(a_u(s),d_u(s)).
\]
这 $D=(r^2)^Q$ 种选取构成正交基：内积分解为 $Q$ 个位置上
内积的乘积。交换 $\Pi_s$ 以 $\vec\psi_s$ 的特征值作用在
$\vec\psi$ 上。若恰有 $j$ 个选取的特征值为 $-1$，则
\[
 W_0\vec\psi=j\vec\psi,\qquad
 V_0\vec\psi=(-1)^j\vec\psi,\qquad
 M_0\vec\psi=(-1)^jf(j)\vec\psi.
\]
因此，置 $f(j)=(-1)^j$ 即可使 $M_0$ 固定这些向量。

对均匀选取的基向量，数 $j$ 的分布为
\[
 W\sim\operatorname{Bin}(Q,p),\qquad p=\frac{r-1}{2r},
\]
因为在每个位置，$r^2$ 个选取中有比例 $p$ 的特征值为 $-1$。
因此大多数基向量的 $j$ 接近 $Qp$。
取靠近 $Qp$ 的 $R$ 个相邻整数的区间 $\mathcal I$，
对 $j\in\mathcal I$ 插值 $f(j)=(-1)^j$。
这给出 $\deg f\le R-1$，并使 $M_0$ 固定每个 $j\in\mathcal I$
的基向量。第 5 步给出端点并检验该区间落在 $\{0,\ldots,Q\}$
内。

\smallskip\noindent
\emph{测量例外方向。}
用剩余基向量的转置作为 $C$ 的行。其分量在 $\{0,1,-1\}$
中，且由正交性，$\ker C$ 由上述被固定的基向量张成。
这证明了 \eqref{eq:outline-matrix-conditions} 的第一个条件。

现在可以看出如何选取 $R$。$W$ 的标准差是
$\sqrt{Qp(1-p)}=\Theta(\sqrt Q)$。取 $R$ 远大于 $\sqrt Q$
使大多数基向量落入区间，只需测量一小部分。
我们仍可取 $R$ 远小于 $Q$，使 Taylor 恢复每点只需
$1+\lfloor 3nR/Q\rfloor=o(n)$ 个系数。
对第 5 步的选取，二项尾界给出
\[
 c=D\Pr[W\notin\mathcal I]
      \le 2D\exp\!\left(-\frac{R^2}{8Q}\right)
      \le\frac{D\tau}{3r},
\]
正如所需。于是低次数给出受限支撑，而插值使修正
$I_D-M_0$ 只依赖这少数测量。

\medskip\noindent
\textbf{第 3 步：定义辅助替代行。}\par\nobreak
\smallskip\noindent
现在把第 2 步的 $\zeta_u$ 分解转化为其替代。
保持由测量确定的项不变。对每个 $M_0[u,w]\ne0$，
把另一项中的系数乘以
\[
 \varphi_{u,w}(T)=\prod_{s\in E_{u,w}}(1-T/s).
\]
该多项式在零点等于一，并在 $\zeta_w$ 落在指派组之外的
每个点为零。用恒等式定义系数行 $\vec P_u(T)$：
\begin{equation}\label{eq:outline-auxiliary-row}
 \langle\vec P_u(T),\vec\zeta\rangle
 =((I_D-M_0)\vec\zeta)_u+
   \sum_{\substack{w\in\mathcal B\\M_0[u,w]\ne0}}
                M_0[u,w]\varphi_{u,w}(T)\zeta_w.
\end{equation}
这直接给出全部三个所需性质：
\begin{itemize}
\item 在 $T=0$ 处每个乘子为一，两项之和为 $\zeta_u$。
\item 每个乘子的次数为 $|E_{u,w}|\le R-1$，故
$\deg\vec P_u\le R-1$。
\item 在 $T=t$ 处，第一项是 $C\vec\zeta$ 的线性函数。
和式中，系数只有在 $a_w(t)=d_u(t)$ 时才能保留，
因此只剩 $\mathcal S_{t,d_u(t)}$ 中的坐标。
\end{itemize}
于是该替代恰好只使用
\eqref{eq:overview-auxiliary-requirement} 所允许的数据。

\medskip\noindent
\textbf{第 4 步：代入原变量。}\par\nobreak
\smallskip\noindent
辅助行已构造完毕。为将其用于公式，选取日程
\[
 a_v\in[r]^Q\quad(v\in[n]),\qquad
 d_e\in[r]^Q\quad(e\in\mathcal F).
\]
在点 $t$，把变量 $v$ 放入组 $a_v(t)$，把子句 $e$
指派到组 $d_e(t)$：
\[
 S_{t,g}=\{v:a_v(t)=g\},\qquad \sigma_e(t)=d_e(t).
\]
出现 $(e,v)$ 使用辅助标号
\[
 u(e,v)=\bigl((a_v(s),d_e(s))\bigr)_{s\in[Q]}\in\mathcal B.
\]
设 $h$ 是任一标号获得的出现数的最大值。我们需要变量组
小且 $h$ 小。对我们的三次出现、宽度三的公式，
引理~\ref{lem:incidence-labels} 在多项式时间内确定性选取
日程，使
\[
 |S_{t,g}|\le2n/r,\qquad h\le2M/D,
\]
只要
\begin{equation}\label{eq:overview-balance-requirements}
 12D^2\le M,\qquad 2Qr^2\le n.
\end{equation}
想法是对日程的随机选取去随机化。若所有 $a_v$ 与 $d_e$
独立均匀取自 $[r]^Q$，则每个标号的期望重数为 $M/D$，
每组 $S_{t,g}$ 的期望大小为 $n/r$。
在 \eqref{eq:overview-balance-requirements} 的假设下，
标准的条件期望法在多项式时间内去随机化该选取，
保证每个标号的重数至多 $2M/D$、每组的大小至多 $2n/r$。
第 5 步对我们的参数选取验证这些假设。

不同的出现可以得到同一标号，却在该坐标要求不同的原变量。
使用 $\mathcal B$ 的 $h$ 份副本，把标号相同的出现放进不同
副本。令 $j(e,v)\in[h]$ 为 $(e,v)$ 所用的副本；
把 $X_v$ 放在其坐标 $u(e,v)$ 处，未用位置放零。
记 $\vec\zeta^{(j)}(\vec X)$ 为副本 $j$ 中所得的填充向量。
每个分量是零或一个原变量。$X_v$ 的重复出现使用同一取值，
因此这些副本不引入新的布尔变量。

把辅助行用于包含该出现的副本：
\begin{equation}\label{eq:overview-transferred-row}
 \langle\vec P_{e,v}(T),\vec X\rangle
 =\left\langle\vec P_{u(e,v)}(T),
                  \vec\zeta^{(j(e,v))}(\vec X)\right\rangle.
\end{equation}
在零点，右端读取含 $X_v$ 的坐标，代入变量保持次数界
$R-1$。在点 $t$，其可自由使用的坐标只含满足
$a_{v'}(t)=d_e(t)$ 的变量 $v'$，即恰为子句指派组
$S_{t,d_e(t)}$ 中的变量。子句 $e$ 的所有出现都使用同一
指派组。

剩余的依赖通过 $h$ 份副本上的 $hc$ 个线性测量
$C\vec\zeta^{(j)}(\vec X)$ 传递。它们可能涉及多个原组。
固定形如 $x_1+x_2$ 的总测量仍会耦合分别含 $x_1$ 与 $x_2$
的两个组。因此我们把每个测量按各组贡献拆分。
对每个测量 $\ell(\vec X)$，写
\[
 \ell(\vec X)=\sum_{g=1}^r
        \ell|_{S_{t,g}}(\vec X|_{S_{t,g}}),
\]
并把这些限制的系数行归入 $C_{t,g}$。
每个辅助测量现在都是组测量的和。
对指派到组 $g$ 的子句，来自 $S_{t,g}$ 的部分是局部线性型，
而其余每一部分都是相应组测量的线性组合。
这正是 \eqref{eq:outline-separation} 的条件。

总行数为 $rhc$。由 $hD\le2M\le6n$，
\[
 \sum_{g=1}^r c_{t,g}
 =rhc\le rh\frac{D\tau}{3r}\le2\tau n.
\]
$C$ 的每行有 $D$ 个 $\{0,1,-1\}$ 中的分量。
代入变量、合并重复副本、限制到一个组，都不会增加系数
绝对值之和。因此 $C_{t,g}$ 的每行都是整数的，且绝对系数
和至多 $D$。结合组大小界，这证明替代型在每个点都容许
$(r,\tau,D)$-design。

\medskip\noindent
\textbf{第 5 步：验证所承诺的参数界。}\par\nobreak
\smallskip\noindent
构造需要两个平衡不等式
\eqref{eq:overview-balance-requirements} 与尾界
$2e^{-R^2/(8Q)}\le\tau/(3r)$。现在按
\eqref{eq:overview-parameter-bounds} 所述的阶选取参数，
并验证这些要求。

令 $\lambda=\lceil\log_2 n\rceil$、
$\nu=\lceil\log_2\lambda\rceil$，并取
\begin{equation}\label{eq:overview-explicit-parameters}
 \begin{aligned}
 r&=\left\lfloor\frac{\sqrt\lambda}{\nu^2}\right\rfloor,
 &Q&=\left\lfloor\frac{\lambda}{128\nu}\right\rfloor,
 &R&=4\lceil\sqrt\lambda\rceil,\\
 \tau&=\frac1{r\lambda},
 &D&=r^{2Q}.
 \end{aligned}
\end{equation}
它们具有 \eqref{eq:overview-parameter-bounds} 所述的阶。
对足够大的 $n$，有 $r\ge2$、$R\ge32$ 与 $6R\le Q$
（因 $R/Q=O(\nu/\sqrt\lambda)\to0$）。逐一检查其余要求：
\begin{itemize}
\item \emph{插值区间。}
取第 2 步中的 $p=(r-1)/(2r)$，令
\[
 A_0=\left\lfloor Qp-\frac{R-1}{2}\right\rfloor,\qquad
 \mathcal I=\{A_0,\ldots,A_0+R-1\}.
\]
由 $Q/4\le Qp\le Q/2$ 与 $6R\le Q$，该区间落在
$\{0,\ldots,Q\}$ 内。此外，$R\ge32$ 保证它包含与 $Qp$
距离在 $R/4$ 之内的每个整数。二项尾界因此给出第 2 步
所用估计 $\Pr[W\notin\mathcal I]\le2e^{-R^2/(8Q)}$。
\item \emph{多项式维数与平衡日程。}
由 $\log_2 r\le\nu/2$，
\[
 \log_2 D=2Q\log_2 r\le\lambda/128,\qquad
 D\le(2n)^{1/128}.
\]
对足够大的 $n$，这给出
\[
 12D^2\le12(2n)^{1/64}\le n\le M,\qquad
 2Qr^2=O(\lambda^2)\le n,
\]
正是第 4 步所需。
\item \emph{少量测量。}
我们有 $R^2/(8Q)\ge256\nu$。由 $\lambda\le2^\nu$ 与
$r^2\le\lambda$，
\[
 2e^{-R^2/(8Q)}\le2e^{-256\nu}
 \le\frac1{3\lambda^2}
 \le\frac1{3r^2\lambda}=\frac{\tau}{3r}.
\]
于是第 2 步给出第 4 步所用的辅助测量界。
\end{itemize}

最后，所有构造维数都被 $n$ 的多项式界定。
引理~\ref{lem:row-coefficients} 给出第 2--3 步矩阵与行的
多项式构造时间，引理~\ref{lem:incidence-labels}
对第 4 步的日程亦然。传递只用 $hD\le6n$ 个坐标位置。
因此我们在多项式时间内构造出替代型：在零点恢复每个原
输入、次数至多 $R-1\le R$，并在每个 $t\in[Q]$
处容许所承诺的 $(r,\tau,D)$-design。
这提供了第~\ref{subsec:preview-parameters} 节所用的界；
引理~\ref{lem:incidence-parameters} 为完整证明的更一般
情形记录了参数检查。

\subsection{一般 \texorpdfstring{$\#k$-SAT}{\#k-SAT} 的较弱算法}
\label{subsec:preview-general-sat}
上述算法已经给出一般 $\#k$-SAT 的亚指数算法：
不交稀疏化之后，我们把每个实例直接归约为每个变量至多
出现三次的 $3$-CNF。两步都保持计数，如下。

固定 $k\ge2$。我们使用稀疏化~\cite{IPZ} 的保持计数版本
~\cite[Appendix~A]{DHMTW}，定量界取自
~\cite[Section~6]{CIP} 与
\cite[Lemma~5.14 and equation~(5.4)]{CalabroThesis}。
如引理~\ref{lem:sparsify} 所记，这把输入计数表示为至多
$\operatorname{poly}_k(n,\sigma^{-1})2^{\sigma n}$ 个公式
$G$ 的计数之和。每个 $G$ 有相同的 $n$ 个声明变量，
每个变量的出现至多为
\[
 B=O_k\!\left((\sigma^{-1}\log(2/\sigma))^{k-1}\right)
\]
次，从而非空子句至多 $Bn$ 个。空子句直接使计数为零。

第一步把 $G$ 变成计数不变的 $3$-CNF。保留规模至多三的
子句。对更长的子句，用二元 OR 门计算其析取，并要求最终
输出为一。对门 $z=u\vee v$，施加
\[
 (\neg u\vee z)\ \wedge\ (\neg v\vee z)\ \wedge\
 (u\vee v\vee\neg z).
\]
这里 $u,v$ 是输入文字或更早的门输出，$z$ 是新的。
这些子句强制与 OR 相等，故门值由原赋值唯一确定。
因此，$G$ 的可满足赋值恰有一种到新公式的扩张，
不满足 $G$ 的赋值则没有。这每子句用 $O(k)$ 个门。

第二步减少出现次数。若变量 $x$ 在该 $3$-CNF 中出现
$s\ge2$ 次，则把其出现替换为不同副本 $x_1,\ldots,x_s$
（保持符号），并加上 Tovey 的蕴含圈~\cite[p.~86]{Tov84}：
\[
 (\neg x_1\vee x_2)\ \wedge\cdots\wedge\
 (\neg x_{s-1}\vee x_s)\ \wedge\ (\neg x_s\vee x_1).
\]
该圈迫使所有副本一致。每个副本出现在一个继承自 $3$-CNF
的子句与两个圈子句中，因此恰在三个子句中。
至多一次出现的变量保持不变，包括未使用的声明变量。
这个标准的变量复制归约~\cite[Proposition~4.7]{Gol08}
同样给出可满足赋值之间的双射。最终公式 $\widehat G$
满足 $\#\widehat G=\#G$ 且
\[
 n\le N:=\#\{\text{$\widehat G$ 的变量}\}=O_k(Bn).
\]
其所有子句（含双文字圈子句）规模至多三，
且每个变量至多出现三次。因此它恰是本概览所述设置的实例。

对每个 $\widehat G$ 应用 \eqref{eq:overview-bounded-time}
并把计数相加。当 $\sigma^{-1}$ 被 $\log n$ 的固定幂界定时
$B$ 亦然，所得时间为
\[
 (n+m)^{O_k(1)}
 2^{O_k\!\left(\sigma n+
          Bn(\log\log n)^2/\sqrt{\log n}\right)}.
\]
例如，把 $\sigma$ 取到 $(\log n)^{-1/(2k)}$ 的常数倍内，得
\[
 (n+m)^{O_k(1)}
 2^{O_k\!\left(n(\log\log n)^{k+1}/(\log n)^{1/(2k)}\right)}
   = (n+m)^{O_k(1)}2^{o(n)}
\]
与多项式空间，逐个处理稀疏化树与归约后的实例。
参数 $\sigma$ 可取有理数。

\subsection{经由对 \texorpdfstring{$\mathbb F_2$}{F\_2} 上稀疏低次多项式方程解的计数得到更快的 \texorpdfstring{$\#k$-SAT}{\#k-SAT} 算法}
\label{subsec:preview-full-case}

第~\ref{subsec:preview-general-sat} 节的运行时间指数是
$O_k(\sigma n+B\rho(n))$，其中
$\rho(n)=n(\log\log n)^2/\sqrt{\log n}$。
若稀疏化能在 $\sigma=\Theta(\rho(n)/n)$ 处给出
$B=O_k(1)$，我们就已得到一般 $\#k$-SAT 的
$2^{O_k(\rho(n))}$ 时间算法（差多项式因子）：复制在每个
叶上只用 $O_k(n)$ 个变量，叶数也落在同一界内。
然而在这个 $\sigma$ 的取法下，保证只有
$B=(\log n)^{O_k(1)}$。

我们分两步得到 $2^{O_k(\rho(n))}$ 界。第一，本概览的计数
算法推广到 $\mathbb F_2$ 上 $O(n)$ 个多项式方程构成的
方程组，每个方程次数至多 $k$ 且单项数为多对数级
（定理~\ref{thm:overview-sparse-polynomials}）。
第二，即使叶 $G$ 有 $Bn$ 个子句，我们也能在原 $n$ 个变量上
构造这样的方程组，其公共零点恰为 $G$ 的可满足赋值。
我们先以一个自含的随机化构造展示该归约，
再解释计数算法。

\medskip\noindent
\textbf{到稀疏多项式方程的随机化归约。}
若 $x$ 违反子句 $e$ 则写 $v_e(x)=1$，否则为 $0$。
子句集 $S$ 上的\emph{奇偶检验}是 $\mathbb F_2$ 上的多项式
方程 $\sum_{e\in S}v_e(x)=0$：当 $x$ 违反 $S$ 中偶数个
子句时通过。每个可满足赋值都通过，因此我们只需抓住每个
不可满足的赋值。若 $G$ 有 $m_G\le n$ 个子句，
每子句一个检验即够。设 $m_G>n$。

短检验容易漏掉单个被违反的子句，但所有只违反该子句的
赋值给出相同回答。因此只有 $m_G$ 个单点模式需要抓住。
使用 $8n$ 个独立检验。对每个检验，有放回地均匀抽取
$L=2B\lceil\log_2(2m_G)\rceil$ 个子句。对每次抽取，
以概率 $1/2$ 独立地把其违反位加入奇偶和；
被包含两次的副本相消。由 $m_G\le Bn$，每个子句在期望上
总共至少被抽取对数次。对固定的 $w\ge1$ 个被违反子句的
集合，令 $u=Lw/m_G$ 为每次检验的期望命中数。
若样本命中该集合，则首个命中处的硬币在其余硬币固定时
以概率 $1/2$ 使奇性为奇。因此拒绝概率为
\[
 p_w=\tfrac12\bigl(1-(1-w/m_G)^L\bigr)\ge\tfrac12(1-e^{-u}).
\]

对少量违反我们计\emph{集合}，对大量违反计\emph{赋值}：
\begin{itemize}
\item 若 $u\le1$，则 $p_w\ge u/4$。大小为 $w$ 的集合至多
$m_G^w$ 个，故至少一个集合逃过所有检验的概率至多为
$m_G^w e^{-8np_w}\le m_G^w e^{-2nLw/m_G}\le16^{-w}$。
对 $w\ge1$ 求和，失败概率至多 $1/15$。
\item 若 $u>1$，则 $p_w>1/4$。对 $2^n$ 个赋值用联合界，
此情形下任一赋值逃过的概率被
$2^n(3/4)^{8n}<1/4$ 界住。
\end{itemize}
于是以大于 $2/3$ 的概率，公共零点恰为 $G$ 的可满足赋值。
每个 $v_e$ 是至多 $k$ 个形如 $x_j$ 或 $1+x_j$ 的因子的
乘积，故每个检验的次数至多 $k$、单项数至多 $2^kL$
（$n$ 的多对数）。
第~\ref{sec:hashes} 节确定性地构造具有同一解集的短检验。
我们先用来自展开码的奇偶校验保证仍然通过的非可满足赋值
必违反许多子句；随机短检验随后以至少 $1/4$ 的概率抓住
每个这样的赋值。对每个候选检验，我们统计同时通过该检验
与已选全部检验的赋值个数，并追加使该计数最小的候选：
每个可满足赋值通过每个候选，因此这至少消除剩余
不可满足赋值的四分之一。

\medskip\noindent
\textbf{计数公共零点。}
回忆第~\ref{subsec:preview-encoding} 节：$x$ 满足子句
$e$ 时 $\Phi_e(x)$ 为 $1$，否则为 $0$。
当我们一次检验整个子句集 $S_i$ 时，该因子应换成什么？
先用普通整数加法统计\emph{本次检验}中被违反的子句数：
\[
 L_i(x)=\sum_{e\in S_i}v_e(x)=\sum_{e\in S_i}\bigl(1-\Phi_e(x)\bigr).
\]
检验恰在 $L_i(x)$ 为偶数时通过。取 $\ell=n+1$，
使至多 $2^n$ 的计数由其模 $2^\ell$ 的剩余确定。
补余 Beigel--Tarui 多项式 $p_\ell$
（引理~\ref{lem:rational-idempotent}）在偶整数上模
$2^\ell$ 为 $1$，在奇整数上为 $0$。
例如，在 $S_i$ 中违反两个子句仍通过检验：
$p_\ell(2)\equiv1\pmod{2^\ell}$。
以 $m_{\mathrm{test}}$ 记检验数，我们的计数表达式变为
\[
 \underbrace{\sum_x\prod_{e\in G}\Phi_e(x)}_{\text{每子句一个因子}}
 \qquad\longrightarrow\qquad
 \underbrace{\sum_x\prod_{i=1}^{m_{\mathrm{test}}} p_\ell(L_i(x))}_{\text{每检验一个因子}}.
\]
左边是 $\#G$。右边模 $2^\ell$ 的最小非负剩余是通过全部
检验的赋值个数。当检验抓住每个不可满足赋值时，
它恰为 $\#G$。同一编码适用于 $\mathbb F_2$ 上任何稀疏
多项式方程 $P_i(x)=0$：取 $L_i(x)$ 为其列出的单项的
整数和。方程成立同样恰在这个和为偶数时。

为复用 Taylor 步骤，我们需要输入次数 $O(n)$ 的乘积。
直接相乘 $m_{\mathrm{test}}$ 个检验因子只有界
$O(km_{\mathrm{test}}\ell)$，可能关于 $n$ 是二次的。
关键是\emph{跨全部检验整体}截断乘积。
第~\ref{subsec:rational-lift} 节把 $p_\ell(y)$ 写成
$e_j(y)h(y)^j$ 项之和，其中 $e_j$ 次数至多一，
$h(y)=y(y-1)$ 在每个整数处为偶。
因此在布尔赋值上模 $2^\ell$，含 $\ell$ 个及以上 $h$
副本的乘积项为零，即使这些副本来自不同检验。
我们只保留总共少于 $\ell$ 个副本的项，并在相乘时追踪
该数目。每个保留项在 $L_i$ 中至多有 $m_{\mathrm{test}}$
个线性因子与 $\ell-1$ 个二次因子。由 $\deg L_i\le k$，
其输入次数至多 $k(m_{\mathrm{test}}+2(\ell-1))=O(n)$。

对替代型，重要的是每个因子读取哪些变量。
第~\ref{subsec:preview-deformation} 节对子句中每个变量
使用一个型，现在我们对检验中每个变量使用一个型，
并在整个检验中复用。每个稀疏方程只读取多对数个变量。
对 $O(n/\sqrt{\log n})$ 个最频繁变量分支后，
每个剩余变量只出现在多对数个检验中，
正是第~\ref{sec:construction} 节所需。
在每个分支上，固定共享测量再次允许我们一次对一个组求和。
由于截断乘积的输入次数为 $O(n)$，次数 $R$ 的替代给出
$\deg_T H=O(nR)$：每点同样 $O(1+nR/Q)$ 个 Taylor 系数
即可恢复 $T=0$ 处 $h$ 的 $0,1,\ldots,\ell-1$ 个副本的
贡献。它们的和取模 $2^\ell$ 的最小非负剩余，
就是精确的分支计数。把分支计数相加即得公共零点个数。

\begin{theorem}[稀疏多项式公共零点的计数]
\label{thm:overview-sparse-polynomials}
固定整数 $k\ge2$ 与 $c\ge1$。设
$P_1,\ldots,P_{m_{\mathrm{test}}}$ 是 $n\ge20$ 个布尔变量上
的 $\mathbb F_2$ 多项式，$m_{\mathrm{test}}\le cn$。
每个 $P_i$ 由至多 $(\log n)^c$ 个次数至多 $k$ 的单项
列表给出。$\{0,1\}^n$ 中公共零点的精确个数可确定性地在
$2^{O_{k,c}(n(\log\log n)^2/\sqrt{\log n})}$ 次位运算与
$n^{O_{k,c}(1)}$ 空间内计算。
\end{theorem}
第~\ref{sec:counting} 与~\ref{sec:construction} 节提供
计数例程与替代型；第~\ref{sec:setup} 节用它们对多项式
方程直接证明该定理。

对为 $G$ 构造的最终方程组，公共零点个数是 $\#G$。
第~\ref{sec:hashes} 节对每个叶只用
$2^{(\log n)^{O_k(1)}}$ 次计数器调用，
保持定理~\ref{thm:overview-sparse-polynomials} 的运行
时间指数。最后，在
第~\ref{subsec:general-reduction} 节取阶为
$(\log\log n)^2/\sqrt{\log n}$ 的 $\sigma$，
使 $B$ 为多对数级且叶数在
$2^{O_k(n(\log\log n)^2/\sqrt{\log n})}$ 之内。
各叶的可满足赋值互不相交，故把计数相加给出
命题~\ref{prop:interpolation-bound} 的界。
\section{经由共享线性摘要的快速计数算法}\label{sec:counting}
第~\ref{sec:counting}--\ref{sec:setup} 节证明
定理~\ref{thm:overview-sparse-polynomials}。
我们从第~\ref{subsec:preview-encoding} 节的和出发：
对每个布尔赋值，把各因子相乘，再把乘积对所有赋值相加。
这里因子可以是任意有理多项式。我们说明在给定替代型与测量时
如何计算该和。第~\ref{sec:construction} 节只用每个因子读取的
变量来构造它们。第~\ref{sec:setup} 节为每个稀疏多项式方程
提供一个因子。

\subsection{多项式因子与要恢复的值}
\label{subsec:product-interface}
设 $\mathcal F$ 为因子的下标集，$I_f\subseteq[n]$ 为 $\Phi_f$
读取的变量。取有理系数时，目标是
\begin{equation}\label{eq:general-product-sum}
 Z=\sum_{x\in\{0,1\}^n}\prod_{f\in\mathcal F}\Phi_f((x_v)_{v\in I_f}).
\end{equation}
对子句满足指示子，$Z$ 就是可满足赋值的个数。
证明对任意交换 $\mathbb Q$-代数 $\mathcal A$ 中的系数同样成立：
此时和与积在 $\mathcal A$ 中进行，插值只用有理标量。
本节读者可始终取 $\mathcal A=\mathbb Q$。

不同因子可对同一变量使用不同替换。
给每个关联 $(f,v)$ 一个自己的未定元 $Y_{f,v}$，并设
\begin{equation}\label{eq:general-product-degree}
 \deg_{\{Y_{f,v}\}}\!
    \left(\prod_{f\in\mathcal F}\Phi_f((Y_{f,v})_{v\in I_f})\right)\le d .
\end{equation}
这是整个乘积在 $\mathcal A$ 上计算的次数。
取次数至多 $R$、满足 $\vec P_{f,v}(0)=\delta_v$ 的替代行
$\vec P_{f,v}(T)$，并作
\begin{equation}\label{eq:general-count-polynomial}
 H(T)=\sum_{x\in\{0,1\}^n}\prod_{f\in\mathcal F}
   \Phi_f\bigl((\langle\vec P_{f,v}(T),x\rangle)_{v\in I_f}\bigr)
   \in\mathcal A[T].
\end{equation}
替换在零点恢复输入，并把次数至多乘以 $R$。
于是 $H(0)=Z$ 且 $\deg_T H\le d_H:=dR$。
我们的应用将如受限出现次数的例子那样取 $d=O(n)$。

\subsection{所需的行、划分与测量}
\label{subsec:layout-interface}
我们使用与第~\ref{subsec:preview-separation} 节相同的 design，
只是对每个因子（而非每个子句）作组指派。
下面的最后一个条件记录精确算术所需的系数界。
记 $\widetilde\xi$ 表示从一个组经补零延拓到全部 $n$ 个坐标。

\begin{definition}\label{def:layout}
对因子 scope $(I_f)_{f\in\mathcal F}$，设 $r\ge2$、$D\ge1$、
$1\le R<Q$ 为整数，$0<\tau\le1/2$。我们要求如下数据与性质。
\begin{enumerate}
\item 对每个 $t\in[Q]$，一个允许空部分的划分
$[n]=S_{t,1}\sqcup\cdots\sqcup S_{t,r}$，且
$|S_{t,g}|\le2n/r$。
\item 对每个因子 $f\in\mathcal F$，一个序列
$\sigma_f\in[r]^Q$。值 $\sigma_f(t)$ 在点 $t$
把因子 $f$ 指派到一个组。
\item 对每个 $t\in[Q]$ 与 $g\in[r]$，一个整数矩阵
$C_{t,g}\in\mathbb Z^{[c_{t,g}]\times S_{t,g}}$，其中
$c_{t,g}$ 为非负整数。每行的绝对值之和至多 $D$，且
\[
 \sum_{g=1}^r c_{t,g}\le2\tau n.
\]
对赋值 $x$，$C_{t,g}(x|_{S_{t,g}})$ 的分量就是组 $g$
上的测量。
\item 对每个 $f\in\mathcal F$ 与 $v\in I_f$，一行
$\vec{P}_{f,v}(T)\in\mathbb Q[T]^{1\times n}$，满足
$\vec{P}_{f,v}(0)=\delta_v$ 与
$\deg \vec{P}_{f,v}\le R$。
对每个 $t\in[Q]$、$g\in[r]\setminus\{\sigma_f(t)\}$ 与
$\xi\in\mathbb Q^{S_{t,g}}$，要求
\[
 C_{t,g}\xi=\vec{0}\quad\Longrightarrow\quad
 \langle \vec{P}_{f,v}(t), \widetilde\xi\rangle=0,
\]
其中 $\widetilde\xi\in\mathbb Q^n$ 是 $\xi$ 的补零延拓。
\item 记 $d_{\mathrm{den}}=(R-1)!Q!$，
$d_{\mathrm{den}} \vec{P}_{f,v}$ 的每个系数都是整数，
其位数为 $O(Q\log Q+\log D)$（隐含常数为绝对常数）。
\end{enumerate}
\end{definition}

为解读核条件，固定 $t$、$f$ 与 $v\in I_f$，
令 $g=\sigma_f(t)$。写 $x_h=x|_{S_{t,h}}$。
若 $x,x'\in\mathbb Q^n$ 对某个 $h\ne g$ 在 $S_{t,h}$
之外一致，且满足 $C_{t,h}x_h=C_{t,h}x'_h$，
则 $x'_h-x_h\in\ker C_{t,h}$。该条件给出
\[
 \langle \vec{P}_{f,v}(t), x'\rangle=\langle \vec{P}_{f,v}(t), x\rangle.
\]
一次对一个组应用之，可知元组
$(\langle \vec{P}_{f,v}(t), x\rangle)_{v\in I_f}$
由 $x_g$ 与 $(C_{t,h}x_h)_{h\ne g}$ 确定。
因此固定全部测量向量后，整个因子成为 $x_g$ 的函数。
这适用于每个多项式因子 $\Phi_f$，包括系数在 $\mathcal A$
中的因子。

布尔向量的每个测量落在 $[-D,D]\cap\mathbb Z$ 中。
我们界定矩阵分量绝对值之和，因为构造可能把同一布尔坐标
放到多个辅助位置。识别出这些位置后，$C_{t,g}$ 的分量
不必属于 $\{0,1,-1\}$。仅凭秩界无法控制可能的测量个数：
单独一行 $(1,2,4,\ldots,2^{n-1})$ 就区分所有 $2^n$
个布尔向量。

\subsection{通过固定测量分离求和}
\label{subsec:measurement-factorization}
固定测量使每个因子只依赖其指派组。下面的引理陈述由此
得到的、在交换环上的计算，因此同样适用于有理值因子与
$\mathcal A$ 中的因子。限制 $x_i=1$ 在 Taylor 系数中
需要，且是各组局部的。

\begin{lemma}\label{lem:measurement-separation}
设 $[n]=S_1\sqcup\cdots\sqcup S_r$ 是一个允许空部分的划分。
对每个 $g\in[r]$，设 $C_g$ 是有 $c_g$ 行、列以 $S_g$
为下标的整数矩阵。设每行的绝对值之和至多为整数 $D\ge0$，
记 $c=\sum_g c_g$。设 $\mathcal A$ 是有单位元的交换环，
$A_g:\mathbb Q^n\to\mathcal A$ 满足
\[
 A_g(x+\xi)=A_g(x)
 \qquad
 (x\in\mathbb Q^n,\ \xi\in\ker C_h,\ h\ne g),
\]
其中 $\xi$ 在 $S_h$ 之外补零。
对任意 $E\subseteq[n]$，和
\[
 Z_E=\sum_{\substack{x\in\{0,1\}^n\\x_i=1\ (i\in E)}}
           \prod_{g=1}^r A_g(x)
\]
可以通过检查至多
\[
 2(2D+1)^c\sum_{g=1}^r2^{|S_g|}
\]
个布尔向量求值，每个向量上至多一次 $A_g$ 求值、
$\poly(n,r,c,\log(D+2))$ 次额外位运算与 $O(1)$
次环运算。除求值器自身的存储与 $O(1)$ 个环元素外，
该过程使用 $\poly(n,r,c,\log(D+2))$ 比特空间。
\end{lemma}
\begin{proof}
对 $y_g\in\mathbb Z^{c_g}$，定义
\[
 \mathcal X_g(y_g)=\{u\in\{0,1\}^{S_g}:C_gu=y_g\},
 \qquad
 \mathcal Y=\prod_{g=1}^r\{-D,\ldots,D\}^{c_g}.
\]
布尔赋值的每个测量元组都属于 $\mathcal Y$，
且 $|\mathcal Y|=(2D+1)^c$。
我们先识别固定测量元组 $\vec y=(y_1,\ldots,y_r)$ 的贡献。
若每个 $\mathcal X_g(y_g)$ 非空，为每个 $g$ 取代表
$\bar x_g\in\mathcal X_g(y_g)$，并令
\[
 \mathcal X_{g,E}(y_g)=
 \{u\in\mathcal X_g(y_g):u_i=1\text{ 对所有 }i\in E\cap S_g\}.
\]
若对每个 $h$ 都有 $x_h\in\mathcal X_h(y_h)$，则
$x_h-\bar x_h\in\ker C_h$。对 $h\ne g$ 逐次应用不变性
假设，得
\[
 A_g\Bigl(\sum_{h=1}^r x_h\Bigr)
 =A_g\Bigl(x_g+\sum_{h\ne g}\bar x_h\Bigr).
\]
测量元组为 $\vec y$ 且以 $E$ 为下标的位固定的赋值构成
$\prod_g\mathcal X_{g,E}(y_g)$。定义
\begin{equation}\label{eq:group-sum}
 \Sigma_g(\vec y)=\sum_{x_g\in\mathcal X_{g,E}(y_g)}
       A_g\Bigl(x_g+\sum_{h\ne g}\bar x_h\Bigr).
\end{equation}
分配律表明这些赋值对 $Z_E$ 的贡献为
$\prod_g\Sigma_g(\vec y)$。若 $\mathcal X_{g,E}(y_g)$
为空，则 $\Sigma_g(\vec y)=0$。
代表 $\bar x_g$ 自身不必满足以 $E$ 为下标的固定位条件。
每个布尔赋值有唯一的测量元组，故
\[
 Z_E=\sum_{\vec y\in\mathcal Y}\prod_{g=1}^r\Sigma_g(\vec y),
\]
其中任何一个 $\mathcal X_g(y_g)$ 为空时乘积取为零。

算法每次处理一个测量元组地求值该恒等式，从
$Z=0_{\mathcal A}$ 开始。对每个 $\vec y\in\mathcal Y$，
先扫描每个 $\{0,1\}^{S_g}$ 找代表
$\bar x_g\in\mathcal X_g(y_g)$；若任一集合为空，
算法跳过该元组。否则初始化 $p=1_{\mathcal A}$。
为计算 $\Sigma_g(\vec y)$，初始化 $s=0_{\mathcal A}$
并再次扫描 $u\in\{0,1\}^{S_g}$。每当
\[
 C_gu=y_g,\qquad u_i=1\quad(i\in E\cap S_g),
\]
算法执行更新
\[
 s\leftarrow s+A_g\Bigl(u+\sum_{h\ne g}\bar x_h\Bigr).
\]
扫描结束时 $s=\Sigma_g(\vec y)$，算法置
$p\leftarrow ps$。全部组处理完后置 $Z\leftarrow Z+p$。
全部元组 $\vec y$ 处理完后，$Z=Z_E$。

对每个 $\vec y$，找代表的扫描与计算
$\Sigma_g(\vec y)$ 的扫描各至多检查
$\sum_g2^{|S_g|}$ 个布尔向量。这给出所述界。
检查一个向量的测量与固定位使用
$\poly(n,r,c,\log(D+2))$ 次位运算：矩阵每行的绝对值之
和至多 $D$，故其部分和有 $O(\log(D+2))$ 位。
每个向量上至多一次因子求值。
保留的状态是 $\vec y$、$(\bar x_1,\ldots,\bar x_r)$、$g$、
$u$、$s$、$p$ 与 $Z$。元组 $\vec y$ 含 $c$ 个 $[-D,D]$
中的整数，代表共含 $n$ 个比特，而 $s,p,Z\in\mathcal A$。
\end{proof}

为在固定 $t$ 处求值 \eqref{eq:general-count-polynomial}，
按 $\sigma_f(t)$ 把变换后的因子分组，
并在 $\mathcal A$ 中以 $E=\varnothing$ 应用该引理。
定义~\ref{def:layout} 给出其不变性假设，故该和为 $H(t)$。
与概览相同，两个指数成本是测量选取与一组内的布尔赋值。

\subsection{由固定布尔代入得到的 Taylor 系数}
\label{subsec:fixed-boolean-substitutions}
$Q$ 个值 $H(1),\ldots,H(Q)$ 未必确定 $H$。
令
\begin{equation}\label{eq:hasse-order}
 q=\left\lfloor\frac{d_H}{Q}\right\rfloor,\qquad J=q+1.
\end{equation}
由 $QJ>d_H$，只需计算
\begin{equation}\label{eq:taylor-truncation}
 \mathcal H_t(z)=\sum_{j=0}^{J-1}([z^j]H(t+z))z^j
 \qquad(t\in[Q]).
\end{equation}
Hermite 插值从这 $QJ$ 个 $\mathcal A$ 中系数恢复 $H$，
从而 $Z=H(0)$。事实上，次数小于 $QJ$ 的多项式的有理
Hermite 矩阵在 $\mathbb Q$ 上可逆，其逆通过有理线性组合
作用在 $\mathcal A$ 的元素上。

固定 $t$。$t+z$ 处的替代型不必尊重 $t$ 处的 design。
写
\[
 \Delta\vec P_{f,v}(z)=\vec P_{f,v}(t+z)-\vec P_{f,v}(t).
\]
每个分量都被 $z$ 整除。对固定的 $x\in\mathbb Q^n$，
引入独立未定元 $\vec X$，把每个输入替换为
\begin{equation}\label{eq:jet-substitution}
 \langle\vec P_{f,v}(t),x\rangle
       +\langle\Delta\vec P_{f,v}(z),\vec X\rangle.
\end{equation}
令
\begin{equation}\label{eq:auxiliary-polynomial}
 \mathcal G_x(\vec X,z)=\prod_{f\in\mathcal F}
 \Phi_f\!\left(
  (\langle\vec P_{f,v}(t),x\rangle+
   \langle\Delta\vec P_{f,v}(z),\vec X\rangle)_{v\in I_f}\right).
\end{equation}
置 $\vec X=x$ 给出
$\sum_x\mathcal G_x(x,z)=H(t+z)$。
改为固定 $\vec X$ 则使每个修正项与 $x$ 无关，
因此 $t$ 处的 design 仍然分离各因子。
此外，$\vec X$ 的每次出现都带一个 $z$ 因子，给出
\begin{equation}\label{eq:degree-cost}
 \deg_{\vec X}([z^j]\mathcal G_x(\vec X,z))\le j.
\end{equation}
下面的恒等式从固定布尔代入与对被求和赋值的限制恢复
代入 $\vec X=x$。

\begin{lemma}\label{lem:boolean-differences}
设 $q\ge0$ 为整数，$\mathcal A$ 是有单位元的交换环。
对每个 $x\in\{0,1\}^n$，设
$p_x(X_1,\ldots,X_n)$ 是 $\mathcal A$ 上次数至多 $q$
的多项式。则
\begin{align}
 \sum_{x\in\{0,1\}^n}p_x(x)
 &=\sum_{\substack{E\subseteq[n]\\|E|\le q}}
   \sum_{U\subseteq E}(-1)^{|E|-|U|}
   \sum_{\substack{x\in\{0,1\}^n\\x_i=1\ (i\in E)}}
                  p_x(\mathbf1_U).
 \label{eq:boolean-differences}
\end{align}
这里 $\mathbf1_U$ 是 $U$ 的指示向量。
\end{lemma}
\begin{proof}
对固定的 $x$，把 $X_i$ 的每个正幂替换为 $X_i$。
这保持布尔求值不变，且不增加次数。
由子集上的 M\"obius 反演，所得多重线性多项式有展开
\[
 \widetilde p_x(\vec{X})=\sum_{E\subseteq[n]}c_{x,E}\prod_{i\in E}X_i,
 \qquad
 c_{x,E}=\sum_{U\subseteq E}(-1)^{|E|-|U|}p_x(\mathbf1_U).
\]
次数界给出：只要 $|E|>q$ 就有 $c_{x,E}=0$。
置 $\vec{X}=x$ 并对 $x$ 求和即得所述恒等式：
单项 $\prod_{i\in E}x_i$ 恰选出对所有 $i\in E$
满足 $x_i=1$ 的赋值。
\end{proof}

对每个 $j<J$，把引理~\ref{lem:boolean-differences}
应用于 $\mathcal A$ 上的 $[z^j]\mathcal G_x$。
全部系数可以用同样的对 $U\subseteq E$、$|E|\le q$，
故在
\[
 \mathcal R=\mathcal A[z]/(z^J)
\]
中一并计算。对固定的对 $(E,U)$，令 $\Phi_{f,U}(x;z)$
为 \eqref{eq:auxiliary-polynomial} 中因子 $f$
在置 $\vec X=\mathbf1_U$ 后的形式。
按指派把这些因子分组：
\begin{equation}\label{eq:group-polynomial}
 A_{g,U}(x;z)=\prod_{f:\sigma_f(t)=g}\Phi_{f,U}(x;z).
\end{equation}
其乘积为 $\mathcal G_x(\mathbf1_U,z)$。
对 $h\ne g$，把 $x$ 加上 $\ker C_{t,h}$ 中的向量，
既不改变这些因子的原输入，也不改变其固定修正项。
因此引理~\ref{lem:measurement-separation} 通过检查至多
\begin{equation}\label{eq:boolean-enumeration}
 2r(2D+1)^{2\tau n}2^{2n/r}
\end{equation}
个布尔向量，计算 $x_i=1$（$i\in E$）的受限和。
合并全部对给出
\begin{equation}\label{eq:full-taylor-sum}
 H(t+z)\equiv
 \sum_{\substack{E\subseteq[n]\\|E|\le q}}
 \sum_{U\subseteq E}(-1)^{|E|-|U|}
 \sum_{\substack{x\in\{0,1\}^n\\x_i=1\ (i\in E)}}
       \mathcal G_x(\mathbf1_U,z)\pmod{z^J}.
\end{equation}
组乘积与组和在 $\mathcal R$ 中作成。结果是
$\mathcal H_t(z)$；在所有 $t$ 处重复即提供插值数据。
算法从不存储关于全部 $n$ 个未定元的多项式：
它先置 $\vec X=\mathbf1_U$，只存储 $\mathcal A$ 中的
$J$ 个系数。

\subsection{三个指数成本}
\label{subsec:three-costs}
对每个 $t\in[Q]$，满足 $|E|\le q$ 的对
$U\subseteq E\subseteq[n]$ 的个数是
\begin{equation}\label{eq:pair-count}
 C_q=\sum_{s=0}^{\min(q,n)}2^s\binom ns.
\end{equation}
特别地，$C_0=1$，且当 $q\ge n$ 时 $C_q=3^n$。
对每个 $0<\eta\le1$，
\begin{equation}\label{eq:finite-difference-cost}
 \eta^qC_q\le\sum_{s=0}^n\binom ns(2\eta)^s=(1+2\eta)^n,
\end{equation}
因为对每个保留的下标 $s\le q$ 都有
$\eta^q\le\eta^s$。

\begin{proposition}[公共计数例程]\label{prop:general-count}
设 $\mathcal A$ 是交换 $\mathbb Q$-代数。
设 \eqref{eq:general-product-sum} 中各因子的乘积
（在 $\mathcal A$ 上计算）满足次数界
\eqref{eq:general-product-degree}，且替代行、划分与测量
满足定义~\ref{def:layout}。取 \eqref{eq:hasse-order}
中的 $q,J$，该例程通过检查至多
\[
 2Qr\,C_q(2D+1)^{2\tau n}2^{2n/r}
\]
个布尔向量计算 $Z$。对每个向量，它在
$\mathcal A[z]/(z^J)$ 中至多求值 $|\mathcal F|$
个因子。每个向量上的额外运算与最终插值需要
$\mathcal A$ 与 $\mathbb Q$ 上关于参数
$n,|\mathcal F|,D,Q,R,J,r$ 与 scope 总长的多项式次运算。
除因子求值器自身的存储外，该例程只保留多项式个代数元素、
有理数与下标比特。
\end{proposition}
\begin{proof}
有 $Q$ 个点，每点 $C_q$ 个对，每个对至多
\eqref{eq:boolean-enumeration} 个布尔向量。
作每个组乘积至多用 $|\mathcal F|$ 次因子求值。
模 $z^J$ 的加法与乘法每次花费 $O(J^2)$ 次 $\mathcal A$
系数上的运算。上述 Taylor 恒等式与插值证明正确性。
依序处理点、对、测量元组与局部赋值，保留其当前下标与
累加器。连同 $QJ$ 个插值系数，这给出所述存储界。
位复杂度还取决于因子求值的成本与系数大小；
对稀疏多项式应用，这些在
第~\ref{subsec:one-query-cost} 节验证。
\end{proof}

当 $d=O(n)$ 时，$d_H=O(nR)$，$q=O(nR/Q)$。
在 \eqref{eq:finite-difference-cost} 中取 $\eta=R/Q$，
得
\[
 \log C_q\le q\log(Q/R)+n\log(1+2R/Q)
     =O\!\left(n\frac RQ\log\frac{2Q}{R}\right).
\]
因此，若因子求值与系数算术有多项式位成本，运行时间就是
多项式因子乘以
\begin{equation}\label{eq:three-costs}
 2^{O\left(n\left[
       \frac RQ\log\frac{2Q}{R}
       +\tau\log(2D+1)+\frac1r\right]\right)}.
\end{equation}
这恰是概览中的三项成本。当因子的 scope 与其最大出现次数
被 $\log n$ 的固定幂界定时，
引理~\ref{lem:incidence-parameters} 的参数选取使其和为
$O((\log\log n)^2/\sqrt{\log n})$。
该例程对更大的 $d$ 同样适用；此时在
命题~\ref{prop:general-count} 的界中保留实际值
$q=\lfloor dR/Q\rfloor$ 即可。

到目前为止我们假定 design 是给定的。
下一节从因子的 scope 构造它。
剩下两节先用所得例程计数给定多项式方程组的零点，
再为公式选取这样的方程组。
\section{构造低次替代型}\label{sec:construction}
计数例程对每对 $(f,v)$（因子 $f$ 读取变量 $v$）需要一个替代型。
对第~\ref{sec:setup} 节的应用，$f$ 是一个多项式方程，
其 scope 是其各单项变量集的并。本节只使用这些 scope。
我们构造定义~\ref{def:layout} 中的行 $\vec{P}_{f,v}(T)$
与测量矩阵 $C_{t,g}$，连同划分 $S_{t,g}$ 与因子指派
$\sigma_f$。主要困难在于保持 $\deg \vec{P}_{f,v}\le R$
且 $R$ 远小于 $Q$。$H$ 的次数界随 $R$ 线性增长，
因此小的 $R/Q$ 减少每点所需的 Taylor 系数个数。
在 $T=0,1,\ldots,Q$ 处独立选行再插值，一般给出次数 $Q$。
我们改为把它们在所有这些点上的取值一起构造出来。

构造先在辅助坐标的完整阵列上给出多项式行。
随后把每个「因子--变量」对放入该阵列的一个副本中。
多个位置可以含同一布尔变量 $x_v$；其取值总相等，
在计数和中不会成为独立变量。

\subsection{辅助坐标}\label{subsec:incidence-layout}
我们先指定更大的坐标集 $\mathcal B$——辅助行
$\vec{P}_u(T)$ 将作用其上——以及把实际的 $(f,v)$
放入其副本所需的关联界。
设 $(I_f)_{f\in\mathcal F}$ 为因子 scope，$I_f\subseteq[n]$，
写
\begin{equation}\label{eq:factor-incidences}
 \begin{gathered}
 \mathcal A=\{(f,v):f\in\mathcal F,\ v\in I_f\},\qquad M=|\mathcal A|,\\
 |I_f|\le w\quad(f\in\mathcal F),\qquad
 \#\{f:v\in I_f\}\le b\quad(v\in[n]).
 \end{gathered}
\end{equation}
于是 $w$ 界住一个因子 scope 中的变量数，$b$
界住含一个变量的因子数。scope 界 $w$ 不必界定因子多项式
的次数。scope 为空的因子是常数，不需要辅助位置；
在每个点把这样的因子指派到组一。
如有必要，对每个 $v\in[n]$ 附加一个 scope 为 $\{v\}$
的互异常数一因子。这既不改变乘积也不改变计数，
且使最大关联至多增加一。以 $b$ 记此次添加后的界，我们设
\[
 n\ge1,\qquad b,w\ge1,\qquad n\le M\le bn.
\]
本节始终设 $r\ge2$、$1\le R<Q$ 为整数，$0<\tau\le1/2$。

在每个点 $t\in[Q]$，变量 $x_v$ 将属于一个组，
而含 $v$ 的因子可能被指派到另一个组。我们将选取
\[
 a_v\in[r]^Q\quad(v\in[n]),\qquad
 d_f\in[r]^Q\quad(f\in\mathcal F),
\]
其中 $a_v(t)$ 指定含 $x_v$ 的组，$d_f(t)$ 指定指派给因子
$f$ 的组。位置 $(f,v)$ 将有标号 $(a_v,d_f)$。
特别地，所有含 $x_v$ 的位置都有同一个第一序列 $a_v$。

在选取 $a_v$ 与 $d_f$ 之前，先取全部可能标号的完整集合：
\begin{equation}\label{eq:auxiliary-labels}
 \mathcal B=([r]^Q)^2,\qquad u=(a_u,d_u),\qquad D=|\mathcal B|=r^{2Q}.
\end{equation}
辅助向量 $\zeta\in\mathbb Q^{\mathcal B}$
为每个对 $u$ 有一个有理坐标 $\zeta_u$。
标号 $u$ 在 $\zeta$ 变动时保持固定。
对 $t\in[Q]$ 与 $g\in[r]$，定义
\begin{equation}\label{eq:auxiliary-groups}
 S^{\mathcal B}_{t,g}=\{u\in\mathcal B:a_u(t)=g\}.
\end{equation}
对固定的 $t$，这些集合把 $\mathcal B$ 划分为 $r$ 个大小
$D/r$ 的集合。

对每个 $u\in\mathcal B$，我们寻求次数至多 $R-1$、
满足 $\vec{P}_u(0)=\delta_u$ 的多项式行
$\vec{P}_u(T)\in\mathbb Q[T]^{1\times\mathcal B}$。
在 $T=t$，取值 $\langle \vec{P}_u(t), \zeta\rangle$
应由
\[
 \zeta|_{S^{\mathcal B}_{t,d_u(t)}}
 \quad\text{与}\quad C^{\mathcal B}\zeta
\]
确定，其中 $C^{\mathcal B}$ 是一个与 $u,t$ 无关的整数矩阵，
分量在 $\{0,1,-1\}$ 中且行数很少。
稍后使用 $u_{f,v}=(a_v,d_f)$ 将使因子 $f$
的每个输入都用组 $d_f(t)$。$a_v,d_f$ 的选取与放入副本
分别推迟到第~\ref{subsec:choosing-labels}
与~\ref{subsec:embedding} 节。

\subsection{矩阵 $M_0$ 与 $C^{\mathcal B}$}
为构造行 $\vec{P}_u(T)$，我们寻求固定 $\ker C^{\mathcal B}$
并满足
\begin{equation}\label{eq:matrix-support-target}
 M_0[u,y]\ne0\quad\Longrightarrow\quad
 \bigl|\{s\in[Q]:a_y(s)\ne d_u(s)\}\bigr|\le R-1.
\end{equation}
的矩阵 $M_0\in\mathbb Q^{\mathcal B\times\mathcal B}$。
恒等式
\begin{equation}\label{eq:coordinate-decomposition}
 \zeta_u=((I-M_0)\zeta)_u+
          \sum_{y\in\mathcal B}M_0[u,y]\zeta_y
\end{equation}
解释了这两个要求。若 $M_0$ 固定 $\ker C^{\mathcal B}$，
则第一项由 $C^{\mathcal B}\zeta$ 确定；我们将使它与 $T$
无关。在和式中，把每个非零 $M_0[u,y]$ 乘以一个在零点为一、
在每个满足 $a_y(t)\ne d_u(t)$ 的 $t$ 处为零的多项式。
界 \eqref{eq:matrix-support-target} 容许次数至多 $R-1$。
于是在 $T=t$，和式将只用坐标 $y\in S^{\mathcal B}_{t,d_u(t)}$。

对 $s\in[Q]$，设 $\vartheta_s:\mathcal B\to\mathcal B$
交换 $a_u(s)$ 与 $d_u(s)$，其余分量不变。
用
\[
 (\Pi_s\zeta)_u=\zeta_{\vartheta_s(u)}.
\]
定义其置换矩阵 $\Pi_s$。
被交换的对 $\vartheta_s(u)$ 属于 $\mathcal B$，
因为该集合含 $[r]^Q$ 中序列的每一对。
矩阵 $\Pi_s$ 两两交换且满足 $\Pi_s^2=I$。其乘积
\[
 V_0=\prod_{s=1}^Q\Pi_s
\]
交换完整序列 $a_u$ 与 $d_u$。于是
$V_0[u,y]\ne0$ 蕴含 $a_y(s)=d_u(s)$ 对每个 $s$；
它以零错配满足 \eqref{eq:matrix-support-target}。
我们将修改 $V_0$，使所得矩阵同时固定一列短的整数测量的核。

定义
\begin{equation}\label{eq:swap-weight}
 W_0=\frac12\sum_{s=1}^Q(I-\Pi_s).
\end{equation}
在 $\Pi_s$ 的公共特征向量 $\psi$ 上，设 $j$ 为满足
$\Pi_s\psi=-\psi$ 的下标 $s$ 的个数。则
\[
 W_0\psi=j\psi,\qquad V_0\psi=(-1)^j\psi,
 \qquad V_0p(W_0)\psi=(-1)^jp(j)\psi
\]
对每个多项式 $p$ 成立。因此只要 $p(j)=(-1)^j$，
$V_0p(W_0)$ 就固定 $\psi$。
我们将构造一个显式公共特征基，其中具有给定 $j$
的向量比例恰为 $\Pr[W=j]$
（$W\sim\operatorname{Bin}(Q,(r-1)/(2r))$）。

对区间
$\mathcal I=\{A_0,\ldots,A_0+R-1\}\subseteq\{0,\ldots,Q\}$，
可以用次数至多 $R-1$ 的多项式强制等式
$f_{\mathcal I}(j)=(-1)^j$ 对 $j\in\mathcal I$ 成立。
我们使用显式插值公式
\begin{equation}\label{eq:swap-matrix-definition}
 f_{\mathcal I}(z)=(-1)^{A_0}\sum_{j=0}^{R-1}(-2)^j
                         \binom{z-A_0}{j},
 \qquad M_0=V_0f_{\mathcal I}(W_0),
\end{equation}
其中
\[
 \binom{z-A_0}{j}=\frac1{j!}\prod_{a=0}^{j-1}(z-A_0-a)
\]
且空乘积为一。对 $0\le q<R$，二项式定理给出
$f_{\mathcal I}(A_0+q)=(-1)^{A_0}(1-2)^q=(-1)^{A_0+q}$。

\begin{lemma}\label{lem:swaps-matrix}
对如上的 $\mathcal B$、$\mathcal I$ 与 $M_0$，
存在矩阵 $C^{\mathcal B}\in\{0,1,-1\}^{[c_{\mathcal B}]\times\mathcal B}$，
使
\[
 c_{\mathcal B}\le D\Pr[W\notin\mathcal I],\qquad
 W\sim\operatorname{Bin}\!\left(Q,\frac{r-1}{2r}\right),
\]
且对每个 $\xi\in\ker C^{\mathcal B}$ 有
$M_0\xi=\xi$。
\end{lemma}
\begin{proof}
对 $(\alpha,\delta)\in[r]^2$，设 $e_{\alpha\delta}$ 为该
有序对的指示函数。交换其两个分量有正交特征基
\[
 e_{\alpha\alpha}\quad(\alpha\in[r]),\qquad
 e_{\alpha\delta}+e_{\delta\alpha},\quad
 e_{\alpha\delta}-e_{\delta\alpha}\quad(1\le\alpha<\delta\le r).
\]
前两类特征值为 $+1$，第三类为 $-1$。
特征值 $+1$ 的向量有 $r(r+1)/2$ 个，$-1$ 的有
$r(r-1)/2$ 个。

对每个 $s\in[Q]$，选取这些函数之一
$\psi_s:[r]^2\to\{0,1,-1\}$，并定义
\[
 \psi(u)=\prod_{s=1}^Q\psi_s(a_u(s),d_u(s))
 \qquad(u\in\mathcal B).
\]
$D=r^{2Q}$ 种选取给出非零正交向量：其内积分解为各
$\psi_s$ 内积的乘积。它们构成 $\mathbb Q^{\mathcal B}$
的一组分量在 $\{0,1,-1\}$ 中的基。
若恰有 $j$ 个所选 $\psi_s$ 的特征值为 $-1$，则
$W_0\psi=j\psi$、$V_0\psi=(-1)^j\psi$。
故只要 $j\in\mathcal I$，$M_0$ 就固定 $\psi$。

为使 $\ker C^{\mathcal B}$ 成为被 $M_0$ 固定的空间，
取 $C^{\mathcal B}$ 的行为 $j\notin\mathcal I$
的那些基向量的转置。每个位置上有比例 $(r-1)/(2r)$
的选取特征值为 $-1$，且各位置独立。
因此作为行保留的 $D$ 个基向量的比例是
$\Pr[W\notin\mathcal I]$，证得行数界。
若 $C^{\mathcal B}\xi=\vec{0}$，则 $\xi$
与每个保留的基向量正交，故是 $j\in\mathcal I$
的那些基向量的线性组合，而它们全部被 $M_0$ 固定。
\end{proof}

对界 $\epsilon>0$，条件
\begin{equation}\label{eq:incidence-tail}
 \Pr\!\left[\operatorname{Bin}\!\left(Q,\frac{r-1}{2r}\right)
                         \notin\mathcal I\right]\le\epsilon
\end{equation}
给出 $c_{\mathcal B}\le\epsilon D$。
放置因子输入后，我们将取 $\epsilon=\tau/(rb)$。
接下来证明 $M_0$ 的另一个要求。

\begin{lemma}\label{lem:swaps-support}
对所有 $u,y\in\mathcal B$，若 $M_0[u,y]\ne0$，则
\[
 \bigl|\{s\in[Q]:a_y(s)\ne d_u(s)\}\bigr|\le R-1.
\]
\end{lemma}
\begin{proof}
在两两交换的矩阵 $\Pi_s$ 中展开 $f_{\mathcal I}(W_0)$。
由 $\deg f_{\mathcal I}\le R-1$ 与 $\Pi_s^2=I$，
每个所得单项是某个满足 $|J|\le R-1$ 的
$J\subseteq[Q]$ 对应的 $\prod_{s\in J}\Pi_s$
的标量倍。乘以 $V_0$ 给出
\[
 V_0\prod_{s\in J}\Pi_s=\prod_{s\notin J}\Pi_s.
\]
该置换矩阵的 $(u,y)$ 分量非零要求
$a_y(s)=d_u(s)$ 对每个 $s\notin J$。
其余位置至多 $R-1$ 个。若不一致的位置多于 $R-1$ 个，
则每个被加项的 $(u,y)$ 分量为零，
故 $M_0[u,y]=0$。
\end{proof}

\subsection{多项式行}
现在用 $M_0$ 构造第~\ref{subsec:incidence-layout} 节
所求的辅助行 $\vec{P}_u(T)$。
引理~\ref{lem:swaps-matrix} 使 $(I-M_0)\zeta$
成为 $C^{\mathcal B}\zeta$ 的函数。
引理~\ref{lem:swaps-support} 界住非零项
$M_0[u,y]\zeta_y$ 使用 $S^{\mathcal B}_{t,d_u(t)}$
之外坐标的点数。现在把该项乘以在这些点为零、在零点为
一的多项式。

对 $M_0[u,y]\ne0$，定义
\begin{equation}\label{eq:coefficient-correction}
 E_{u,y}=\{s\in[Q]:a_y(s)\ne d_u(s)\},\qquad
 \varphi_{u,y}(T)=\prod_{s\in E_{u,y}}\left(1-\frac Ts\right).
\end{equation}
空乘积为一。则 $\deg\varphi_{u,y}\le R-1$、
$\varphi_{u,y}(0)=1$，且只要 $a_y(t)\ne d_u(t)$ 就有
$\varphi_{u,y}(t)=0$。定义
\begin{equation}\label{eq:swaps-row}
 \vec{P}_u(T)=\delta_u(I-M_0)+
 \sum_{\substack{y\in\mathcal B\\M_0[u,y]\ne0}}
 M_0[u,y]\varphi_{u,y}(T)\delta_y.
\end{equation}
该行作为 $\zeta$ 的线性型作用，其系数是单个参数 $T$
的多项式。

\begin{lemma}\label{lem:swaps-forms}
\eqref{eq:swaps-row} 中的行次数至多 $R-1$，
且对每个 $u\in\mathcal B$ 满足 $\vec{P}_u(0)=\delta_u$。
对 $t\in[Q]$ 与 $g\in[r]$，设 $C^{\mathcal B}_{t,g}$
为 $C^{\mathcal B}$ 限制到列 $S^{\mathcal B}_{t,g}$。
对每个满足 $g\ne d_u(t)$ 的 $u,t,g$ 与每个
$\xi\in\mathbb Q^{S^{\mathcal B}_{t,g}}$，其补零延拓
$\widetilde\xi$ 满足
\[
 C^{\mathcal B}_{t,g}\xi=\vec{0}
 \quad\Longrightarrow\quad \langle \vec{P}_u(t), \widetilde\xi\rangle=0.
\]
更一般地，$\langle \vec{P}_u(t), \zeta\rangle$
由 $\zeta|_{S^{\mathcal B}_{t,d_u(t)}}$ 与
$C^{\mathcal B}\zeta$ 确定。
\end{lemma}
\begin{proof}
次数界来自 $\varphi_{u,y}$ 的次数界。在零点，
\[
 \vec{P}_u(0)=\delta_u(I-M_0)+\delta_uM_0=\delta_u.
\]
把 \eqref{eq:swaps-row} 作用于 $\zeta$ 给出
\begin{equation}\label{eq:swaps-form-decomposition}
 \langle \vec{P}_u(T), \zeta\rangle=((I-M_0)\zeta)_u+
 \sum_{\substack{y\in\mathcal B\\M_0[u,y]\ne0}}
 M_0[u,y]\varphi_{u,y}(T)\zeta_y.
\end{equation}
在 $T=t$，所有满足 $a_y(t)\ne d_u(t)$ 的被加项为零。
若 $\widetilde\xi$ 支撑在 $S^{\mathcal B}_{t,g}$
（$g\ne d_u(t)$）上，则每个剩余被加项在 $\widetilde\xi$
上为零。假设 $C^{\mathcal B}_{t,g}\xi=\vec{0}$
也给出 $C^{\mathcal B}\widetilde\xi=\vec{0}$，
故由引理~\ref{lem:swaps-matrix} 有
$(I-M_0)\widetilde\xi=\vec{0}$。
这就证明了该蕴含。

对最后的断言，在 $C^{\mathcal B}$ 下测量相同的两个向量
相差其核的一个元素。于是 $(I-M_0)\zeta$
由 $C^{\mathcal B}\zeta$ 确定。
\eqref{eq:swaps-form-decomposition} 中在 $T=t$
的其余项只用 $\zeta|_{S^{\mathcal B}_{t,d_u(t)}}$。
\end{proof}

固定 $t$ 并记 $\zeta_g=\zeta|_{S^{\mathcal B}_{t,g}}$。
$C^{\mathcal B}$ 的限制满足
\begin{equation}\label{eq:auxiliary-measurements}
 C^{\mathcal B}\zeta=\sum_{g=1}^r C^{\mathcal B}_{t,g}\zeta_g.
\end{equation}
因此固定每个向量 $C^{\mathcal B}_{t,g}\zeta_g$
既固定 $C^{\mathcal B}\zeta$，又对 $r$ 个辅助坐标组
施加各自的条件。当全部 $C^{\mathcal B}_{t,g}\zeta_g$
固定时，$\langle \vec{P}_u(t), \zeta\rangle$
由 $\zeta|_{S^{\mathcal B}_{t,d_u(t)}}$ 确定。
把 $(f,v)$ 放到标号 $(a_v,d_f)$ 的坐标处，
我们将在实际布尔变量上得到相应的分离。

\begin{remark}
我们演示 \eqref{eq:swaps-row} 如何把指派组之外的依赖
限制到一个线性测量。取 $Q=R=2$（在运行时间界所用
$R<Q$ 体制之外）与 $r=2$。
用 $\{0,1\}$ 作组下标，取 $\mathcal I=\{0,1\}$，
故 $f_{\mathcal I}(z)=1-2z$。
从标号 $(00,11)$ 出发，两次交换生成四个坐标：
\[
 \zeta_0=\zeta_{(00,11)},\quad \zeta_1=\zeta_{(10,01)},\quad
 \zeta_2=\zeta_{(01,10)},\quad \zeta_{12}=\zeta_{(11,00)}.
\]
两次交换都保持这四个坐标，故可把计算限制在其上。
在此限制上，$M_0=\Pi_1+\Pi_2-\Pi_1\Pi_2$
固定
\[
 c(\zeta)=\zeta_0-\zeta_1-\zeta_2+\zeta_{12}
\]
的核。以 $u=(00,11)$ 为下标的行给出
\[
 \langle \vec{P}_u(T), \zeta\rangle=\zeta_0-\frac T2\zeta_1-T\zeta_2.
\]
它在零点等于 $\zeta_0$，而
\[
 \langle \vec{P}_u(1), \zeta\rangle=c(\zeta)+\frac12\zeta_1-\zeta_{12},\qquad
 \langle \vec{P}_u(2), \zeta\rangle=c(\zeta)-\zeta_2-\zeta_{12}.
\]
在 $t=1$，$\zeta_1,\zeta_{12}$ 的标号满足
$a_y(1)=1=d_u(1)$。在 $t=2$，$\zeta_2,\zeta_{12}$
的标号满足 $a_y(2)=1=d_u(2)$。
因此固定 $c(\zeta)$ 后，$\langle \vec{P}_u(t), \zeta\rangle$
只依赖其指派组中的坐标。
这里 $c(\zeta)$ 在四坐标限制上扮演 $C^{\mathcal B}\zeta$
的角色。完整构造必须对每个 $u\in\mathcal B$
以至多 $\epsilon D$ 个测量行给出该性质。
\end{remark}

\subsection{选取变量组与因子组}\label{subsec:choosing-labels}
现在从第~\ref{subsec:incidence-layout} 节选取族
$(a_v)_{v\in[n]}$ 与 $(d_f)_{f\in\mathcal F}$，并置
\[
 u_{f,v}=(a_v,d_f)\qquad((f,v)\in\mathcal A).
\]
这些选取有两个要求。第一，若某标号出现 $c_u$ 次，
其关联在辅助阵列副本中需要 $c_u$ 个不同位置。
因此所需副本数为
\[
 c_u=\#\{(f,v)\in\mathcal A:u_{f,v}=u\},\qquad
 h=\max_{u\in\mathcal B}c_u.
\]
第二，实际变量组 $\{v:a_v(t)=g\}$ 必须小。
只界定 $h$ 界住的是一组内的关联个数，
而非其布尔变量个数。下面的引理同时控制两个量。

\begin{lemma}\label{lem:incidence-labels}
设 $(I_f)_{f\in\mathcal F}$ 满足
\eqref{eq:factor-incidences}，$n\le M\le bn$，
并对整数 $r\ge2$、$Q\ge1$ 置 $D=r^{2Q}$。设
\begin{equation}\label{eq:incidence-balance-hypotheses}
 2(b+w)D^2\le M,\qquad 2Qr^2\le n.
\end{equation}
可以确定性选取 $a_v,d_f\in[r]^Q$
（对 scope 为空的因子取 $d_f(t)=1$），使
\[
 \begin{aligned}
 \max_{u\in\mathcal B}\#\{(f,v)\in\mathcal A:(a_v,d_f)=u\}
 &\le\frac{2M}{D},\\
 \#\{v\in[n]:a_v(t)=g\}&\le\frac{2n}{r}
 \qquad(t\in[Q],\ g\in[r]).
 \end{aligned}
\]
构造花费 $\poly(n,M,|\mathcal F|,D,Q,r,b,w)$ 次位运算。
\end{lemma}
\begin{proof}
我们先估计随机标号的效果，然后再确定性固定其取值。
在此估计中，所有 $a_v$ 与所有 $I_f$ 非空的 $d_f$
独立均匀取自 $[r]^Q$。对空的 $I_f$，恒置 $d_f$ 为一。
每个关联标号 $u_{f,v}=(a_v,d_f)$ 均匀分布于 $\mathcal B$。
置 $n_{t,g}=\#\{v:a_v(t)=g\}$，考虑
\[
 \Phi=\sum_{u\in\mathcal B}(c_u-M/D)^2,
 \qquad
 \Xi=\sum_{t=1}^Q\sum_{g=1}^r(n_{t,g}-n/r)^2.
\]
其作为指示子和的展开为
\[
 \begin{aligned}
 \Phi&=\sum_{(f,v),(f',v')\in\mathcal A}
       \mathds{1}_{\{(a_v,d_f)=(a_{v'},d_{f'})\}}-\frac{M^2}{D},\\
 \Xi&=\sum_{t=1}^Q\sum_{v,v'\in[n]}
       \mathds{1}_{\{a_v(t)=a_{v'}(t)\}}-\frac{Qn^2}{r}.
 \end{aligned}
\]
两个关联标号除非共享因子或变量，否则独立。
对每个 $(f,v)$，至多有 $b+w$ 个关联共享 $f$ 或 $v$
（含 $(f,v)$ 自身）。独立对的相等概率为 $1/D$；
其余对的概率至多为一。故
\[
 \mathbb E\Phi\le M(b+w).
\]
对每个 $t$ 与 $v\ne v'$，$a_v(t),a_{v'}(t)$
独立均匀取自 $[r]$。第二个展开因此给出
\[
 \mathbb E\Xi=Qn(1-1/r)\le Qn.
\]

为确定性获得标号，我们逐个固定 $a_v$ 与非空 $I_f$
的 $d_f$。每步选取一个不增加
\[
 \frac{\Phi}{M(b+w)}+\frac{\Xi}{Qn}
\]
的条件期望的取值。这样的选取存在，因为当前期望是下一个
$a_v$ 或 $d_f$ 各可能取值上的平均。
最终值至多为一的两倍。由于两个被加项都非负，
\[
 c_u\le\frac MD+\sqrt{2M(b+w)},\qquad
 n_{t,g}\le\frac nr+\sqrt{2Qn}.
\]
\eqref{eq:incidence-balance-hypotheses} 的两个不等式
使两个根号项分别至多 $M/D$ 与 $n/r$。

下面的枚举在 $\poly(n,M,|\mathcal F|,D,Q,r,b,w)$
次位运算内计算条件期望。
置 $K=r^Q$，则 $D=K^2$。$\Phi$ 展开中的每个指示子至多
依赖四个独立选取 $a_v,d_f,a_{v'},d_{f'}$；
重复的下标只贡献一个选取。$\Xi$ 中的每个指示子至多依赖
$a_v,a_{v'}$。在任何部分固定之后，对一个指示子中剩余
选取枚举至多 $K^4=D^2$ 个赋值。其条件概率的分母整除
$K^4$。对 $M^2+Qn^2$ 个指示子求和即算出两个期望。
计入归一化 $M(b+w)$ 与 $Qn$，这些计算中的全部有理值的
分子与分母都有
\[
 O\!\left(\log(2D)+\log(2M)+\log(2n)+\log(2Q)+\log(2+b+w)\right)
\]
位。需要固定的选取至多 $n+|\mathcal F|$ 个，
每个有 $K$ 个候选值。由 $K^4=D^2$，
整个构造花费 $\poly(n,M,|\mathcal F|,D,Q,r,b,w)$
次位运算。
\end{proof}

\subsection{把变量副本放入辅助阵列}\label{subsec:embedding}
现在把辅助行 $\vec{P}_u(T)$ 与矩阵 $C^{\mathcal B}$
转移到定义~\ref{def:layout} 所要求的、原变量上的行
$\vec{P}_{f,v}(T)$ 与矩阵 $C_{t,g}$。
设 \eqref{eq:incidence-balance-hypotheses} 成立并应用
引理~\ref{lem:incidence-labels}。
以 $h=\max_u c_u$，用 $[h]\times\mathcal B$
为辅助坐标的 $h$ 份副本编下标。
对每个 $(f,v)\in\mathcal A$，指派 $c(f,v)\in[h]$
使所有对 $(c(f,v),u_{f,v})$ 互异。
这是可能的，因为没有标号出现超过 $h$ 次，且
\[
 hD\le2M\le2bn.
\]
定义线性映射
$\iota:\mathbb Q^n\longrightarrow\mathbb Q^{[h]\times\mathcal B}$：
\[
 (\iota x)_{c,u}=
 \begin{cases}
 x_v,&(c,u)=(c(f,v),u_{f,v})\text{ 对某个 }(f,v)\in\mathcal A,\\
 0,&\text{否则}.
 \end{cases}
\]
被占据的位置互异，故该定义无歧义。
不同位置可以含同一个 $x_v$，未用的位置固定为零。
对 $c\in[h]$，定义投影
\[
 \pi_c:\mathbb Q^{[h]\times\mathcal B}\longrightarrow\mathbb Q^{\mathcal B},
 \qquad (\pi_c\eta)_u=\eta_{c,u}.
\]
现在置
\begin{equation}\label{eq:actual-rows}
 \begin{gathered}
 \vec{P}_{f,v}(T)=\vec{P}_{u_{f,v}}(T)\pi_{c(f,v)}\iota,\\
 S_{t,g}=\{v\in[n]:a_v(t)=g\},\qquad \sigma_f(t)=d_f(t).
 \end{gathered}
\end{equation}
行 $\vec{P}_{f,v}(T)$ 作用在实际变量 $x\in\mathbb Q^n$
上。在 $T=0$，它提取位置 $(c(f,v),u_{f,v})$
处的值，即 $x_v$。于是 $\vec{P}_{f,v}(0)=\delta_v$，
次数至多 $R-1$。引理~\ref{lem:incidence-labels}
给出 $|S_{t,g}|\le2n/r$。

用
\begin{equation}\label{eq:actual-measurements}
 \widetilde C=\operatorname{diag}
   \bigl(\underbrace{C^{\mathcal B},\ldots,C^{\mathcal B}}_{h\text{ 份副本}}\bigr),
 \qquad \widehat C=\widetilde C\iota,
 \qquad C_{t,g}=\widehat C|_{S_{t,g}}.
\end{equation}
定义测量矩阵。
该限制保留 $S_{t,g}$ 中的列与全部行（含零行）。
于是 $\widehat Cx$ 在每个副本的坐标从 $x$ 填充后
记录其辅助测量。
矩阵 $C_{t,g}$ 是整数的。其分量不必属于 $\{0,1,-1\}$，
因为含同一 $x_v$ 的多个位置上的系数被相加。
$\widetilde C$ 的每行至多 $D$ 个非零分量，
全为一或 $-1$。乘以 $\iota$ 在公共列中添加分量并丢弃
未用位置，故不会增加其绝对值之和。
因此 $C_{t,g}$ 的每行绝对行和至多 $D$，
其在布尔赋值上的值落在 $[-D,D]\cap\mathbb Z$ 中。

若 \eqref{eq:incidence-tail} 成立，则每个 $C_{t,g}$
有 $c_{t,g}=hc_{\mathcal B}$ 行。故
\begin{equation}\label{eq:incidence-measurement-count}
 \sum_{g=1}^r c_{t,g}=rh c_{\mathcal B}
 \le rh\epsilon D\le2r\epsilon M\le2r\epsilon bn.
\end{equation}
取 $\epsilon=\tau/(rb)$ 使其至多 $2\tau n$。

为验证分离性，固定 $t\in[Q]$、$g\in[r]$ 与满足
$g\ne\sigma_f(t)$ 的 $(f,v)\in\mathcal A$。
设 $\xi\in\mathbb Q^{S_{t,g}}$ 满足
$C_{t,g}\xi=\vec{0}$，补零延拓为
$\widetilde\xi\in\mathbb Q^n$。
\eqref{eq:actual-measurements} 给出
$\widetilde C\iota\widetilde\xi=\vec{0}$。
于是对每个 $c\in[h]$，
\[
 C^{\mathcal B}\pi_c\iota\widetilde\xi=\vec{0}.
\]
$\pi_c\iota\widetilde\xi$ 的每个非零坐标含一个满足
$a_{v'}(t)=g$ 的变量 $v'$。其辅助标号的第一序列为
$a_{v'}$，故该向量支撑在 $S^{\mathcal B}_{t,g}$ 上。
在副本 $c(f,v)$ 中以 $d_{u_{f,v}}(t)=d_f(t)=\sigma_f(t)$
应用引理~\ref{lem:swaps-forms}，给出
\[
 \langle \vec{P}_{f,v}(t), \widetilde\xi\rangle
 =\langle \vec{P}_{u_{f,v}}(t), \pi_{c(f,v)}\iota\widetilde\xi\rangle=0.
\]
这正是定义~\ref{def:layout} 中的核条件。
因子 $f$ 的全部输入使用同一指派组 $d_f(t)$，
而计数和仍恰遍历 $n$ 个布尔位。

\subsection{系数界与完成的构造}
\eqref{eq:actual-rows} 中的代入行
$\vec{P}_{f,v}(T)\in\mathbb Q[T]^{1\times n}$
的分量是带理系数的多项式。现在界定存储这些系数所需的
位数；这些界将在第~\ref{subsec:one-query-cost} 节
控制精确算术。我们先对 \eqref{eq:swaps-row} 定义的
辅助行 $\vec{P}_u(T)$ 给出公分母与系数界。
由 \eqref{eq:actual-rows} 从 $\vec{P}_u(T)$
转到 $\vec{P}_{f,v}(T)$ 只在公共列中添加系数或丢弃
未用列，故同一分母与绝对值之和的界仍然适用。

\begin{lemma}\label{lem:row-coefficients}
设 $d_{\mathrm{den}}=(R-1)!Q!$。则对每个 $u\in\mathcal B$
有 $d_{\mathrm{den}}\vec{P}_u(T)\in\mathbb Z[T]^{1\times D}$。
行 $d_{\mathrm{den}}\vec{P}_u(T)$ 中全部系数的绝对值之
和至多 $2^{O(Q\log Q)}$。对 $d_{\mathrm{den}}\vec{P}_{f,v}(T)$
同样成立。矩阵 $M_0,C^{\mathcal B}$ 与行 $\vec{P}_u(T)$
可在 $\poly(D,Q,r)$ 次位运算内构造。
\end{lemma}
\begin{proof}
对 $0\le j<R$，$f_{\mathcal I}(W_0)$ 的相应被加项为
\[
 (-1)^{A_0}(-2)^j\binom{W_0-A_0I}{j}
 =\frac{(-1)^{A_0+j}}{j!}
   \prod_{a=0}^{j-1}\bigl(2W_0-2(A_0+a)I\bigr).
\]
由 $2W_0=QI-\sum_s\Pi_s$，右端每个矩阵因子都是整数的。
且 $j!$ 整除 $(R-1)!$。
因此乘以置换矩阵 $V_0$ 表明 $(R-1)!M_0$ 是整数的。
$\varphi_{u,y}$ 的系数的分母整除
$\prod_{s\in E_{u,y}}s$，而后者整除 $Q!$。
\eqref{eq:swaps-row} 证明所述公分母。

对系数估计，用矩阵范数
$\|A\|_\infty=\max_i\sum_j|A[i,j]|$
与多项式范数 $\|\sum_j a_jT^j\|_1=\sum_j|a_j|$。
由 $\|W_0\|_\infty\le Q$ 与上式各因子满足
$0\le A_0+a\le Q$，每个 $2W_0-2(A_0+a)I$
的范数至多 $4Q$。于是
\[
 \|M_0\|_\infty\le R(4Q)^{R-1},\qquad
 \|\varphi_{u,y}\|_1
 \le\prod_{s\in E_{u,y}}(1+1/s)\le2^Q.
\]
行 $d_{\mathrm{den}}\vec{P}_u$ 中全部系数的绝对值之
和因此至多
\[
 d_{\mathrm{den}}\bigl(1+(1+2^Q)R(4Q)^{R-1}\bigr).
\]
由 $R<Q$，其对数为 $O(Q\log Q)$。
对 $u=u_{f,v}$，把 $d_{\mathrm{den}}\vec{P}_u(T)$
乘以 $\pi_{c(f,v)}\iota$ 在公共列中添加系数并丢弃未用列。
由 \eqref{eq:actual-rows}，所得行是
$d_{\mathrm{den}}\vec{P}_{f,v}(T)$，故满足同一界。

枚举引理~\ref{lem:swaps-matrix} 的特征基、求值
$f_{\mathcal I}(W_0)$、并作成多项式 $\varphi_{u,y}$
需要 $\poly(D,Q,r)$ 次算术运算。使用上述公分母，
这些矩阵与多项式计算中存储的系数的分子与分母都有
$O(Q\log Q)$ 位。因此构造花费 $\poly(D,Q,r)$
次位运算。
\end{proof}

构造有三个定量要求。
\eqref{eq:incidence-balance-hypotheses} 的两个不等式
界定辅助副本数与每组实际变量数。
取 $\epsilon=\tau/(rb)$ 的尾界经由
\eqref{eq:incidence-measurement-count} 给出总测量界。
我们汇总结果供计数算法使用。

\begin{lemma}\label{lem:layout}
设整数 $n,b,w\ge1$，$(I_f)_{f\in\mathcal F}$
为 $[n]$ 的规模至多 $w$ 的子集，
每个 $v\in[n]$ 至多属于 $b$ 个集合。
置 $M=\sum_f|I_f|$ 并设 $n\le M\le bn$。
设 $r\ge2$、$1\le R<Q$ 为整数，$0<\tau\le1/2$，
$D=r^{2Q}$。设
\[
 2(b+w)D^2\le M,\qquad 2Qr^2\le n,
\]
且区间
$\mathcal I=\{A_0,\ldots,A_0+R-1\}\subseteq\{0,\ldots,Q\}$
满足
\[
 \Pr[W\notin\mathcal I]\le\frac{\tau}{rb},\qquad
 W\sim\operatorname{Bin}\!\left(Q,\frac{r-1}{2r}\right).
\]
则可以在 $\poly(n,M,|\mathcal F|,D,Q,r,b,w)$ 次位运算内
确定性构造定义~\ref{def:layout} 中的全部数据。
\end{lemma}
\begin{proof}
引理~\ref{lem:swaps-matrix} 构造满足
$c_{\mathcal B}\le\tau D/(rb)$ 的 $C^{\mathcal B}$。
引理~\ref{lem:swaps-support} 与~\ref{lem:swaps-forms}
给出辅助多项式行。
引理~\ref{lem:incidence-labels} 选取变量组与因子组，
满足 $hD\le2M$ 与 $|S_{t,g}|\le2n/r$。
\eqref{eq:actual-rows} 与 \eqref{eq:actual-measurements}
中的映射随后给出满足 $\vec{P}_{f,v}(0)=\delta_v$、
次数至多 $R-1\le R$ 与所需核蕴含的行。
$C_{t,g}$ 的每行绝对行和至多 $D$，且
\eqref{eq:incidence-measurement-count} 对每个 $t\in[Q]$
给出 $\sum_{g=1}^r c_{t,g}\le2\tau n$。
引理~\ref{lem:row-coefficients} 提供公分母与系数界。
合起来，各构造花费 $\poly(n,M,|\mathcal F|,D,Q,r,b,w)$
次位运算，所得阵列占用
$\poly(n,M,|\mathcal F|,D,Q,r,b,w)$ 比特。
\end{proof}

\subsection{选取参数}\label{subsec:incidence-parameters}
在应用中，最大 scope $w$ 与含一个变量的因子最大数 $b$
都将被 $\log n$ 的固定幂界定。现在在该体制下选取满足
引理~\ref{lem:layout} 的参数。
计数成本的指数取决于三个量
\[
 \frac RQ\log\frac{2Q}{R},\qquad
 \tau\log(2D+1),\qquad \frac1r.
\]
它们分别对应恢复 Taylor 系数、枚举测量取值、
与在一组内枚举赋值。下面的选取同时界定三者。

\begin{lemma}\label{lem:incidence-parameters}
固定整数 $\beta\ge1$。设 $n\ge20$，
$\lambda=\lceil\log_2 n\rceil$、
$\nu=\lceil\log_2\lambda\rceil$，
并设 $b,w\le\lambda^\beta$、$n\le M\le bn$。
对所有足够大的 $n$，参数
\begin{equation}\label{eq:parameter-QR}
 r=\left\lfloor\frac{\sqrt\lambda}{\nu^2}\right\rfloor,
 \qquad Q=\left\lfloor\frac{\lambda}{128\nu}\right\rfloor,
 \qquad R=(\beta+3)\lceil\sqrt\lambda\rceil,
\end{equation}
\[
 \tau=\frac1{r\lambda},\qquad
 \epsilon=\frac1{br^2\lambda},\qquad D=r^{2Q},
\]
连同证明中指定的区间，满足引理~\ref{lem:layout}
的全部数值假设。它们还满足
\[
 R\ge32,\qquad 6R\le Q,\qquad D\le2^{\lambda/128},
\]
以及
\begin{equation}\label{eq:balanced-cost}
 \frac RQ\log\frac{2Q}{R}+\tau\log(2D+1)+\frac1r
 =O_\beta\!\left(\frac{\nu^2}{\sqrt\lambda}\right).
\end{equation}
给定数值 $n,M,b,w$，这些参数与数值假设可以在
$\poly_\beta(\log n)$ 次位运算内精确计算并检验。
对固定的 $\beta$，只有有限多个 $n$ 通不过这些检验。
\end{lemma}
\begin{proof}
对足够大的 $n$，参数的阶为
\[
 r=\Theta(\sqrt\lambda/\nu^2),\qquad
 Q=\Theta(\lambda/\nu),\qquad R=\Theta_\beta(\sqrt\lambda).
\]
特别地，最终有 $r\ge2$、$R\ge32$ 与 $6R\le Q$。
由 $\log_2 r\le\nu/2$，
\[
 \log_2 D=2Q\log_2 r
 \le2\frac{\lambda}{128\nu}\frac\nu2=\frac{\lambda}{128}.
\]
对足够大的 $n$，两个平衡不等式随之成立：
\[
 2(b+w)D^2\le4\lambda^\beta2^{\lambda/64}\le n\le M,
 \qquad 2Qr^2\le n.
\]
对第二个：$Q,r$ 被 $\lambda$ 的固定幂界定，
而 $n>2^{\lambda-1}$。第一个由同样的比较得出。

令
\[
 \mu=\frac{Q(r-1)}{2r},\qquad
 A_0=\left\lfloor\mu-\frac{R-1}{2}\right\rfloor,
 \qquad \mathcal I=\{A_0,\ldots,A_0+R-1\}.
\]
这里 $Q/4\le\mu\le Q/2$。由 $R\le Q/6$，区间落在
$\{0,\ldots,Q\}$ 内。其中心与 $\mu$ 相差小于一，
而 $R\ge32$ 保证它包含与 $\mu$ 距离在 $R/4$ 内的每个
整数。对 $W\sim\operatorname{Bin}(Q,(r-1)/(2r))$，
Hoeffding 不等式给出
\[
 \Pr[W\notin\mathcal I]
 \le2e^{-R^2/(8Q)}
 \le2e^{-16(\beta+3)^2\nu}
 \le2^{-(\beta+2)\nu}
 \le\lambda^{-(\beta+2)}
 \le\frac1{br^2\lambda}.
\]
最后一个不等式用 $b\le\lambda^\beta$ 与 $r^2\le\lambda$。
这证明了取 $\epsilon=\tau/(rb)$ 的
\eqref{eq:incidence-tail}。且一旦 $r\ge2$ 就有
$0<\tau\le1/2$。

最后，
\[
 \frac RQ=O_\beta(\nu/\sqrt\lambda),\qquad
 \log(2Q/R)=O(\nu),\qquad
 \tau\log(2D+1)=O(1/r).
\]
连同 $1/r=O(\nu^2/\sqrt\lambda)$，这些证明
\eqref{eq:balanced-cost}。

最后我们证明 \eqref{eq:parameter-QR}
中的参数可以在所述位界内精确构造。
算法找出 $\lambda=\lceil\log_2 n\rceil$
（即使 $n\le2^\lambda$ 的最小整数），
再以同样方式找出 $\nu=\lceil\log_2\lambda\rceil$。
它把 $r$ 取为满足 $r^2\nu^4\le\lambda$ 的最大整数，
用整数除法计算 $Q$，用整数平方根搜索找
$\lceil\sqrt\lambda\rceil$ 得 $R$。
若 $r<2$ 或 $Q\le R$，算法在此阶段报告失败。
否则它用整数或有理算术构造 $\tau,\epsilon,D$
与 $\mathcal I$ 的端点。然后检验
$\mathcal I\subseteq\{0,\ldots,Q\}$、两个平衡不等式、
$R\ge32$ 与 $6R\le Q$。
\eqref{eq:incidence-tail} 中的尾概率分母为 $(2r)^Q$，
故其所需界也可用整数算术检验。
这些计算花费 $\poly_\beta(\log n)$ 次位运算。
前述估计表明所有检验对足够大的 $n$ 都通过。
\end{proof}
\section{稀疏多项式方程公共零点的计数}\label{sec:setup}
现在证明定理~\ref{thm:overview-sparse-polynomials}。
输入是 $\mathbb F_2$ 上 $n$ 个布尔变量、$m_{\mathrm{test}}\le cn$
个多项式方程 $P_i(x)=0$ 构成的方程组 $\mathcal P$，
$k,c$ 固定。每个 $P_i$ 是至多 $(\log n)^c$ 个次数至多 $k$
的单项的列表。目标是整数
\begin{equation}\label{eq:common-zero-count}
 N(\mathcal P)=\#\{x\in\{0,1\}^n:P_i(x)=0\text{ 对每个 }i\}.
\end{equation}
首先我们构造每方程一个因子的乘积，其在布尔立方上的取值
记录公共零点指示子。困难在于使其总输入次数关于 $n$
保持线性。随后为这些因子构造替代型，并用
第~\ref{sec:counting} 节对乘积求和。
第~\ref{sec:hashes} 节将所得算法用于任意给定的方程组，
包括仍接受某公式不可满足赋值的方程组。

\subsection{单个方程的整数多项式}\label{subsec:rational-lift}
在布尔输入上，可把每个正幂 $x_v^j$ 替换为 $x_v$。
把重复单项按模二合并。设 $\mathcal M_i$ 是由此得到的
$P_i$ 的变量集列表，定义其整数提升为
\[
 L_i(\vec X)=\sum_{M\in\mathcal M_i}\prod_{v\in M}X_v .
\]
空单项为一，空和为零。$L_i(x)$ 统计取值为一的单项个数，
故 $P_i(x)=0$ 恰在 $L_i(x)$ 为偶数时成立。
写 $I_i=\bigcup_{M\in\mathcal M_i}M$。
虽然 $L_i$ 可读取 $k(\log n)^c$ 个变量，
其次数仍至多 $k$。

令 $\ell=n+1$。由 $0\le N(\mathcal P)\le2^n<2^\ell$，
只需得到其模 $2^\ell$ 的剩余。
我们先用一个模二的幂的整数多项式表示偶性指示子。
直接相乘 $m_{\mathrm{test}}$ 个这样的指示子只能得到二次
次数界；下一条引理的形式将使我们能\emph{跨全部方程}
控制次数。

令 $h(Y)=Y(Y-1)$。它在每个整数 $Y$ 处取偶数值。
把指示子写成 $h$ 的幂，就在乘积中为每个 $h$ 副本
暴露出一个二的因子。

\begin{lemma}\label{lem:rational-idempotent}
对每个整数 $\ell\ge1$，可以构造多项式
$e_0,\ldots,e_{\ell-1}\in\mathbb Z[Y]$（每个次数至多一），
使
\begin{equation}\label{eq:idempotent-expansion}
 p_\ell(Y)=\sum_{j=0}^{\ell-1}e_j(Y)h(Y)^j,
 \qquad
 p_\ell(y)\equiv\mathds{1}_{\{2\mid y\}}\pmod{2^\ell}
 \quad(y\in\mathbb Z).
\end{equation}
$e_j$ 的系数有 $\poly(\ell)$ 位，构造使用 $\poly(\ell)$
次位运算与 $\poly(\ell)$ 比特空间。
\end{lemma}
\begin{proof}
设 $F_\ell$ 为 Beigel 与 Tarui 的模放大多项式
~\cite[Section~2.3]{BT94}。Chan 与 Williams
把 $F_\ell$ 作用于子句违反指示子的和，
确定性计数可满足赋值~\cite[Theorem~2.5 and proof of Theorem~1.6]{CW21}。
我们需要偶奇偶性的指示子，故取其补：
\begin{equation}\label{eq:binary-indicator}
 p_\ell(Y)=1-F_\ell(Y)=(1-Y)^\ell
       \sum_{j=0}^{\ell-1}\binom{\ell+j-1}{j}Y^j.
\end{equation}
该有限和是 $(1-Y)^\ell$ 模 $Y^\ell$ 的逆。
事实上，$(1+Y+Y^2+\cdots)^\ell$ 中 $Y^j$ 的系数统计 $j$
分入 $\ell$ 部分的弱组合，等于
$\binom{\ell+j-1}{j}$。因此
\[
 p_\ell(Y)\equiv1\pmod{Y^\ell},
 \qquad
 p_\ell(Y)\equiv0\pmod{(1-Y)^\ell}.
\]
对偶数 $y$，第一个同余给出
$p_\ell(y)\equiv1\pmod{2^\ell}$。
对奇数 $y$，第二个给出 $p_\ell(y)\equiv0\pmod{2^\ell}$。

多项式 $p_\ell$ 有整系数且次数小于 $2\ell$。
用首一二项式 $h=Y^2-Y$ 逐次除之，把它唯一写成
$\sum_{j<\ell}e_jh^j$，其中余式 $e_j$ 是次数小于二的
整系数多项式。附录~\ref{app:arithmetic} 的
引理~\ref{lem:rational-lift-height} 界住其系数，
并证明该构造的 $\poly(\ell)$ 位运算界。
\end{proof}

\subsection{保持乘积次数线性}
完整乘积 $\prod_{i=1}^{m_{\mathrm{test}}} p_\ell(L_i(\vec{X}))$
的次数至多 $km_{\mathrm{test}}(2\ell-1)$。
每个因子 $h(y)$ 在整数自变量处贡献一个二的因子，
故含至少 $\ell$ 个这样的因子的项可被模 $2^\ell$ 丢弃。
这使我们能\emph{跨全部方程}只保留 $h$ 的总幂次小于
$\ell$ 的项。引入形式变量 $V$ 记录该幂次，对每个方程
使用同一个 $V$。其指数记录在整数自变量处整除相应项的
一个有保证的二的幂。它不是要被求和的附加布尔变量。
这里我们选取第~\ref{subsec:product-interface} 节的
系数代数：
\[
 \mathcal A_\ell=\mathbb Q[V]/(V^\ell).
\]
在该代数中每次乘法都丢弃 $s\ge\ell$ 的幂 $V^s$。
乘积作成后，用线性映射
\[
 \mathcal S_\ell\!\left(\sum_{s=0}^{\ell-1}b_sV^s\right)
 =\sum_{s=0}^{\ell-1}b_s
\]
把其保留系数相加。该映射不必保持乘积，
因此只在因子相乘之后使用。
对多项式系数 $b_s$ 我们沿用同一记号。

\begin{lemma}\label{lem:total-order}
设整数 $m_{\mathrm{test}}\ge0$、$\ell\ge1$。
对整数 $y_1,\ldots,y_{m_{\mathrm{test}}}$，置
\[
 B_i(V)=\sum_{j=0}^{\ell-1}e_j(y_i)h(y_i)^jV^j.
\]
则
\[
 \mathcal S_\ell\!\Bigl(\prod_{i=1}^{m_{\mathrm{test}}} B_i(V)\bmod V^\ell\Bigr)
 \equiv\prod_{i=1}^{m_{\mathrm{test}}}\mathds{1}_{\{2\mid y_i\}}\pmod{2^\ell}.
\]
对任何非负整数 $d_*$，若把 $y_i$ 换成总次数至多 $d_*$
的 $\mathbb Q$ 上多项式，则每个保留的 $V^s$
系数（从而其和）的总输入次数至多为
$d_*(m_{\mathrm{test}}+2(\ell-1))$。
\end{lemma}
\begin{proof}
截断之前，对 $V^s$ 系数有贡献的项形如
\[
 \prod_{i=1}^{m_{\mathrm{test}}} e_{j_i}(y_i)h(y_i)^{j_i},
 \qquad \sum_{i=1}^{m_{\mathrm{test}}} j_i=s.
\]
在整数自变量处它被 $2^s$ 整除。因此丢弃每个
$s\ge\ell$ 的系数后，系数和在模 $2^\ell$
下不变。不截断时，该和是 $\prod_i p_\ell(y_i)$，
由引理~\ref{lem:rational-idempotent} 有所需同余。

对多项式自变量，每个保留项满足
$\sum_i j_i\le\ell-1$。由 $\deg e_{j_i}\le1$
与 $\deg h=2$，其次数至多
\[
 \sum_{i=1}^{m_{\mathrm{test}}} d_*(1+2j_i)
 \le d_*(m_{\mathrm{test}}+2(\ell-1)).
\]
该界逐系数成立，对其和亦然。
\end{proof}

为说明截断为何使用跨全部方程的全序，取 $\ell=2$。则
\[
 p_2(Y)=1-3Y^2+2Y^3=(1-Y)+(2Y-1)h(Y).
\]
对两个方程，模 $V^2$ 相乘后再作 $\mathcal S_2$，
保留的是
\[
 (1-y_1)(1-y_2)
 +(2y_1-1)h(y_1)(1-y_2)
 +(1-y_1)(2y_2-1)h(y_2).
\]
其次数至多四。被略去的项
$(2y_1-1)(2y_2-1)h(y_1)h(y_2)$ 次数为六，
但在整数输入处被四整除。

\subsection{用给定的替代型计数}\label{subsec:count-polynomial}
对方程 $i$，向第~\ref{sec:counting} 节提供因子
\begin{equation}\label{eq:lifted-factor}
 \mathcal B_i((X_v)_{v\in I_i},V)
 =\sum_{j=0}^{\ell-1}e_j(L_i(\vec X))h(L_i(\vec X))^j V^j
 \quad\text{（在 }\mathcal A_\ell[(X_v)_{v\in I_i}]\text{ 中）}.
\end{equation}
在一个因子内部，$X_v$ 的每次出现都使用同一替代型；
不同因子可对 $X_v$ 使用不同的型。
即使每个因子得到各自独立的输入，$L_i$
在这些输入上的次数仍至多 $k$。
引理~\ref{lem:total-order} 把乘积
\emph{在 $\mathcal A_\ell$ 中}的输入次数界定为
$d=k(m_{\mathrm{test}}+2(\ell-1))$。
于是该乘积恰满足公共例程的次数假设
\eqref{eq:general-product-degree}。
添加 scope 为 $\{v\}$、取值为一的因子，
保证所有声明变量都出现在构造中。

设 $H(T)\in\mathcal A_\ell[T]$ 是这些因子的
\eqref{eq:general-count-polynomial}。其在 $T=0$
处的 $V$ 系数是整数，因为所有替换在零点都恢复其布尔
输入。由同一条引理，对其取值求和给出
\begin{equation}\label{eq:origin-residue}
 \deg_T H\le d_H:=kR(m_{\mathrm{test}}+2(\ell-1)),\qquad
 \mathcal S_\ell(H(0))\equiv N(\mathcal P)\pmod{2^\ell}.
\end{equation}
因此 $\mathcal S_\ell(H(0))$ 的最小非负剩余
就是精确计数。插值在有理数上进行；
该同余只在 $T=0$ 的布尔输入处需要。

\subsection{求值因子的成本}\label{subsec:one-query-cost}
\begin{proposition}\label{prop:count}
给定 $n$ 个变量上的次数至多 $k$ 的多项式方程列表，
以及对其因子满足定义~\ref{def:layout} 的 design。
设 $S\ge1$ 界住单项列表与因子 scope 的总长，
取任何满足 $2^\ell>2^n$ 的 $\ell$。
用 \eqref{eq:origin-residue} 的 $d_H$ 与
\eqref{eq:hasse-order} 的 $q,J$，
$N(\mathcal P)$ 可确定性地在
\[
 \poly(n,m_{\mathrm{test}},\ell,D,Q,R,J,r,k,S)\,
 C_q(2D+1)^{2\tau n}2^{2n/r}
\]
次位运算与 $\poly(n,m_{\mathrm{test}},\ell,D,Q,R,J,r,k,S)$
比特工作空间内计算。
\end{proposition}
\begin{proof}
对 $\mathcal A_\ell$ 中的系数应用
命题~\ref{prop:general-count}。在求值点 $t$，
例程使用 $\mathcal A_\ell[z]/(z^J)$，
故存储的元素恰有 $J\ell$ 个有理系数。
要在 \eqref{eq:jet-substitution} 的固定布尔代入
$\vec X=\mathbf1_U$ 处求值因子 $i$，
先把每个列出的单项通过相乘其至多 $k$ 个被代入输入求值，
再把结果相加得 $L_i$。把该值用于
\eqref{eq:lifted-factor}，逐次乘以 $h(L_i)$
得到其幂。这在所列参数上花费多项式次有理运算，
且从不枚举整个 scope $I_i$ 上的赋值。
多项式 $e_j$ 由引理~\ref{lem:rational-idempotent}
一次性算出。

对位复杂度，各行有公分母 $d_{\mathrm{den}}=(R-1)!Q!$。
每个被代入单项的分母整除 $\Delta=d_{\mathrm{den}}^k$，
每个提升因子的分母整除 $\Gamma=\Delta^{2\ell-1}$，
其乘积的分母整除 $\Omega=\Gamma^{m_{\mathrm{test}}}$。
于是指数多个被加项可在同一个分母上累加。
附录~\ref{app:arithmetic} 把这些和中的位长界定为
\begin{equation}\label{eq:sum-bit-bound}
 B_{\mathrm{sum}}=\poly(n,m_{\mathrm{test}},\ell,k,Q,\log(D+2),\log(S+2)).
\end{equation}
对每个 $0\le s<\ell$，我们从其 $QJ$ 个 Taylor 系数
插值有理多项式 $[V^s]H(T)$。
该附录验证此插值以及 \eqref{eq:origin-residue}
的最终整数剩余的多项式位成本。
$O(QJ\ell)$ 个存储系数、design 以及
命题~\ref{prop:general-count} 中的依序枚举
给出所述空间界。
\end{proof}

\subsection{构造替代型并完成证明}\label{subsec:sparse-queries}
每个稀疏多项式读取的变量很少，但一个变量可能出现在许多
方程中。我们先固定出现次数最多的少数变量的取值。
在所得的每个分支上，即可构造
命题~\ref{prop:count} 所需的 design。

\begin{proof}[定理~\ref{thm:overview-sparse-polynomials} 的证明]
设 $\mathcal P$ 为给定的多项式方程组。
空方程列表给出 $2^n$ 个公共零点；常值方程 $1=0$
给出零个。丢弃方程 $0=0$。为每个声明变量 $v$
附加一个 scope 为 $\{v\}$、取值为一的中性因子，并写
\[
 M=\sum_{f\in\mathcal F}|I_f|
\]
表示「因子--变量」关联总数。
输入界蕴含 $M=n(\log n)^{O_{k,c}(1)}$，
且每个 scope 有 $(\log n)^{O_{k,c}(1)}$ 个变量。
不过某个特定变量可能出现在许多 scope 中。

令 $\lambda=\lceil\log_2 n\rceil$ 与
\[
 b_0=\left\lceil\frac Mn\right\rceil\lceil\sqrt\lambda\rceil.
\]
设 $S\subseteq[n]$ 为属于多于 $b_0$ 个 scope 的变量。
对每个 $b\in\{0,1\}^S$，固定 $x|_S=b$
并计数 $[n]\setminus S$ 上的补全；我们称这个受限问题为
一个分支。由 $|S|\le M/b_0\le n/\sqrt\lambda$，
分支至多 $2^{n/\sqrt\lambda}$ 个。
每个分支有相同的自由变量数 $n'$，且
\[
 n'\ge n-\frac n{\sqrt\lambda}\ge\frac n2
 \qquad(n\ge20).
\]
在 $\mathbb F_2$ 上把固定位代入每个 $P_i$，
并在作整数提升前按模二合并重复单项。
次数与单项数都不增加。常值方程可再次立即处理。
各分支划分布尔立方，故其公共零点计数相加为
$N(\mathcal P)$。

为应用第~\ref{sec:construction} 节，
从 scope 中移除固定变量，并把自由变量重标记为 $[n']$。
记 $I'_f\subseteq[n']$ 为受限 scope，
$M'=\sum_f|I'_f|$ 为其关联总数。
设 $b$ 为自由变量的最大因子关联数，
$w=\max_f|I'_f|$。则 $b\le b_0$，
且 $b$ 与 $w$ 都被 $\log n'$ 的固定幂界定。
剩余的中性因子保证 $M'\ge n'$。
令
\[
 \lambda'=\lceil\log_2 n'\rceil,\qquad
 \nu'=\lceil\log_2\lambda'\rceil,
\]
并选取满足 $b,w\le(\lambda')^\beta$ 的最小整数
$\beta\ge1$。对固定的 $k,c$，对一切足够大的 $n$，
$\beta$ 的可能取值被一个常数界定。
我们把引理~\ref{lem:incidence-parameters} 应用于
$[n']$ 上的受限 scope $(I'_f)$（关联计数 $M'$），取
\[
 r=\left\lfloor\frac{\sqrt{\lambda'}}{(\nu')^2}\right\rfloor,
 \qquad Q=\left\lfloor\frac{\lambda'}{128\nu'}\right\rfloor,
 \qquad R=(\beta+3)\lceil\sqrt{\lambda'}\rceil.
\]
若 $r<2$ 或 $Q\le R$，则先使用下文的穷举搜索退路，
再作 $\tau$、$\epsilon$ 与奇偶插值区间 $\mathcal I$。
否则置
\[
 \tau=\frac1{r\lambda'},\qquad
 \epsilon=\frac1{br^2\lambda'},\qquad D=r^{2Q},
\]
并执行该引理其余的可行性检验。
对一切足够大的 $n$，引理~\ref{lem:layout}
提供所需的行与测量矩阵，满足
\[
 r=\Theta\!\left(\frac{\sqrt{\log n'}}{(\log\log n')^2}\right),
 \qquad Q=\Theta\!\left(\frac{\log n'}{\log\log n'}\right),
 \qquad R=O_{k,c}(\sqrt{\log n'}),
 \qquad D\le(2n')^{1/128}.
\]

在每个分支上都可以使用原精度 $\ell=n+1$，
因为即使不受限的计数也小于 $2^\ell$。
由 $n'\ge n/2$，$\ell$ 与 $m_{\mathrm{test}}$
都是 $O_c(n')$。于是 $d_H=O_{k,c}(n'R)$、
$J=O_{k,c}(n'R/Q)+1$，而 \eqref{eq:three-costs}
在维数 $n'$ 下的指数为
\[
 O_{k,c}\!\left(\frac{n'(\log\log n')^2}{\sqrt{\log n'}}\right)
 =O_{k,c}\!\left(\frac{n(\log\log n)^2}{\sqrt{\log n}}\right).
\]
全部代入数据、系数阵列与插值方程组的表示总长为
$n^{O_{k,c}(1)}$ 比特。把每分支的运行时间乘以至多
$2^{n/\sqrt\lambda}$，所述指数界不变，
故把分支答案相加给出所需的运行时间。

为对每个输入规模实现该构造，算法用整数算术检验其数值
假设。除上面的初步守卫外，还检验
引理~\ref{lem:incidence-labels} 的
$2(b+w)D^2\le M'$ 与 $2Qr^2\le n'$、
$\mathcal I$ 的端点、$R\ge32$、$6R\le Q$，
以及可精确检验的 \eqref{eq:incidence-tail}
中的二项尾界。若任何检验失败，算法放弃分支计算，
改为在全部 $2^n$ 个赋值上直接检验原方程。
对固定的 $k,c$，刚证明的界与
引理~\ref{lem:incidence-parameters} 表明失败只可能对
有限多个 $n$ 发生。因此该退路被吸收进定理的常数中。

为在所述空间界内合并分支计数，我们一次处理一个分支。
初始化整数累加器 $Z=0$，把每个分支的答案加进去。
在每个分支上，构造与
命题~\ref{prop:count} 使用 $n^{O_{k,c}(1)}$
比特空间。不交性保证全程 $0\le Z\le2^n$，
最终值为 $N(\mathcal P)$，如所断言。
\end{proof}
\section{把 \texorpdfstring{$\#k$-SAT}{\#k-SAT} 归约到公共零点计数}\label{sec:hashes}\label{subsec:sparse-hashes}
给定至多 $Bn$ 个子句的稀疏 $k$-CNF $F$（$B$ 为 $n$
的多对数），我们在其原 $n$ 个变量上构造 $O(n)$
个稀疏多项式方程，使其公共零点恰为其可满足赋值。
每子句一个方程可能需要 $Bn$ 个方程。我们改为让每个方程
检验一列短子句之违反的奇偶性。我们用
定理~\ref{thm:overview-sparse-polynomials} 既选取这些方程，
也得到最终计数。

设 $m$ 为编入下标的子句数，$v_e(x)$ 为子句 $e$
的违反位。作为 $\mathbb F_2$ 上的多项式，
$v_e$ 是正文字的因子 $1+x_v$ 与负文字的因子 $x_v$
的乘积。重言子句有 $v_e=0$，空子句有 $v_e=1$。
删除重复文字后，$\deg v_e\le k$ 且 $v_e$
展开为至多 $2^k$ 个单项。置 $v(x)=(v_1(x),\ldots,v_m(x))$。
对有 $m$ 列的二值矩阵 $A$，第 $i$ 行指定方程
\begin{equation}\label{eq:hash-query}
 P_i(x):=\sum_{e:A_{ie}=1} v_e(x)=0\qquad\text{（在 }\mathbb F_2\text{ 上）}.
\end{equation}
记该方程组为 $\mathcal P_A$，其公共零点个数为
$N(A)=N(\mathcal P_A)$。每个可满足赋值通过每一行。
故 $N(A)=\#F+D(A)$，其中 $D(A)$ 统计仍然通过的
不可满足赋值。我们不知道 $\#F$，但它对每个候选新行都相同：
使 $N$ 最小的行也使 $D$ 最小。

我们用对 $(e,c)$ 的列表表示检验行，其中 $e\in[m]$、
$c\in\{0,1\}$；子句下标重复时系数按模二合并。
列表长度计入全部对，包括重复下标与系数为零的对。
列表长度为 $s$ 的行在相消之前至多读取 $ks$ 个变量。

\subsection{排除小违反集的初始检验}

若一个赋值只违反众多子句中的一个，则抽样一列短子句的
检验不太可能注意到它。因此我们首先构造矩阵 $A^{(0)}$，
它拒绝每个支撑小的非零违反向量。
通过这些初始检验的每个不可满足赋值必违反许多子句。
于是，抽样一列短子句的进一步检验能以常数概率拒绝每个
这样的赋值。

我们使用标准的基于边的 Tanner 码：把位指派到展开图的边
上，并要求与每个顶点关联的位构成某个指定局部线性码的字
（局部约束见 Tanner~\cite[Section~II]{Tanner81}，
展开图构造见 Sipser 与 Spielman~\cite[Sections~IV and~VI]{SS96}）。
通常的局部距离与展开论证给出非零码字权重的所需下界。
下面的证明确定性地构造图与局部码，其参数满足我们对检验
个数与长度的界。

\begin{lemma}\label{lem:sparse-checks}
设正整数 $n,m,B$ 满足 $m\le Bn$，并置
\[
 \mu=\lfloor\log_2 m\rfloor+1,\qquad
 P=20000B^2\mu^2,\qquad D_0=320000B^2\mu^3.
\]
存在 $A^{(0)}$ 的确定性构造：一个有 $m$ 列、至多 $3n$
行的二值矩阵，每行非零分量至多 $D_0$ 个，且
\[
 z\in\{0,1\}^m\setminus\{\vec{0}\},\quad
 A^{(0)}z=\vec{0}\pmod2
 \quad\Longrightarrow\quad
 |\operatorname{supp}z|\ge m/P.
\]
构造使用 $\exp(\poly(B,\log(m+2)))$ 次位运算与
$\poly(n,m,B)$ 比特工作空间。
\end{lemma}

\begin{proof}
我们构造一个边数在 $m$ 与 $2m$ 之间的二部图，
用其边作为临时二值坐标。每个顶点处的检验将迫使与某条
非零边关联的每个顶点有许多条非零关联边；
随后一个混合估计将迫使总权重大。
最后我们保留 $m$ 个边坐标，把其余固定为零。
令
\[
 d_G=\min(m,D_0),
 \qquad
 q_G=\min\{2^a:a\in\mathbb Z_{\ge0},\ 2^a d_G\ge m\}.
\]
则 $m\le q_Gd_G<2m$。写 $q_G=2^{\ell_G}$，
取 $\{0,1\}^{\ell_G}$ 的两份副本作顶点集。
对每个左顶点 $u$ 与每个端口 $j\in[d_G]$，
放一条从 $u$ 到 $u\oplus s_j$ 的边，
其中 $\oplus$ 表示按位模二加，
移位 $s_1,\ldots,s_{d_G}\in\{0,1\}^{\ell_G}$
将在下文选取。边以对 $(u,j)$ 为下标；
平行边保持互异。每侧每个顶点在每个端口 $j$
都有一条关联边，故图是 $d_G$ 正则的。

为得到混合估计 \eqref{eq:hash-mixing}，
我们选取移位使非空坐标集的每个奇偶性都近乎平衡。
对非零向量 $b\in\{0,1\}^{\ell_G}$，
位 $b\cdot s_j\bmod2$ 是 $s_j$ 中 $b$ 取一的坐标的
奇偶性。取 $\varepsilon=1/(200B\mu)$，要求
\begin{equation}\label{eq:hash-walsh}
 \Bigl|\sum_{j=1}^{d_G}(-1)^{b\cdot s_j}\Bigr|
 \le\varepsilon d_G
 \qquad
 \bigl(b\in\{0,1\}^{\ell_G}\setminus\{\vec{0}\}\bigr).
\end{equation}
\eqref{eq:hash-walsh} 中的和是奇偶为偶的下标个数
减去为奇的下标个数。于是对每个非零 $b$，
每种奇偶性出现在比例介于 $(1-\varepsilon)/2$
与 $(1+\varepsilon)/2$ 之间的下标处。
具有该性质的移位列表称为 $\varepsilon$
偏置的（$\varepsilon$-biased）~\cite[Definitions~2.1--2.2]{NN93}。

若 $q_G=1$，则没有条件要检验。否则 $d_G=D_0$。
为证明合适的移位存在，我们独立均匀随机地选取。
对每个固定的非零 $b$，符号 $(-1)^{b\cdot s_j}$
独立均匀分布于 $\{-1,1\}$。
Hoeffding 不等式把该 $b$ 违反
\eqref{eq:hash-walsh} 的概率界定为
$2\exp(-\varepsilon^2d_G/2)$。对 $b$
用联合界，总失败概率至多
\[
 2q_G\exp(-\varepsilon^2d_G/2)
 =2q_G\exp(-4\mu)<1,
\]
其中 $q_G<2m<2^{\mu+1}$。因此，穷举移位元组并检验
\eqref{eq:hash-walsh}，可在
$\exp(O(d_G\log(m+2)))$ 时间内找到有效元组。

接下来从 \eqref{eq:hash-walsh} 推导混合估计。
设 $\mathsf A_G$ 为 $q_G\times q_G$ 矩阵，
$\mathsf A_G[u,v]=|\{j\in[d_G]:u\oplus s_j=v\}|$；
其分量统计从左顶点 $u$ 到右顶点 $v$ 的边数。
两侧有同一下标集 $\{0,1\}^{\ell_G}$，
其上函数 $\chi_b(u)=(-1)^{b\cdot u}$
构成正交基。对每个 $b,u\in\{0,1\}^{\ell_G}$，
\[
 (\mathsf A_G\chi_b)(u)
 =\sum_{j=1}^{d_G}\chi_b(u\oplus s_j)
 =\chi_b(u)\sum_{j=1}^{d_G}(-1)^{b\cdot s_j}.
\]
于是每个 $\chi_b$ 都是特征向量：常值函数
$\chi_{\vec{0}}$ 有特征值 $d_G$，而
\eqref{eq:hash-walsh} 把所有其他特征值的绝对值界定在
$\varepsilon d_G$。对左顶点集 $U$ 与右顶点集 $V$，
令 $e_G(U,V)$ 统计二者之间的边数（含重数）。
把 $U,V$ 的指示向量分解为常值部分与正交部分，
应用 Cauchy--Schwarz 得
\begin{equation}\label{eq:hash-mixing}
 e_G(U,V)\le
 \frac{d_G|U||V|}{q_G}
 +\varepsilon d_G\sqrt{|U||V|}.
\end{equation}
为指定局部码，置
\[
 c_G=\left\lceil\frac{d_G}{8B}\right\rceil,
 \qquad \delta=\frac1{100B\mu}.
\]
我们需要一个 $c_G$ 行 $d_G$ 列的二值矩阵
$A_{\mathrm{loc}}$，使其核（模二）中每个非零字的权重
至少 $\delta d_G$。为证明这样的矩阵存在，
考虑分量独立均匀取位的矩阵。对每个固定的非零字，
其某个非零分量使每行中的奇偶性均匀；
由行的独立性，该字落在核中的概率为 $2^{-c_G}$。
置 $h=\lceil\delta d_G\rceil-1$。若 $h=0$，
则没有权重小于 $\delta d_G$ 的非零字。
否则其个数至多为
\[
 \sum_{j=1}^{h}\binom{d_G}{j}
 \le(h+1)(2d_G)^h
 \le(2d_G)^{2h}
 <2^{d_G/(25B)}.
\]
最后一个不等式用 $h<\delta d_G$、$d_G\le m<2^\mu$
与 $\mu+1\le2\mu$。由 $c_G\ge d_G/(8B)$，
联合界表明合适的 $A_{\mathrm{loc}}$ 存在。
枚举 $2^{c_Gd_G}$ 个二值矩阵并检验全部权重小于
$\delta d_G$ 的字，可在 $\exp(\poly(B,\log(m+2)))$
次位运算内找到一个。

在两侧每个顶点处，把 $A_{\mathrm{loc}}$ 的 $c_G$
行按端口次序施加到其关联边坐标上。
这些约束定义了一个局部码为
$C_{\mathrm{loc}}=\ker_{\mathbb F_2}A_{\mathrm{loc}}$
的 Tanner 码。其关联图对我们图的每条边有一个位节点，
对每个顶点有一个局部码节点；每个位参与其两个端点处的
约束。考虑被所有这些检验接受的非零边字。
设 $U,V$ 为其非零边在两侧触碰的顶点，$h'$
为其权重。$U\cup V$ 中每个顶点至少关联
$\delta d_G$ 条非零边，故
$h'\ge\delta d_G\max(|U|,|V|)$。因此
\[
 \delta d_G\sqrt{|U||V|}
 \le h'\le e_G(U,V)
 \le\frac{d_G|U||V|}{q_G}
       +\varepsilon d_G\sqrt{|U||V|}.
\]
$U$ 与 $V$ 都非空。除以 $d_G\sqrt{|U||V|}$
得 $\sqrt{|U||V|}\ge(\delta-\varepsilon)q_G$，
从而
\[
 h'\ge\delta(\delta-\varepsilon)q_Gd_G
 =\frac{q_Gd_G}{20000B^2\mu^2}\ge\frac mP.
\]

通过保留 $q_Gd_G$ 个边坐标中的任意 $m$ 个
（把它们等同于 $[m]$）并把其余坐标置零，缩短该
Tanner 码。把全部顶点检验限制到保留坐标上即得
$A^{(0)}$。被 $A^{(0)}$ 接受的字经补零延拓为被接受的
边字，故具有所需权重。每行非零分量至多 $d_G\le D_0$
个。行数为 $2q_Gc_G$。若 $q_G=1$，则 $d_G=m$ 且
\[
 2c_G\le\frac{m}{4B}+2\le\frac n4+2\le3n.
\]
若 $q_G>1$，则 $d_G=D_0\ge8B$，且
\[
 2q_Gc_G\le\frac{q_Gd_G}{4B}+2q_G
 \le\frac{q_Gd_G}{2B}<\frac mB\le n.
\]
对移位、局部矩阵与检验字的搜索依序进行。
存储当前选取与最终的行使用 $\poly(n,m,B)$
比特工作空间。
\end{proof}

\subsection{通过计数公共零点选取附加方程}
初始检验之后，仍然通过检验的每个不可满足赋值至少违反
$m/P$ 个子句。这使短随机检验能以常数概率抓住每个这样的
赋值。为使选取确定性，我们使用条件期望法背后的平均原理
~\cite{Raghavan88}。

\begin{lemma}\label{lem:greedy-hashes}
取引理~\ref{lem:sparse-checks} 的 $P,D_0$ 与
$A^{(0)}$，一个确定性过程追加 $3n$ 行列表长度 $P$
的行，得到满足 $N(A)=\#F$ 的矩阵 $A$。
它至多使用 $3n(2m)^P$ 次公共零点计数。
每个候选方程组至多 $6n$ 个方程，
每个方程至多涉及 $2^k\max(D_0,P)$ 个次数至多 $k$
的单项。
\end{lemma}
\begin{proof}
在当前矩阵为 $A$ 的阶段，枚举有序列表
\[
 \gamma=((e_1,c_1),\ldots,(e_P,c_P))
 \in\mathcal C:=([m]\times\{0,1\})^P.
\]
列表 $\gamma$ 表示二值行 $h_\gamma$：
\[
 (h_\gamma)_e=
 \Bigl(\sum_{j:e_j=e}c_j\Bigr)\bmod2
 \qquad(e\in[m]).
\]
列表共 $(2m)^P$ 个；表示同一行的不同列表仍分别计数。
把 $h_\gamma$ 追加到 $A$ 所得的矩阵记作
$A;h_\gamma$。

固定被 $A$ 接受的一个不可满足赋值 $x$。
初始行仍在 $A$ 中，故
$|\operatorname{supp}v(x)|\ge m/P$。
对均匀列表 $\gamma$，$e_1,\ldots,e_P$
都不属于该支撑的概率至多
\[
 (1-1/P)^P\le e^{-1}<1/2.
\]
在命中的条件下，选取一个命中位置 $j$。
固定除 $c_j$ 外的所有系数后，均匀位 $c_j$
使 $h_\gamma\cdot v(x)$ 模二均匀。
当某些子句下标重复时仍然如此。
于是对每个这样的 $x$，至少四分之一的候选列表给出
拒绝 $x$ 的行。

把该界对仍然通过的不可满足赋值求和，得
\[
 \mathbb E_\gamma D(A;h_\gamma)\le\tfrac34 D(A).
\]
对每个候选，通过计数其多项式方程组的公共零点计算
$N(A;h_\gamma)$。追加取得最小值的候选，
并列时按枚举次序取先者。
由每个 $\gamma$ 都有
$N(A;h_\gamma)=\#F+D(A;h_\gamma)$，
所选行留下至多 $\tfrac34 D(A)$ 个不可满足赋值。
$3n$ 轮之后至多剩
$(3/4)^{3n}2^n=(27/32)^n<1$ 个；
其个数是整数，故为零。
该过程运行固定轮数，从不需要知道 $D(A)$ 或 $\#F$。

初始矩阵至多 $3n$ 行，每个候选至多用 $3n$
个附加行求值。由 \eqref{eq:hash-query}，
列出 $s$ 个子句的行在原变量上展开为至多 $2^k s$
个单项；重复项在 $\mathbb F_2$ 上相消。
这些观察给出所述规模与查询界。
\end{proof}

\subsection{对单个稀疏公式计数}\label{subsec:count-sparse}
设 $n\ge20$、$1\le m\le Bn$，其中 $B$
被只依赖 $k$ 的 $\log n$ 固定幂界定。
用引理~\ref{lem:sparse-checks} 构造初始矩阵。
这使用 $\exp((\log n)^{O_k(1)})$ 次位运算与多项式空间。
$P$ 与 $D_0$ 都是 $(\log n)^{O_k(1)}$，
故引理~\ref{lem:greedy-hashes} 的每个候选都是
定理~\ref{thm:overview-sparse-polynomials}
在单个固定 $c=c(k)$ 下的实例。
这同样适用于被丢弃的候选与最终方程组：
计数器只需要方程的个数、次数与稀疏度。
对固定的 $k$，展开每行并合并单项花费多项式时间与空间。

至多 $3n(2m)^P+1=\exp((\log n)^{O_k(1)})$
次调用即够，含一次最终计数。每次调用花费
$2^{O_k(n(\log\log n)^2/\sqrt{\log n})}$ 次位运算。
由于 $\log n$ 的任何固定幂都是
$o(n(\log\log n)^2/\sqrt{\log n})$，
包括构造与候选比较在内的总时间为
\begin{equation}\label{eq:sparse-formula-cost}
 (n+m)^{O_k(1)}2^{O_k(n(\log\log n)^2/\sqrt{\log n})}.
\end{equation}
由贪心引理，最终计数是 $\#F$。
无子句或含空子句的公式直接处理。
候选被依序检查；除计数器的多项式工作空间外，
只需保留当前各行、一个候选、最佳行及其计数。

\subsection{回到一般公式}\label{subsec:general-reduction}

我们用不交稀疏化把可满足赋值分拆到同一 $n$
个声明变量上的稀疏公式中。前述算法对每个这样的公式计数。
参数 $\sigma$ 控制一个权衡：取小会减少公式数，
但增大其出现界 $B$。我们的求解器容许 $B$
按 $\log n$ 的固定幂增长，
因此可以让 $\sigma$ 趋于零，
同时把公式数与每公式成本都保持在所需时间界内。

\begin{lemma}\label{lem:sparsify}
固定 $k\ge2$。给定 $n$ 个声明变量、$m$
个子句的 $k$-CNF 与有理数 $0<\sigma\le1$，
一个确定性算法在同一变量上构造一棵 $k$-CNF 树。
设 $L_\sigma$ 为所给 $\sigma$ 的分子与分母的总二进制
编码长度。一个赋值满足输入当且仅当它满足某个叶公式，
且它不能同时满足两个不同的叶。
该树有
\[
 \operatorname{poly}_k(n,\sigma^{-1})\,2^{\sigma n}
\]
个节点，深度为 $\poly_k(n,\sigma^{-1})$。
在每个叶中，每个变量至多出现 $B$ 次，
其中 $B$ 是显式计算出的整数，满足
\[
 B=O_k\!\Bigl(\bigl[\sigma^{-1}\log(2/\sigma)\bigr]^{k-1}\Bigr).
\]
输入归一化使用 $\poly_k(n,m)$ 次位运算。
构造数值参数使用 $\poly_k(L_\sigma)$ 次位运算。
此后，处理每个节点使用 $\poly_k(n,\log(2/\sigma))$
次位运算。输入计数是叶计数之和。
\end{lemma}
附录~\ref{app:sparsify} 证明该引理，
包括为叶计数相加提供正当性的不交性。

\begin{proposition}\label{prop:interpolation-bound}
对每个固定的 $k\ge2$，插值算法对 $n\ge20$
个声明变量、$m$ 个子句的 $k$-CNF
计数可满足赋值，使用 $(n+m)^{O_k(1)}2^{O_k(\rho(n))}$
次位运算与 $(n+m)^{O_k(1)}$ 比特空间。
\end{proposition}
\begin{proof}
我们先以 $(n+m)^{O_k(1)}$ 时间归一化输入：
删除重复文字、重言子句、重复子句以及包含其他子句的
子句。我们保留全部 $n$ 个声明变量。
空子句给出计数零，无子句的公式给出计数 $2^n$。
$n<20$ 用穷举搜索处理。
对 $n\ge20$，接着选取稀疏化参数。定义
\[
 \lambda_n=\lceil\log_2 n\rceil,\qquad
 \nu_n=\lceil\log_2\lambda_n\rceil,\qquad
 s=\left\lfloor\frac{\sqrt{\lambda_n}}{\nu_n^2}\right\rfloor.
\]
我们用整数比较计算 $s$：
取满足 $s^2\nu_n^4\le\lambda_n$ 的最大非负整数。
若 $s<2$，用穷举搜索；
$\sqrt{\lambda_n}/\nu_n^2$ 趋于无穷，
故只有有限多个这样的规模。
否则以 $\sigma=1/s$ 应用引理~\ref{lem:sparsify}。
树至多 $2^{n/s}\operatorname{poly}_k(n,s)$
个节点，其叶的出现界为
\[
 B=O_k\bigl((s\log s)^{k-1}\bigr)=(\log n)^{O_k(1)}.
\]
每个叶使用原 $n$ 个变量，非空子句至多 $Bn$ 个。
我们用第~\ref{subsec:count-sparse} 节的稀疏公式算法
对每个叶计数，其运行时间由
\eqref{eq:sparse-formula-cost} 给出。
无子句或含空子句的叶直接处理。
该应用中的子句数是叶的子句数；
最终界中的 $m$ 仍是输入子句数。

对 $s\ge2$，其定义中的取整给出
\[
 \frac ns=O\!\left(\frac{n\nu_n^2}{\sqrt{\lambda_n}}\right)
 =O(\rho(n)).
\]
于是叶数与每叶计数时间的指数都是 $O_k(\rho(n))$。
其多项式因子，连同归一化与树节点上的工作，
被 $(n+m)^{O_k(1)}$ 界住。位运算总数因此为
\[
 (n+m)^{O_k(1)}2^{O_k(\rho(n))}.
\]
有限多个穷举搜索情形被吸收进常数中。
全部参数、初始检验与附加行都由上述确定性过程计算；
每次计数查询都使用
定理~\ref{thm:overview-sparse-polynomials}
的公共零点算法。

为界定完整算法使用的空间，我们深度优先遍历稀疏化树，
一次对一个叶计数。当前路径有 $\poly_k(n,s)$
个节点，每个归一化公式有 $\poly_k(n)$
比特编码。由 $s\le\sqrt{\lceil\log_2 n\rceil}$，
存储该路径并对一个叶计数使用 $\poly_k(n,m)$ 比特。
不交性使已处理叶计数之和保持在零与 $2^n$ 之间，
故其累加器至多用 $n+1$ 比特。
因此总空间为 $\poly_k(n,m)$ 比特。
\end{proof}

\appendix
\section{可满足赋值的分拆}\label{app:sparsify}

固定 $k\ge2$ 与有理数 $0<\sigma\le1$。
我们通过把公式反复分裂成两种情形来证明
引理~\ref{lem:sparsify}。两个孩子使用同样的 $n$
个变量，且父公式的每个可满足赋值恰满足一个孩子。
我们沿用选取共享文字集的规则与~\cite[Section~6]{CIP}
的插入记账。使两种情形互不相同的精细化在精确计数中是
标准的~\cite[Appendix~A]{DHMTW}。

在本附录中，子句是文字的集合，表示其析取。
其\emph{宽度}是其所含文字数。
我们删除重言子句、重复子句与包含其他子句的子句。
于是剩余子句在包含关系下构成一个反链。
我们保留全部 $n$ 个声明变量，且不传播单元子句。
空子句给出计数零，空公式给出计数 $2^n$。

\subsection{分支规则}

固定整数 $M\ge2$。我们将在每个叶得到出现界
$B=O_k(M^{k-1})$，然后选取 $M$ 使树至多
$\poly_k(n,\sigma^{-1})2^{\sigma n}$ 个节点。
置
\[
 \theta_j=M^j\qquad(0\le j<k).
\]
记 $F$ 为当前子句族，对每个文字集 $H$ 定义
\[
 \deg_c(H)=\#\{C\in F:|C|=c,\ H\subseteq C\}.
\]
只要宽度 $c$ 的子句的非空真子集 $H$ 满足
\begin{equation}\label{eq:spar-violation}
 \deg_c(H)\ge\theta_{c-|H|},
\end{equation}
我们就分支。算法选取最小的可能 $c$，再选取最大的可能
$|H|$；其余并列确定性打破。
于是，一个留下 $j=c-|H|$ 个其他文字的非空共享文字集
$H$，必须在至少 $M^j$ 个子句中出现我们才在其上分支。

设 $H\cup P_1,\ldots,H\cup P_z$ 为包含所选 $H$
的全部宽度 $c$ 子句，$G$ 为其余子句。置 $j=c-|H|$。
这里 $P_a$ 是从第 $a$ 个被选子句中移除 $H$
后剩下的文字集，故每个 $P_a$ 宽度为 $j$ 且
$z\ge\theta_j$。两个孩子是
\begin{equation}\label{eq:spar-children}
 F_A=G\wedge H,\qquad
 F_B=G\wedge\bigwedge_{a=1}^zP_a
          \wedge\bigwedge_{\ell\in H}\{\neg\ell\}.
\end{equation}
我们在每个孩子中删除被其新子句弄得冗余的子句，
以恢复反链条件。集合 $P_a$ 可以互相重叠。
两个孩子的可满足赋值互不相交：
$F_A$ 要求 $H$ 的某个文字为真，
而 $F_B$ 要求 $H$ 的每个文字为假。

\begin{lemma}\label{lem:spar-branch}
\eqref{eq:spar-children} 中的两个孩子划分其父公式的
可满足赋值。每个新子句的宽度小于 $c$。
$A$ 孩插入一个新子句，$B$ 孩插入至少 $\theta_j$
个宽度 $j$ 的新子句。所有新插入的子句都在即时归一化中
存活。任一孩子中被删除的每个旧子句都包含一个新插入的
真子集。
\end{lemma}
\begin{proof}
每个赋值恰满足条件「$H$ 的某个文字为真」与
「$H$ 的每个文字为假」之一。
在前一情形，全部被选子句被满足，剩下 $G\wedge H$。
在后一情形，每个被选子句化归为 $P_a$，
而单元子句 $\{\neg\ell\}$ 强制该情形。
因此两个孩子划分父公式的可满足赋值。

旧反链不含被 $H$ 或任何 $P_a$ 包含的子句：
那样的子句将是 $H\cup P_a$ 的真子集。
于是 $H$ 与各互异的等宽子句 $P_a$ 都是新的。
任何 $P_a$ 不含 $H$ 的变量，因为其父子句非重言。
因此，单元子句 $\{\neg\ell\}$（$\ell\in H$）
都不是任何 $P_a$ 的子集。
已经存在的单元子句 $\{\neg\ell\}$ 被保留，
不计为新插入。其余新单元子句与子句 $P_a$ 全部存活。
每个被选旧子句包含孩子 $A$ 的新插入子句 $H$，
或孩子 $B$ 的新插入子集 $P_a$。
按归一化规则，每个其他被删子句都包含一个新插入的
真子集。
\end{proof}

\subsection{统计插入}
为界定树的大小，我们统计沿一条路径引入的每个子句
（包括后来被删除的子句）。输入中已有的子句不计入。
由引理~\ref{lem:spar-branch}，每步插入一个子句，
因此界定插入总数也就界定路径长度。

在此记账中，子句在分支引入它时被标记。
我们先界定仍存在且含给定文字集的被标记子句数。
最小宽度优先级控制分支之前存在的子句；
最大 $|H|$ 优先级控制该分支引入的子句。

\begin{lemma}\label{lem:spar-marked-degree}
在每个节点，对 $1\le i<j<k$ 与每个大小 $i$
的文字集 $S$，仍存在的、含 $S$ 的宽度 $j$
被标记子句少于 $2\theta_{j-i}$ 个。
\end{lemma}
\begin{proof}
初始没有被标记子句。只需考虑插入宽度 $j$
子句的一个分支。其所选宽度 $c$ 大于 $j$。
插入之前，含 $S$ 的当前宽度 $j$ 子句少于
$\theta_{j-i}$ 个：否则最小宽度优先级会选取至多 $j$
的宽度。

选孩子 $A$ 只插入一个子句。
设孩子 $B$ 插入宽度 $j$ 的子句 $P_a$。
若没有 $P_a$ 含 $S$，则含 $S$ 的新宽度 $j$
子句数为零。否则 $S$ 落在某个花瓣内，
因此与 $H$ 不交。此时，$P_a$ 含 $S$
当且仅当被选子句 $H\cup P_a$ 含 $H\cup S$。
由 $0<i<j$，$H\cup S$ 大于 $H$，
且是某个被选宽度 $c$ 子句的真子集。
最大共享集的优先级因此给出
\[
 \#\{a:S\subseteq P_a\}
 =\deg_c(H\cup S)<\theta_{c-|H\cup S|}
 =\theta_{j-i}.
\]
其他新子句 $\{\neg\ell\}$ 的宽度为一。
于是宽度 $j$ 的新旧子句合起来满足所述界。
删除不能增大它。
\end{proof}

\begin{lemma}\label{lem:spar-insertion-bound}
沿任何路径或路径前缀，宽度 $j$ 插入的总数 $T_j$
满足 $T_j\le4(3M)^{j-1}n$（$1\le j<k$）。
\end{lemma}
\begin{proof}
$n=0$ 的情形是显然的，设 $n\ge1$。
对当前路径前缀，设 $L_j$ 为仍存在的被标记宽度 $j$
子句数，$D_j$ 为其子句已被删除的宽度 $j$
插入数。则 $T_j=L_j+D_j$。
单元子句从不被删除，故 $T_1\le2n\le4n$。
对 $j\ge2$，把引理~\ref{lem:spar-marked-degree}
对 $2n$ 个单文字求和，得 $L_j\le4n\theta_{j-1}$。

每当一个宽度 $j$ 的被标记子句被删除，
就把该删除指派给它所含的一个新引入的真子集；
引理~\ref{lem:spar-branch} 保证这样的子集存在。
若新子句的宽度 $i<j$，引理~\ref{lem:spar-marked-degree}
表明在分支之前它被含于少于 $2\theta_{j-i}$
个宽度 $j$ 的被标记子句。于是宽度 $i$
子句的一次插入至多收到这么多次删除。
对新子句的各宽度 $i$ 求和得
$D_j\le2\sum_{i=1}^{j-1}\theta_{j-i}T_i$。因此
\begin{equation}\label{eq:spar-insertions}
 T_j=L_j+D_j\le4n\theta_{j-1}
       +2\sum_{i=1}^{j-1}\theta_{j-i}T_i.
\end{equation}
第一项统计仍存在的子句；和式统计已被删除的。
除以 $nM^{j-1}$ 并归纳，得
\[
 \frac{T_j}{nM^{j-1}}
 \le4+2\sum_{i=1}^{j-1}4\cdot3^{i-1}
 =4\cdot3^{j-1},
\]
如所需。
\end{proof}

\subsection{树的大小}
现在把引理~\ref{lem:spar-insertion-bound}
转化为树节点数的界。它同时控制路径长度与路径在
\eqref{eq:spar-children} 中选 $B$ 孩的次数。
把这两个界分别写作 $Cn$ 与 $\Gamma n/M$，其中
\[
 C=4\sum_{j=0}^{k-2}(3M)^j,\qquad
 \Gamma=4\sum_{j=0}^{k-2}3^j=2(3^{k-1}-1).
\]
引理~\ref{lem:spar-insertion-bound} 给出至多 $Cn$
的深度。一次选 $|P_a|=j$ 的 $B$ 孩插入至少 $M^j$
个宽度 $j$ 的子句，故一条路径至多含
$4\cdot3^{j-1}n/M$ 次这样的选取。
总计，一条路径至多 $\Gamma n/M$ 次选 $B$。

少选 $B$ 的路径有短的描述。
由于分支规则是确定性的，节点由其深度与路径选 $B$
的位置确定；其余位置都选 $A$。
对 $M\ge\Gamma$，这把节点数界定为
\[
 (Cn+1)\sum_{a\le\lfloor\Gamma n/M\rfloor}\binom{Cn}{a}
 \le (Cn+1)^2
       2^{(\Gamma n/M)\log_2(4CM)}.
\]
这里 $\Gamma n/M\le n\le Cn/4$。
对正整数 $a$，用 $\binom{Cn}{a}\le(4Cn/a)^a$；
该上界在 $a\le Cn/4$ 上递增。
被加项至多 $Cn+1$ 个。$\Gamma n/M<1$
的情形是显然的。

我们选取 $M$ 使该界中的指数至多 $\sigma n$。
按递增次序检验二的幂 $M\ge\Gamma$。
对每个选取，计算 $C$ 与 $h=\lceil\log_2(4CM)\rceil$，
当
\[
 \Gamma h\le\sigma M
\]
时停止。所有比较都是精确的。
由 $C=O_k(M^{k-2})$ 且 $\Gamma$ 只依赖 $k$，
第一个成功的选取满足
\[
 M=O_k\bigl(\sigma^{-1}\log(2/\sigma)\bigr).
\]
树因此至多 $(Cn+1)^2\,2^{\sigma n}$ 个节点。

在叶处，\eqref{eq:spar-violation}
对子句的每个非空真子集都不成立。
当 $c\ge2$ 时，固定文字属于少于 $\theta_{c-1}$
个宽度 $c$ 的子句，且至多属于一个单元子句。
故每个变量至多出现
\[
 B=2\sum_{j=0}^{k-1}M^j
   \le4M^{k-1}
   =O_k\Bigl(\bigl[\sigma^{-1}\log(2/\sigma)\bigr]^{k-1}\Bigr).
\]
于是每个叶都有每变量出现次数的所需界。
反复的两情形划分也表明每个原可满足赋值恰出现在一个叶，
故叶计数之和给出原计数。

为在节点处实现分支规则，我们对其当前公式计算
$\deg_c(H)$。对每个宽度 $c$ 子句，枚举其非空真子集
$H$，并把以 $(c,H)$ 为下标的计数器加一。
非空归一化子句至多 $O_k(n^k)$ 个可能，
每个至多 $2^k$ 个子集。子句计数器使用
$O_k(\log(n+2))$ 比特，阈值使用 $O_k(\log M)$
比特。因此，计算计数、选分支并归一化其孩子在每个节点
使用 $\poly_k(n,\log M)$ 次位运算与比特工作空间。
$M$ 的界给出 $\log M=O_k(\log(2/\sigma))$。

取引理~\ref{lem:sparsify} 中的 $L_\sigma$，
参数搜索检验 $O_k(L_\sigma)$ 个二的幂。
其整数操作数有 $O_k(L_\sigma)$ 位，
故构造数值参数使用 $\poly_k(L_\sigma)$
次位运算与比特工作空间。

\section{精确算术与空间}\label{app:arithmetic}
第~\ref{sec:construction} 节界住了代入行
$\vec{P}_{f,v}(T)$ 的有理系数。我们通过界定提升因子、
用于计算 $H$ 的 Taylor 多项式的和、以及最终插值，
来完成命题~\ref{prop:count} 的位成本与空间估计。
要点是每个和使用固定的分母。其分子随被加项个数增长，
故其位长只随该数的对数增长。
对有限多个未定元上的任何多项式
$G=\sum_\alpha c_\alpha \vec{X}^{\alpha}$，
写 $\|G\|_1=\sum_\alpha|c_\alpha|$
表示其系数范数。

\subsection{提升指示子的整系数}
\label{subsec:rational-lift-height}
我们先界定引理~\ref{lem:rational-idempotent}
给出的奇偶多项式的系数。它们出现在
\eqref{eq:lifted-factor} 的每个方程因子 $\mathcal B_i$
中，故其大小进入求值这些因子的位成本界。
对整数 $\ell\ge1$，回忆 $h(Y)=Y(Y-1)$ 与
引理~\ref{lem:rational-idempotent} 的多项式：
\[
 p_\ell(Y)=(1-Y)^\ell\sum_{j=0}^{\ell-1}
                    \binom{\ell+j-1}{j}Y^j
          =\sum_{j=0}^{\ell-1}e_j(Y)h(Y)^j,
 \qquad \deg e_j\le1.
\]
该引理证明奇偶同余，并用对 $h$ 的逐次首一除法构造
$e_j$。这里我们界定此构造中出现的整数。

\begin{lemma}\label{lem:rational-lift-height}
写 $e_j(Y)=a_j+b_jY$（$0\le j<\ell$）。
则 $a_j,b_j\in\mathbb Z$ 且
\[
 |a_j|,|b_j|\le2^{C\ell^2}\qquad(0\le j<\ell)
\]
（$C$ 为绝对常数）。用对 $h$ 的逐次除法构造这些多项式
使用 $\poly(\ell)$ 次位运算与比特工作空间。
\end{lemma}
\begin{proof}
系数范数不等式 $\|fg\|_1\le\|f\|_1\|g\|_1$
连同二项恒等式
\[
 \sum_{j=0}^{\ell-1}\binom{\ell+j-1}{j}
 =\binom{2\ell-1}{\ell-1}
\]
给出
\[
 \|p_\ell\|_1\le2^\ell\binom{2\ell-1}{\ell-1}\le2^{3\ell}.
\]
对 $h=Y^2-Y$ 的首一除法保持整系数。
每次首项消元使当前商与余式的系数范数之和至多增大三倍。
全部逐次除法共 $O(\ell^2)$ 次消元。
因此每个中间系数的绝对值至多 $2^{O(\ell^2)}$，
这证明对 $a_j,b_j$ 的所述界。其带符号二进制编码
从而使用 $O(\ell^2)$ 位。每个系数向量有 $O(\ell)$
个分量。逐次除法对 $O(\ell^2)$ 位的整数执行
$O(\ell^2)$ 次整数算术运算，
给出所述 $\poly(\ell)$ 位运算与空间界。
\end{proof}

\subsection{布尔赋值上求和的位复杂度}
\label{subsec:rational-streaming-height}
现在界定 \eqref{eq:full-taylor-sum} 中计算
$\mathcal H_t(z)=H(t+z)\bmod z^J$ 的每个中间值。
从代入行 $\vec{P}_{f,v}(T)$ 出发，
我们界定被代入的因子、它们对每组布尔变量的和、
以及命题~\ref{prop:count} 中 Taylor 数据的累加器
$T_t$。
固定 $n$ 个布尔变量上的多项式方程组 $\mathcal P$，
其 $m_{\mathrm{test}}$ 个方程次数至多 $k$，
精度 $\ell$ 满足 $2^\ell>2^n$，
并给定其方程因子的替代 design。
设行 $\vec{P}_{f,v}$、划分 $S_{t,g}$ 与矩阵 $C_{t,g}$
满足定义~\ref{def:layout}，参数为 $Q,R,r,D,\tau$。
回忆 $H$ 的次数界与每个节点保留的 Taylor
系数个数：
\[
 d_H=kR\bigl(m_{\mathrm{test}}+2(\ell-1)\bigr),\qquad
 J=\left\lfloor\frac{d_H}{Q}\right\rfloor+1.
\]
Taylor 变量 $z$ 与共享截断标记 $V$
的所有多项式都按其模 $(z^J,V^\ell)$
的唯一代表存储。移除这些次数范围之外的项
不会增大系数范数。

代入行有公分母
\[
 d_{\mathrm{den}}=(R-1)!Q!.
\]
我们使用整数 $B_P\ge1$，使每行
$d_{\mathrm{den}}\vec{P}_{f,v}(T)$
的每个整系数的绝对值至多 $2^{B_P}$。
定义~\ref{def:layout} 给出
$B_P=O(Q\log Q+\log D)$。
在固定节点 $t\in[Q]$ 与固定修正集 $U\subseteq[n]$
处，因子 $f$ 在变量 $v$ 的输入是
\[
 \theta_{f,v}(z;x)=\langle \vec{P}_{f,v}(t), x\rangle+
       \bigl\langle \vec{P}_{f,v}(t+z)-\vec{P}_{f,v}(t),\mathbf1_U\bigr\rangle,
 \qquad x\in\{0,1\}^n.
\]
设 $L_i(z;x)$ 为方程 $i$ 的列出单项在这些代入后的和。
每个 $d_{\mathrm{den}}\theta_{i,v}(z;x)$
是系数范数至多
\[
 B_0=d_{\mathrm{den}}+3n(R+1)2^{B_P}(Q+1)^R
\]
的整系数多项式。事实上，$\theta_{i,v}$
的三次行求值在清除分母后各贡献至多
$n(R+1)2^{B_P}(Q+1)^R$。
该界对组和中的代表同样成立。

置 $\Delta=d_{\mathrm{den}}^k$。
对次数 $u\le k$ 的单项，把其代入值乘以 $\Delta$
给出范数至多 $d_{\mathrm{den}}^{k-u}B_0^u\le B_0^k$
的整系数多项式（含常值单项）。
由 $S$ 界住列表总长，
\[
 B=\max\{\Delta,SB_0^k\}
 \quad\text{满足}\quad \|\Delta L_i(z;x)\|_1\le B .
\]
于是 $\log d_{\mathrm{den}}=O(Q\log(Q+2))$、
$\log B_0=O(Q\log(Q+2)+\log(nD+2))$，
且 $\log B=O(k\log B_0+\log(S+2))$。

方程 $i$ 的被代入因子是
\[
 \mathcal B_i(z,V;x)=\sum_{j=0}^{\ell-1}
       e_j\bigl(L_i(z;x)\bigr)h\bigl(L_i(z;x)\bigr)^jV^j
       \pmod{(z^J,V^\ell)}.
\]
由 $e_j$ 是整数的且次数至多一，其代入的分母整除
$\Delta$；$h(L_i)^j$ 的分母整除 $\Delta^{2j}$。
于是，令
\[
 \Gamma=\Delta^{2\ell-1},
\]
$\Gamma\mathcal B_i(z,V;x)$ 的每个系数都是整数。
为界定所得整数分子，置
\[
 E_\ell=\max\{1,\|e_0\|_1,\ldots,\|e_{\ell-1}\|_1\},\qquad
 B_1=\ell E_\ell(B+\Delta+2)^{2\ell-1}.
\]
对 $A_i=\Delta L_i$，两个被代入因子有整数分子
$\Delta e_j(L_i)$ 与 $A_i(A_i-\Delta)$，
范数分别至多 $E_\ell(B+\Delta)$ 与 $B(B+\Delta)$。
把每项垫到分母 $\Gamma$ 后得
\[
 \|\Gamma\mathcal B_i\|_1\le B_1,\qquad B_1\ge\Gamma,
 \qquad
 \log B_1=O\!\left(\ell^2+
       \ell\log(B+\Delta+2)+\log(\ell+1)\right).
\]

接下来把这些估计应用于实际的组和。
写 $S_g=S_{t,g}$、$C_g=C_{t,g}$，
设 $c_g$ 为 $C_g$ 的行数。于是
$c=\sum_g c_g\le2\tau n$，
且布尔测量 $C_gx_g$ 的每个分量落在
$[-D,D]\cap\mathbb Z$ 中。
设 $m_{\mathrm{test},g}$ 为在 $t$ 被指派到组 $g$
的方程因子数，故
$\sum_g m_{\mathrm{test},g}=m_{\mathrm{test}}$。
其乘积 $A_{g,U}(x;z)$
（\eqref{eq:group-polynomial}）
表示在 $\mathbb Q[z,V]/(z^J,V^\ell)$ 中。
其分母整除 $\Gamma^{m_{\mathrm{test},g}}$，
在该分母上的分子范数至多
$B_1^{m_{\mathrm{test},g}}$。
中性常数一因子不影响该界。

对满足 $|E|<J$ 的对 $U\subseteq E\subseteq[n]$
与可达测量元组 $\vec y=(y_1,\ldots,y_r)$，
选取满足 $C_h\bar x_h=y_h$ 的布尔代表
$\bar x_h$。\eqref{eq:group-sum} 中的组和为
\[
 \Sigma_g(\vec y)=
 \sum_{\substack{x_g\in\{0,1\}^{S_g}\\ C_gx_g=y_g\\
                  (x_g)_v=1\ (v\in E\cap S_g)}}
 A_{g,U}\!\Bigl(x_g+\sum_{h\ne g}\bar x_h;z\Bigr).
\]
这里各组上的向量补零延拓。
该和至多 $2^{|S_g|}$ 项，全在同一分母上，故
\[
 \left\|\Gamma^{m_{\mathrm{test},g}}\Sigma_g(\vec y)\right\|_1
 \le 2^{|S_g|}B_1^{m_{\mathrm{test},g}}.
\]
对 $1\le g\le r$，定义部分乘积
$\Pi_g(\vec y)=\prod_{h=1}^g\Sigma_h(\vec y)$
（每次乘法后截断）。
其分母整除
\[
 \Gamma^{\sum_{h\le g}m_{\mathrm{test},h}}.
\]
因此每个完整贡献 $\Pi_r(\vec y)$
可放到固定分母
\[
 \Omega=\Gamma^{m_{\mathrm{test}}}
\]
上，且满足
\[
 \|\Omega\Pi_r(\vec y)\|_1
 \le2^nB_1^{m_{\mathrm{test}}}.
\]

设 $T_t(z,V)$ 为把全部对与测量元组上的这些完整贡献
带符号 $(-1)^{|E|-|U|}$ 相加的累加器。
由 \eqref{eq:full-taylor-sum}，其最终值是取值于
$\mathcal A_\ell$ 的 Taylor 多项式
$\mathcal H_t(z)$。
一个节点处的带符号更新数 $M_*$ 界定为
\[
 M_*\le
 \biggl(\sum_{s=0}^{\min(J-1,n)}2^s\binom ns\biggr)(2D+1)^c
 \le3^n(2D+1)^c.
\]
因此 $T_t$ 的每个中间值（不只是最终值）满足
\[
 \|\Omega T_t\|_1
 \le M_*\,2^nB_1^{m_{\mathrm{test}}}.
\]
整系数多项式 $\Omega T_t$ 的每个系数有
\[
 O\!\left(\log(M_*+1)+n+m_{\mathrm{test}}\log B_1\right)
\]
位。该估计同样以上述固定分母界定全部组和与部分乘积的
整数分子。分母满足
\[
 \log\Omega
 =k m_{\mathrm{test}}(2\ell-1)\log d_{\mathrm{den}}
 =O\!\left(k(m_{\mathrm{test}}\ell+1)Q\log(Q+2)\right).
\]

每个贡献至多相乘 $m_{\mathrm{test}}$
个提升方程因子。继续处理对或测量元组只是添加新贡献，
并不把它们乘入之前的贡献。
这正是指数多个迭代只给位估计贡献 $\log(M_*+1)$
的原因。每个存储的多项式至多
$J\ell\le(d_H+1)\ell$ 个系数。
枚举保留其当前下标与布尔代表、一个组和、一个部分乘积、
以及累加器 $T_t$，而不存储全体赋值或全体测量元组的表。
这些估计给出 \eqref{eq:sum-bit-bound}
与每次有理运算的多项式位成本。
$QJ\ell$ 个 Taylor 系数与工作状态容于
$\poly(n,m_{\mathrm{test}},\ell,D,Q,R,J,r,k,S)$
比特。

\subsection{插值与整数答案}
\label{subsec:rational-readout}
对 $0\le s<\ell$，设 $H_s(T)=[V^s]H(T)\in\mathbb Q[T]$。
第~\ref{subsec:rational-streaming-height} 节界定
对 $t\in[Q]$、$0\le j<J$ 计算出的数据
$[z^jV^s]\mathcal H_t(z)=[z^j]H_s(t+z)$。
固定 $s$。我们对这一坐标插值，
对全部 $\ell$ 个坐标使用同一有理 Hermite 矩阵。
置 $M_H=QJ>d_H$。在单项基
$1,T,\ldots,T^{M_H-1}$ 下，Hermite 插值矩阵是整数矩阵
\[
 \mathsf V_{(t,j),u}=
 \begin{cases}\binom uj t^{u-j},&u\ge j,\\0,&u<j,\end{cases}
 \qquad t\in[Q],\quad 0\le j<J,\quad 0\le u<M_H.
\]
由第~\ref{sec:counting} 节的 Hermite
插值命题，它非奇异。其分量有 $O(M_H\log(Q+2))$
位。设 $\mathbf y$ 为按 $(t,j)$ 排序的计算出的系数
$[z^jV^s]\mathcal H_t(z)$ 的向量，
$\mathbf h$ 为 $H_s$ 的单项系数向量（补零到长
$M_H$）。则 $\mathsf V\mathbf h=\mathbf y$。
向量 $\mathbf b=\Omega\mathbf y$ 是整数的；
记 $B_{\mathrm{data}}$ 为其分量位长的界。
第~\ref{subsec:rational-streaming-height} 节的界给出
\[
 B_{\mathrm{data}}
 \le\poly(n,m_{\mathrm{test}},\ell,k,Q,\log(D+2),\log(S+2)).
\]

只需常系数 $h_0=H_s(0)$。设 $\mathsf V^{(0)}$
由 $\mathsf V$ 把其 $u=0$ 列换成 $\mathbf b$
而得，并置
\[
 N_H=\det\mathsf V^{(0)},\qquad D_H=\det\mathsf V\ne0.
\]
对 $\mathsf V(\Omega\mathbf h)=\mathbf b$
用 Cramer 法则给出
\begin{equation}\label{eq:exact-readout}
 H_s(0)=h_0=\frac{N_H}{\Omega D_H}.
\end{equation}
我们用免分数高斯消元
~\cite[Sections~I--II and~IV--V]{Bareiss68}
计算这两个行列式。
对 $M_H\times M_H$ 整数矩阵，该算法使用
$O(M_H^3)$ 次整数算术运算；主元选取使其也能处理奇异
矩阵。保留的消元分量是输入矩阵的带符号子式。
每次更新对这些分量作乘积与差，随后作精确整数除法。
因此，$\mathsf V$ 与 $\mathsf V^{(0)}$
的子式的界也界定用于计算 $D_H$ 与 $N_H$
的中间整数。
$(\mathsf V\mid\mathbf b)$ 的每个子式的绝对值至多
\[
 2^{O\left(M_H\left(B_{\mathrm{data}}+
                M_H\log(Q+2)+\log(M_H+1)\right)\right)}.
\]
事实上，其行列式是至多 $M_H!$ 个、每个至多 $M_H$
个分量的乘积之和。
因此免分数消元中的中间整数在带符号二进制记法下有
\[
 O\!\left(M_H\bigl(B_{\mathrm{data}}+
                   M_H\log(Q+2)+\log(M_H+1)\bigr)\right)
\]
位。计算 $N_H,D_H$
于是使用 $\poly(M_H,B_{\mathrm{data}})$
次位运算与比特工作空间。

在 $T=0$ 处每个 $H_s(0)$ 由构造是整数，
故 \eqref{eq:exact-readout} 中的除法对每个 $s$
都是精确的。算出这 $\ell$ 个整数后，
置 $u=\sum_{s=0}^{\ell-1}H_s(0)$。
由 \eqref{eq:origin-residue}，最小非负剩余
\[
 u-2^\ell\left\lfloor\frac{u}{2^\ell}\right\rfloor=N(\mathcal P)
\]
就是精确计数。对相继的坐标复用同一工作区，
全部 $\ell$ 次插值、求和与取余使用
$\poly(M_H,B_{\mathrm{data}},\log(\Omega+1),\ell)$
次位运算与多项式比特工作空间。

\subsection{技术概览的算术}
\label{subsec:overview-arithmetic}
我们还为第~\ref{sec:technical-overview} 节的
受限出现次数算法论证位复杂度与空间界。
引理~\ref{lem:row-coefficients}
把替代系数的位长界定为 $n$ 的多项式。
求解第~\ref{subsec:preview-separation} 节的
多项式规模有理线性方程组，给出位长同样为多项式的
测量系数 $\vec\lambda_{e,v,h}$。
设 $\delta$ 为全部这些系数与全部替代系数
（跨所有求值点）的分母之积。
它们只有多项式个，故 $\delta$
的位长也是多项式的。每个子句乘积至多用 $M\le3n$
个线性输入；因此受限和中全部系数可用公分母
$\delta^{3n}$ 累加。位长为多项式的输入系数的
多项式个乘积，以及指数多个这样的项在该公分母上的和，
仍有多项式位长。因此上述精确 Hermite 插值对这些
Taylor 数据也有多项式位成本。最后，枚举可依序执行，
只保留构造、当前下标、系数阵列、运行中的和、
以及插值所需的多项式个 Taylor 系数。
这给出多项式空间与仅多项式的位运算开销。

\raggedbottom
\section{证伪 ETH 的其他推论}\label{app:consequences}
定理~\ref{thm:main} 给出一个电路下界，以及求逆常数输出
局部性布尔映射的亚指数时间算法。固定密度情形给出
推论~\ref{cor:quantitative-circuits} 的定量电路下界。
两个电路论证都只用我们 $3$-SAT 算法的判定版本；
精确原像计数使用其完整的计数保证。

\subsection{$\mathrm E^{\mathrm{NP}}$ 的电路下界}\label{subsec:circuit-lower-bounds}
对下文的定性界，记 $\mathrm{SIZE}(O(n))$
为拥有 $O(n)$ 个有界扇入 AND、OR、NOT 门的非一致电路的
语言类。

\begin{corollary}\label{cor:consequence-circuits}
我们有
\[
 \mathrm E^{\mathrm{NP}}\not\subseteq\mathrm{SIZE}(O(n)).
\]
\end{corollary}
\begin{proof}
Williams 证明：若 ETH 为假，则
$\mathrm E^{\mathrm{NP}}$ 没有线性规模电路
~\cite[Theorem~6.1]{Wil13}。把该蕴含应用于
定理~\ref{thm:main} 提供的确定性亚指数时间 $3$-SAT
算法。
\end{proof}

这一联系是从更快的可满足性算法导出电路下界的一个实例。
为编码有 $s$ 个门、$n$ 个输入的布尔电路，
为每个门引入一个变量，用其输入强制其取值，
并添加一个要求输出位等于 $1$ 的单元子句。
所得 $3$-CNF 有 $O(n+s)$ 个变量与子句，
且恰在电路输出为一时可满足。
因此我们的定理在 $2^{o(n)}$ 时间内判定线性规模电路的
可满足性。Williams 的论证把这样的加速与对见证的假想小
电路结合，与某个非确定性时间层级矛盾。
该推论给出 $\mathrm E^{\mathrm{NP}}$
中没有线性规模电路族的语言；
它不断言 NP 的超多项式电路下界。

\begin{proof}[推论~\ref{cor:quantitative-circuits} 的证明]
我们使用 Williams 的见证电路论证
~\cite[Theorem~6.1]{Wil13}，连同 Jahanjou、Miles 与
Viola 的局部归约~\cite[Theorem~1 and Section~6]{JMV13}。

由 \v{Z}\'ak 的非确定性时间层级~\cite{Zak83}
（取~\cite[Theorem~2.1 in the full version]{Wil13}
的表述形式），固定多带图灵机上的一元语言
$H\in\operatorname{NTIME}(2^n)\setminus\operatorname{NTIME}(2^{n/2})$。
对一切足够大的 $n$，把局部归约以其固定机器作用于 $1^n$，
局部性参数取 $r=2$。时间限额 $T=2^n n^{O(1)}$
涵盖在归约中模拟机器的多对数开销，
给出 $q=n+O(\log n)$ 个子句下标位。
该归约在多项式时间内构造一个解码器，
在每个下标上给出三个变量标号及其符号。
令 $\varphi_n$ 由在全部 $2^q$ 个下标上解码出的子句组成。
它可满足当且仅当 $1^n\in H$。

解码器的 $O(n)$ 个输出位各有一棵常数深度的决策树，
故每个只依赖至多固定数 $w$ 个下标位。
在汉明权重至多 $w$ 的全部串上求值解码器
即可恢复每个支撑。事实上，见证对某一坐标依赖的两个串
可以在支撑之外置零而保持其相异的输出。
同样的求值给出每个支撑上的真值表。
这在多项式时间内构造出 $O(n)$ 个扇入二的门的解码器。

把每个变量标号垫到 $dn$ 位（$d$ 为固定正整数），
并声明全部 $2^{dn}$ 个变量。若对此类 $n$ 有 $|i|=dn$
且 $\varphi_n$ 可满足，令 $W(i)$ 为其按字典序第一个
可满足赋值的第 $i$ 位（变量按其标号排序）。
其他情形一律置 $W(i)=0$。$i$ 的长度确定 $n$。
为判定 $W$，写出 $2^q$ 个子句，问 NP 谕示公式是否可满足；
若可满足，通过询问「当前前缀接零是否仍有可满足扩张」
恢复第一个赋值。至多 $2^{dn}$ 个长度 $2^{O(n)}$
的前缀查询即够，故 $W\in\mathrm E^{\mathrm{NP}}$。

置 $a(t)=\sqrt{\log_2 t}$。检验
定理~\ref{thm:main} 固定密度部分的计数是否为正，
可在时间 $V^b2^{A V/a(V)}$（固定 $A,b>0$）
内判定有 $V$ 个变量、至多 $5V$ 个子句的 $3$-SAT。
对数底的变化被吸收进 $A$。求解器使用多项式空间，
故多带模拟可在每次位访问时扫描其工作区，
成本 $V^{O(1)}$ 被吸收进 $b$。

固定正有理数 $\eta<1/(12A)$，并对 $n\ge2$
置 $s_n=\lfloor\eta n\lfloor a(n)\rfloor\rfloor$。
该阈值可用 $n$ 的二进制长度与
$\lfloor\log_2 n\rfloor$ 的整数平方根计算。
设对所有足够大的 $n$ 有 $C_W(dn)\le s_n$。
为判定 $H$：拒绝非一元输入。在 $1^n$ 上，
猜测一个 $dn$ 个下标位、至多 $s_n$ 个门的电路。
把解码器的三个变量标号馈入该电路的三份副本，
并检验所索引的子句是否被违反。
若检验电路不可满足则接受。
当 $1^n\in H$ 时，在长度 $dn$ 处计算 $W$
的电路给出一个接受的猜测。
当 $1^n\notin H$ 时，每个猜测都违反一个子句。

检验电路有 $q$ 个输入与至多 $3s_n+O(n)$ 个门。
上述编码为每个输入与门各用一个变量，
每门至多四个宽度至多三的子句以强制其真值表，
再用一个单元子句钉住输出。
用未用变量垫齐后得 $U_n=3s_n+Bn$ 个变量与至多
$5U_n$ 个子句（$B$ 为固定整数）。
由 $U_n=(3\eta+o(1))n a(n)$ 与 $a(U_n)\sim a(n)$，
检验耗时 $2^{(3A\eta+o(1))n}$。
猜测与构造电路花费多项式时间，
故对一切足够大的 $n$，$H$
可在时间 $2^{n/2}$ 内非确定性地判定。
有限多个例外长度可以硬编码，
这与 $H$ 的选取矛盾。
于是对无穷多个 $n$ 有 $C_W(dn)>s_n$。
由 $s_n\sim\eta n\sqrt{\log_2 n}$ 与 $d$ 固定，
取 $L=W$ 与任何 $0<c<\eta/(d\sqrt{\log 2})$
即证明 \eqref{eq:quantitative-circuit-lower-bound}。
\end{proof}


\subsection{求逆常数输出局部性的布尔映射}\label{subsec:local-inversion}
固定常数 $d$。考虑布尔映射
$G:\{0,1\}^n\to\{0,1\}^m$，其第 $i$ 个输出位形如
\[
 G(\vec{x})_i=g_i(\vec{x}|_{S_i}),
 \qquad S_i\subseteq[n],\quad |S_i|\le d.
\]
$G$ 的输入描述列出集合 $S_i$ 与布尔函数 $g_i$
的真值表。于是每个输出只检查常数个输入位，
尽管集合 $S_i$ 可以任意重叠。
给定目标 $\vec{b}\in\{0,1\}^m$，
求逆即要求任意满足 $G(\vec{x})=\vec{b}$
的 $\vec{x}$。

\begin{corollary}\label{cor:consequence-local-inversion}
对每个固定的 $d$ 与 $n\ge20$，
数 $|G^{-1}(\vec{b})|$ 可以被精确计算，
且当原像存在时可以找到其一，
确定性时间为
\[
 (n+m)^{O_d(1)}2^{O_d(\rho(n))}
\]
，空间为 $(n+m)^{O_d(1)}$。
\end{corollary}
\begin{proof}
对每个 $i$ 与每个满足 $g_i(\vec{\alpha})\ne b_i$
的局部赋值 $\vec{\alpha}\in\{0,1\}^{S_i}$，
加入子句
\[
 \bigvee_{j\in S_i}\ell_{j,\alpha_j},
 \qquad
 \ell_{j,0}=X_j,\quad \ell_{j,1}=\neg X_j.
\]
该子句恰禁止局部赋值 $\vec{\alpha}$。
若 $S_i$ 为空，其约束是常值的，可直接检验。
剩余子句的合取是原 $n$ 个声明变量上的 CNF，
宽度至多 $d$、子句至多 $2^d m$ 个。
其可满足赋值恰为 $G^{-1}(\vec{b})$；
不需要辅助变量或重数修正。
以 $k=\max\{2,d\}$ 应用定理~\ref{thm:main}。
通常的逐变量自归约用至多 $n$ 次进一步计数查询找到一个
原像，这只改变多项式因子。
每次查询可在下一次之前丢弃，
保持多项式空间。
\end{proof}

条件计数还允许在原像集非空时采样一个完全均匀的原像，
期望时间界相同（差多项式因子）~\cite{JVV86}。
特别地，该推论适用于 Goldreich 候选单向函数的常数
局部性版本，其在输入位的重叠子集上求值局部谓词
~\cite{Gol00}。对多项式个输出位，
它给出对像中每个输出的 $2^{o(n)}$ 时间求逆算法，
排除这些族的 $2^{\Omega(n)}$ 求逆难度。
这并不排除对多项式时间算法的单向性。

\section{精确计数的更尖锐界}\label{app:quantitative}

我们通过收缩布尔张量网络证明定理~\ref{thm:main}。
两个运行时间反映网络中的线数。有 $m\le Dn$
个子句的公式给出 $O_{k,D}(n)$ 条线的网络。
对任意公式，我们使用第~\ref{sec:hashes} 节的奇偶检验。
引理~\ref{lem:total-order} 的次数界使我们能用
$O_k(n\log\log n)$ 条线为多项式提升之乘积中保留的
单项编下标。在证明其所用的收缩定理之前，
我们先给出两种编码。

\subsection{布尔张量收缩}
\label{subsec:tensor-interface}

布尔张量网络是一个二部多重图 $G$，
其每条边承载一个位。每个顶点 $v$
有一个关联边的有序列表 $\delta(v)$
与一个张量 $T_v$：对这些边的每个布尔赋值，
指派一个环元素。平行边有不同端口。
网络的\emph{收缩}是
\[
 Z=\sum_{x\in\{0,1\}^{E(G)}}\prod_{v\in V(G)} T_v(x|_{\delta(v)}).
\]
孤立顶点贡献其标量分量。
我们以边数度量网络的规模。

\begin{theorem}[布尔张量收缩]\label{thm:tensor-contraction}
固定正常数 $C,A,b$。设 $s\ge2$ 为整数，
$\mathcal L$ 为一个至多 $Cs$ 条边、最大度至多
$\lceil s^{9/10}\rceil$ 的二部张量网络的图与张量描述的
总长度。设顶点数被 $s$ 的固定多项式界定，且分量属于
\[
 \mathcal B_{P_0,\ell}
   =(\mathbb Z/2^{P_0}\mathbb Z)[V]/(V^\ell),
 \qquad 1\le P_0,\ell\le Cs.
\]
设一致地对每个顶点与其端口的每个布尔选取，
其分量可以从描述在时间 $(s+\mathcal L)^b2^{A s^{9/10}}$
与空间 $(s+\mathcal L)^b$ 内计算。
则收缩可以确定性地在
\[
 (s+\mathcal L)^{O(1)}
  2^{O(s/\sqrt{\log s})}
\]
次位运算与 $(s+\mathcal L)^{O(1)}$ 比特空间内计算。
常数只依赖假设中的固定界。特别地，容许 $\ell=1$。
\end{theorem}

第~\ref{app:tensor} 节的算法用对线的递归替代单一全局
Taylor 展开。每次调用时它变形线的等同关系，
使在选定的插值节点处只有同一组内的端口相互作用。
在该节点附近，组间的相互作用在称为模（mode）
的变量中展开。权重 $j$ 的模的系数被 $2^j$ 整除，
故只有总权重低于工作精度的单项才有贡献。
随节点的二进赋值调整 Hermite 重数即可界定该精度。
有 $M$ 条线、保留 $u$ 个赋值层时，模单项的个数是
$2^{O(M/\sqrt{u})}$，且每个单项分解为至多 $M/8$
条线的收缩。在递归深度 $d$ 处，
$M=O(s8^{-d})$、$u=\Theta((\log s)/2^d)$。
诸指数因此几何求和为 $O(s/\sqrt{\log s})$。

\subsection{固定子句密度}\label{subsec:sharp-fixed-density}

\begin{proof}[定理~\ref{thm:main} 固定密度界的证明]
给每个子句一个顶点，承载其布尔满足表，
子句中每个不同变量一个端口。重言式可以丢弃，
空子句给出计数零。对有 $d\ge1$ 个端口的变量，
要求这些位的取值一致。$d=1$ 时挂上一元张量
$(1,1)$；$d=2$ 时挂上一个二元相等张量。
$d\ge3$ 时取一棵有 $d$ 个外端口、内部顶点为三价相等
张量的树。这样的树有 $d-2$ 个内部顶点，
其端口上的公共取值唯一延拓到其全部内部边。

子句端口至多 $km$ 个。图有 $O_k(n+m)$ 条边与顶点，
最大度至多 $\max\{k,3\}$。把每条边用二元相等张量
细分一次。旧顶点与新顶点构成二部网络的两侧。
细分边上两位必须相等，故这保持收缩。
出现变量上的每个可满足赋值贡献一，
其余每个赋值贡献零。

取 $s=n+m$，在 $\mathbb Z/2^{n+1}\mathbb Z$
上应用定理~\ref{thm:tensor-contraction}，
取 $P_0=n+1$、$\ell=1$。对固定的 $k$，
局部表规模有界。次数与精度界对所有足够大的 $s$
成立，有界的实例可直接枚举。
若有 $u$ 个声明变量未被使用，
把收缩乘以 $2^u$（模 $2^{n+1}$）。
其最小非负剩余是 $\#F$，因为 $0\le\#F\le2^n$。
这包括空公式：空网络的收缩为一。
时间为 $s^{O_k(1)}2^{O_k(s/\sqrt{\log s})}$，
空间为多项式。当 $m\le Dn$ 时，
$n\le s\le(1+D)n$，给出指数
$O_{k,D}(n/\sqrt{\log n})$。
\end{proof}

\subsection{任意子句密度}\label{subsec:sharp-dense}

第~\ref{sec:hashes} 节的稀疏化与贪心行选取把该情形
化归为稀疏多项式方程的公共零点计数。
我们保留第~\ref{sec:setup} 节的奇偶提升。
其次数界使我们能用 $O(n\log\log n)$ 条布尔线
以「检验--变量」关联的子集为单项编码。

\begin{lemma}[用收缩计数稀疏多项式方程组]
\label{lem:tensor-sparse-query}
固定整数 $k,c\ge1$。设 $P_1,\ldots,P_a$
为 $n\ge20$ 个声明布尔变量上的 $\mathbb F_2$
多项式。设 $a\le cn$，且每个 $P_i$
列为至多 $\lceil\log_2 n\rceil^c$
个次数至多 $k$ 的单项。公共零点个数可确定性地在
$n^{O_{k,c}(1)}2^{O_{k,c}(n\log\log n/\sqrt{\log n})}$
次位运算与 $n^{O_{k,c}(1)}$ 比特空间内计算。
\end{lemma}

\begin{proof}
令 $\lambda=\lceil\log_2 n\rceil$、
$\nu=\lceil\log_2\lambda\rceil$。
在每个单项内把变量的正幂替换为一次幂，
并把重复单项按模二合并。检验 $i$
的 scope 由其剩余单项中的变量组成，
规模至多 $k\lambda^c$。
设 $N_0$ 为 scope 大小之和，
对属于检验 $i$ 的变量 $v$ 把对 $(i,v)$ 计一次。
于是 $N_0\le ak\lambda^c$。
以一切可能方式固定属于多于
\[
 D_*=\max\{1,\lceil N_0/n\rceil\}\lceil\sqrt{\lambda}\rceil
\]
个 scope 的变量。这样的变量至多 $n/\sqrt{\lambda}$
个。其赋值把原计数分拆为至多 $2^{n/\sqrt{\lambda}}$
个分支计数。在每个分支上代入固定位，
并再次合并重复单项。
检验 $i$ 的自由 scope $I_i$ 规模仍至多 $k\lambda^c$，
且每个自由变量出现在至多
$D_*=\lambda^{O_{k,c}(1)}$ 个 scope 中。
不在任何 scope 中的自由变量保持声明。
常值一方程使分支计数为零。可以丢弃常值零方程；
若都不剩，分支的答案为其自由变量数的二的幂。
下文中 $a$ 表示剩余方程数，$n$ 表示原变量数。

\paragraph{关联上的多项式。}
给每对 $(i,v)$（$v\in I_i$）自己的未定元 $z_{iv}$。
写 $z_i=(z_{iv})_{v\in I_i}$。
把归一化分支多项式的剩余单项以整系数求和
（用 $z_{iv}$ 代 $x_v$），作成
$\widetilde L_i(z_i)$。我们使用环
\[
 \mathcal B=(\mathbb Z/2^p\mathbb Z)[V]/(V^\ell),
 \qquad p=\ell=n+1 .
\]
取 $h(Y)=Y(Y-1)$ 与引理~\ref{lem:rational-idempotent}
的整系数多项式 $e_j$。把其系数模 $2^p$
化简，并作成因子
\[
 B_i(z_i,V)=\sum_{j=0}^{\ell-1}
   e_j(\widetilde L_i(z_i))\,h(\widetilde L_i(z_i))^j V^j .
\]
设 $\overline B_i$ 为其在 $z_i$ 中的多重线性代表：
在 $V$ 的每个系数中把关联未定元的每个正幂替换为一次幂
而得。它在 $z_i$ 的每个布尔赋值上与 $B_i$ 一致。
在 $\mathcal B[z]$ 中定义
$F(z,V)=\prod_i\overline B_i(z_i,V)$。
不同检验的关联未定元互不相交，
故即使在把变量 $v$ 的全部副本代入同一布尔 $x_v$
之前，$F$ 也是多重线性的。

$\overline B_i$ 中 $V^j$ 的系数的关联次数至多
$k(1+2j)$：由 $\deg e_j\le1$、$\deg h=2$，
且多重线性化只会降低单项次数。
$\mathcal B$ 中的乘法只保留满足
$\sum_i j_i\le\ell-1$ 的项。
因此 $F$ 的每个关联单项的次数至多
\begin{equation}\label{eq:sharp-incidence-degree}
 K=k\bigl(a+2(\ell-1)\bigr)=O_{k,c}(n).
\end{equation}
设 $\mathcal S_\ell$ 在 $\mathcal B$ 中对
$1,V,\ldots,V^{\ell-1}$ 的系数求和。
在 $z_{iv}=x_v$ 处，多重线性化的因子
$\overline B_i$ 对归一化分支多项式与 $B_i$ 一致。
整数 $\widetilde L_i((x_v)_{v\in I_i})$
与原 $P_i$ 在 $x$ 及固定位处的取值有相同奇偶性。
对该整数应用引理~\ref{lem:total-order}
表明分支计数模 $2^p$ 同余于
\begin{equation}\label{eq:sharp-branch-residue}
 \mathcal S_\ell\left(\sum_x F\bigl((z_{iv}=x_v)_{iv},V\bigr)\right),
\end{equation}
其中 $x$ 遍历全部自由布尔赋值。
我们在截断环中相乘检验因子之后应用系数求和映射。

\paragraph{为单项编下标。}
把检验与自由变量分别划分为
$J=\lceil n^{1/4}\rceil$ 个组，
两侧各组的大小相差至多一。
记 $\mathcal T_g$ 为检验组，$\mathcal X_h$
为变量组。每组含 $O_c(n^{3/4})$
个对象，故其关联数在任一侧至多为
$n^{3/4}\lambda^{O_{k,c}(1)}$。
设 $N_{gh}$ 为从 $\mathcal T_g$ 到 $\mathcal X_h$
的关联数，$N=\sum_{g,h}N_{gh}\le N_0$。

对关联单项，设 $d_{gh}$ 为其连接组 $g,h$
的被选关联数。枚举满足 $0\le d_{gh}\le N_{gh}$
且 $\sum_{g,h}d_{gh}\le K$ 的整数数组 $(d_{gh})$。
至多有
\[
 (K+1)^{J^2}=2^{O_{k,c}(\sqrt n\log n)}
\]
个数组，可用 $J^2$ 个计数器枚举。
对固定数组，每一对有
\[
 C_{gh}=\binom{N_{gh}}{d_{gh}}
\]
个指定大小的子集。在两个组顶点之间画
$\lceil\log_2 C_{gh}\rceil$ 条布尔线，
按字典序为其编下标。$C_{gh}=1$
时两个端点都用唯一子集（可能是整个关联集），
不用线。任何未用的位串给出张量分量零。
解码有效下标时，按次序取关联，
用二项系数统计「含或不含每条关联」的可能补全。
这些有多项式位长，故解码使用多项式时间与空间。

全部对上子集的并指定 $N$ 个未定元上次数至多 $K$
的一个无平方因子单项。计入该次数界的全部单项，
由隔板法得上界 $\binom{N+K}{K}$。
因此该数组在网络中的线数 $M$ 满足
\begin{align}
 M &\le J^2+\log_2\prod_{g,h} C_{gh}
      \le J^2+\log_2\binom{N+K}{K} \notag\\
   &\le J^2+K\log_2\!\left(\frac{\mathrm e(N+K)}{K}\right)
      = O_{k,c}(n\nu).
 \label{eq:sharp-wire-budget}
\end{align}
这里 $K=\Theta_{k,c}(n)$，$N\le n\lambda^{O_{k,c}(1)}$。
且 $\lceil\log_2 C_{gh}\rceil\le N_{gh}$，
故每个顶点的度至多是其关联数
$n^{3/4}\lambda^{O_{k,c}(1)}$。

\paragraph{张量分量与收缩。}
线位的一个有效赋值在每个检验处确定一个子集
$S_i\subseteq I_i$。写
$z_{i,S_i}=\prod_{v\in S_i}z_{iv}$。
检验组 $g$ 处的分量是其所示单项在
$\prod_{i\in\mathcal T_g}\overline B_i(z_i,V)$
中的系数。由 $z_i$ 互不相交，这在 $\mathcal B$
中等于
\[
 \prod_{i\in\mathcal T_g}[z_{i,S_i}]\overline B_i(z_i,V).
\]
对变量组 $h$，对 $v\in\mathcal X_h$
置 $b_v=|\{i:v\in S_i\}|$。其分量为
\[
 \sum_{x\in\{0,1\}^{\mathcal X_h}}
        \prod_{v\in\mathcal X_h}x_v^{b_v}
   =2^{|\{v\in\mathcal X_h:b_v=0\}|}\pmod{2^p}.
\]
事实上，正指数要求 $x_v=1$，零指数使 $x_v$ 自由。
这包括不在任何检验 scope 中的变量。

$F$ 的每个单项恰出现在一个数组中，
且在每个组对处恰有一个有效下标。
其在 $F$ 中的系数是检验组系数之积。
把含同一 $v$ 的关联置为相等后，
其在布尔赋值上的和是变量组分量的乘积。
由 \eqref{eq:sharp-incidence-degree}，
对这些数组求和即穷尽全部单项。
因此 \eqref{eq:sharp-branch-residue} 中
$\mathcal S_\ell$ 之内的 $\mathcal B$
元素是这些二部网络收缩之和。

为求值一个检验侧系数，设 $\mathbf1_T$
为 $I_i$ 上支撑为 $T\subseteq S_i$ 的布尔向量。
在 $\mathbf1_T$ 处，多重线性多项式求值为以 $T$
的子集为下标的系数之和，且在该处与 $B_i$ 一致。
反演该子集和给出
\[
 [z_{i,S_i}]\overline B_i(z_i,V)=
 \sum_{T\subseteq S_i}(-1)^{|S_i|-|T|}B_i(\mathbf1_T,V)
 \qquad\text{（在 }\mathcal B\text{ 中）}.
\]
每项花费多项式次位运算。从列出的单项求值
$\widetilde L_i(\mathbf1_T)$，
使用引理~\ref{lem:rational-idempotent}
的多项式 $e_j$（多项式时间内可构造）。
把 $B_i(\mathbf1_T,V)$ 模 $2^p$
的 $\ell$ 个系数存储。
需处理的子集至多 $2^{|S_i|}\le 2^{k\lambda^c}$ 个。
一次对一个检验求和，并在组上把所得系数在 $\mathcal B$
中相乘。这在 $2^{\lambda^{O_{k,c}(1)}}$
次位运算与多项式空间内求值一个检验组分量。
变量侧分量与对无效线下标的检验花费多项式时间。
图、度数组与分量过程都有多项式规模描述。

以规模参数 $s=n\nu$、根精度 $p=n+1$、截断阶
$\ell=n+1$ 应用定理~\ref{thm:tensor-contraction}。
由 \eqref{eq:sharp-wire-budget}，图有 $2J$
个顶点与 $O_{k,c}(s)$ 条线。
对一切足够大的 $n$，其最大度至多
$\lceil s^{9/10}\rceil$，其精度与截断阶为 $O(s)$，
其分量求值时间至多 $2^{O_{k,c}(s^{9/10})}$。
每次收缩于是花费
$n^{O_{k,c}(1)}2^{O_{k,c}(s/\sqrt{\log s})}$
次位运算与多项式空间。这里
$s/\sqrt{\log s}=O(n\nu/\sqrt\lambda)$。
高关联分支与度数组把运行时间至多乘以
\[
 2^{\,n/\sqrt{\lambda}+O_{k,c}(\sqrt n\log n)}.
\]
该指数的两项都是 $O_{k,c}(n\nu/\sqrt{\lambda})$。
依序处理分支与数组。在每个分支上，
把 $\mathcal B$ 中的收缩相加并应用
$\mathcal S_\ell$。把分支剩余模 $2^{n+1}$
相加给出原计数，因为它介于 $0$ 与 $2^n$ 之间。
这证明了所述界，有界的例外规模用直接枚举处理。
\end{proof}

\begin{proof}[定理~\ref{thm:main} 任意密度界的证明]
如第~\ref{subsec:general-reduction} 节那样
在原 $n$ 个声明变量上归一化公式。
空子句与无子句的公式有直接答案。
对 $n\ge20$，以
$\sigma=1/\lceil\sqrt{\lambda}\rceil$
（$\lambda=\lceil\log_2 n\rceil$）
应用引理~\ref{lem:sparsify}。
树有 $n^{O_k(1)}2^{O(n/\sqrt{\lambda})}$
个节点与多项式深度。
每个叶的出现界为
$B=O_k((\sigma^{-1}\log(2/\sigma))^{k-1})=\lambda^{O_k(1)}$，
非空子句至多 $Bn$ 个。

对有 $m_G$ 个子句的非平凡叶，使用
引理~\ref{lem:sparse-checks} 的初始奇偶校验与
引理~\ref{lem:greedy-hashes} 的贪心过程。
每个初始行的支撑大小至多 $D_0$，
每个附加行的列表长度为 $P$。
两个界都是 $\lambda^{O_k(1)}$。
由 \eqref{eq:hash-query}，每行给出一个次数至多 $k$、
至多 $2^k\max\{P,D_0\}$ 个单项的方程，
且每个被查询的方程组至多 $6n$ 个方程。
对一切足够大的 $n$，单个常数 $c=c(k)$
同时界定方程数与变量数之比以及单项界中 $\lambda$
的指数。在每个贪心步，
用引理~\ref{lem:tensor-sparse-query}
对每个候选行 $h_\gamma$ 计算 $N(A;h_\gamma)$，
并追加取得最小值的一个。
求值器也适用于其方程组仍接受不可满足赋值的被丢弃候选。
对所选矩阵作最后一次调用给出 $N(A)$，
由引理~\ref{lem:greedy-hashes} 它等于叶计数。

构造初始检验与枚举候选在求值之外花费
$2^{\lambda^{O_k(1)}}$ 次位运算。
每叶至多 $3n(2m_G)^P+1=2^{\lambda^{O_k(1)}}$
次求值。由于 $\lambda$ 的每个固定幂都是
$o(n\log\log n/\sqrt{\log n})$，
这些成本与叶数都被吸收进所述时间界
\[
 (n+m)^{O_k(1)}2^{O_k(n\log\log n/\sqrt{\log n})}.
\]
这里 $m$ 计入输入子句。深度优先遍历树，
一次处理一个查询。叶解集互不相交，
故其精确计数用 $(n+1)$ 比特累加器相加。
连同每次查询与每条树路径的多项式空间，
这给出所述的多项式空间界。
\end{proof}

单一算法同时达到两个界。
取 $\lambda=\lceil\log_2 n\rceil$、
$\nu=\lceil\log_2\lambda\rceil$：
当 $m\le n\nu$ 时用子句网络，否则用奇偶查询。
第一分支中 $n\le s=n+m=O(n\log\log n)$，
故即使密度增长，其指数也是
$O_k(n\log\log n/\sqrt{\log n})$。
对每个固定的 $D$ 与一切足够大的 $n$，
每个满足 $m\le Dn$ 的输入都走该分支，
$s=O_D(n)$。

\subsection{收缩定理的证明}
\label{app:tensor}

在定理~\ref{thm:tensor-contraction} 的证明中，
$\log$ 表示以二为底的对数，$\ln$ 表示自然对数。

\subsubsection{系数环与精确算术}
\label{subsec:tensor-arithmetic}

写 $\lambda=\lceil\log_2 s\rceil$、$e=4\lambda$。
取一个次数为 $e$ 的首一多项式
$f\in\mathbb Z[U]$，其首项系数之下的系数在 $\{0,1\}$
中，且其模二约化不可约。
这可以在 $s$ 的多项式时间内确定性完成：
枚举 $2^e$ 个候选，并用 $\mathbb F_2$ 上每个次数为正、
至多 $e/2$ 的首一多项式检验可除性。
有限域有每个正次数的不可约多项式~\cite{vzGG13}。
置
\[
 \mathcal O=\mathbb Z_{(2)}[U]/(f),\qquad
 \mathcal K=\mathbb Q[U]/(f),\qquad
 \mathcal A=\mathcal O[V]/(V^\ell),\qquad
 \mathcal A_p=\mathcal A/(2^p).
\]
这里 $\mathbb Z_{(2)}$
由分母为奇数的有理数组成。
模二不可约蕴含在 $\mathbb Q$ 上不可约，
故 $\mathcal K$ 是域。
对 $\mathcal K$ 的非零元素，设 $v_2$
为其在 $1,U,\ldots,U^{e-1}$
中坐标的寻常二进赋值的最小值。
$\mathcal O$ 模二的约化是域 $\mathbb F_{2^e}$。
对最小坐标赋值取商因此表明
\[
 v_2(xy)=v_2(x)+v_2(y)\qquad(x,y\ne0).
\]
$\mathcal O$ 中模二非零的每个元素都是单位，
因为其乘法矩阵的行列式在 $\mathbb Z_{(2)}$
中可逆，故逆矩阵的分量在该环中。

$\mathcal A_p$ 中的算术在 $s,p,\ell$
上花费多项式时间与空间。
我们使用基 $1,U,\ldots,U^{e-1}$
与 $1,V,\ldots,V^{\ell-1}$，
并把每个坐标表示为它在零与 $2^p-1$
之间的整数剩余。标量单位的求逆可先模二求逆
再用 Newton 提升完成，或用其乘法矩阵完成，
使用精确的多项式与整数运算~\cite{vzGG13}。
$\mathcal A$ 中元素的可除性与精度总是逐坐标理解。

把 $2^e$ 个剩余域元素的二值坐标提升，
即得两两之差为单位的插值点。
因此，在这些点的任意所需数量上的 Lagrange
插值在 $\mathcal A_p$ 上有效。
我们将使用少于 $s^4$ 个点。
初始分量属于 $\mathcal B_{P_0,\ell}$
的收缩，在 $\mathcal A_{P_0}$ 中计算之后仍属于该子环，
因为它是那些分量的多项式。

精度增大时，我们逐个提升每个张量分量。
对原来按模 $2^p$ 指定的分量，
其\emph{典范提升}是由其全部坐标的整数代表得到的
$\mathcal A$ 中元素。
典范提升不必保持加法或乘法。
每次递归调用按该调用的精度对其指定的张量作收缩。
为定义其变形，它先作出这一逐分量的提升选取。
该调用中后续的全部求值都指称这些相同的提升。
更明确地，每个张量描述都带有其创建精度 $p_c$
（原始张量为 $P_0$）。
在其后的每个精度处，其分量按模 $2^{p_c}$
求值，并使用这些相同的整数代表。
精度沿递归路径严格递增，
故这些代表在其后每个精度处也是典范代表。
特别地，一个包张量在整个递归中保持其创建时的代表。

\subsubsection{参数与碰撞装箱}
\label{subsec:tensor-packing}

下文全部估计对足够大的 $s$ 成立，
阈值只依赖固定常数。
较小规模的固定有限范围用直接求和处理。
参数不等式可以直接检验。置
\[
 H=\left\lceil\frac{s}{\lfloor\sqrt\lambda\rfloor}\right\rceil,
 \qquad r=64,\qquad \varepsilon=\frac1{128\cdot1024},
 \qquad \Delta=\lceil s^{9/10}\rceil.
\]
有 $M\le H$ 条边的调用是叶，直接求值。
在深度 $d$ 的活动调用处（根在深度零），使用
\begin{equation}\label{eq:tensor-parameters}
 u=u_d=\left\lfloor\frac{\varepsilon\lambda}{2^d}\right\rfloor,
 \qquad a=a_d=2\sum_{j=0}^d u_j,\qquad
 a_{-1}=0,\qquad L=2^a,\qquad D_{\rm arr}=r^{2a}.
\end{equation}
我们将安排每个孩子至多 $M/8$ 条边、最大度至多
$\Delta$。这些断言蕴含，在每条活动路径上，
\[
 H<M\le Cs\,8^{-d},\qquad
 d\le\frac{\log_2\lambda}{6}+O_C(1).
\]
因此，对一切足够大的 $s$，
\begin{equation}\label{eq:tensor-parameter-bounds}
 u=\Theta_C(\lambda/2^d),\quad
 \lambda^{4/5}\le u\le M,\quad
 a\le4\varepsilon\lambda,\quad
 L,D_{\rm arr}\le s^{1/1000}.
\end{equation}
特别地，在每次活动调用处
\begin{equation}\label{eq:tensor-packing-conditions}
 M\ge2D_{\rm arr}^{\,2}\Delta,\qquad
 M\ge L D_{\rm arr},\qquad 2L\le\Delta.
\end{equation}
例如，第一个不等式右端是
$O(s^{9/10+2/1000})=o(H)$。
其余不等式同理。

\begin{lemma}[装箱成包]\label{lem:tensor-packets}
在 \eqref{eq:tensor-packing-conditions} 下，
把网络每侧的顶点用 $[r]^a$ 的元素标记。
存在可确定性构造的标记与把边嵌入 $L$
个槽位的包的方案，满足以下性质。
\begin{enumerate}
\item 包至多 $3M/L$ 个。每条边在其行端与列端占据
      同一槽位。未用的槽位没有端口。
\item 在每个包的每个坐标中，行符号与列符号的无序对
      是固定的。当两符号不同时，槽位号的第 $i$
      位指定它们在坐标 $i$ 中的取向
      （$0\le i<a$）。
\item 对每个坐标与每个符号，每侧属于带该符号顶点的
      实端口至多 $3M/r$ 个。
\end{enumerate}
构造在 $s+\mathcal L$ 上使用多项式时间与空间。
\end{lemma}
\begin{proof}
对每对标号，设 $X_\gamma$ 为具有该对的边数。
独立均匀的顶点标号给出
\[
 \mathbb E\sum_\gamma (X_\gamma-M/D_{\rm arr})^2\le2M\Delta.
\]
事实上，把该期望展开为有序边对上的标号一致概率之和
减去 $1/D_{\rm arr}$。端点不交的对贡献零。
共享一个端点的有序对（含相等的边）至多 $2M\Delta$
个，每个贡献至多一。
条件期望给出具有该界的标记。
为实现之，尝试下一个顶点的 $r^a$ 个标号，
并通过枚举边对计算条件期望。
每个概率只涉及至多四个顶点，
故这是精确的多项式时间计算。
这就是通常的条件期望法~\cite{Raghavan88}。

每个负载至多 $M/D_{\rm arr}+\sqrt{2M\Delta}$。
把全部对—标类垫到共同的槽数 $H_{\rm pad}$
（最大负载向上取整到 $L$ 的倍数）。
假设蕴含
\[
 H_{\rm pad}\le M/D_{\rm arr}+\sqrt{2M\Delta}+L
              \le3M/D_{\rm arr}.
\]
固定每个坐标中符号的无序对。
若有 $a'$ 个坐标的符号不等，其取向给出 $2^{a'}$
个对标号。把这些垫后的槽位分装成包：
每包用每个这样的标号取 $2^{a-a'}$ 个槽位。
这是可能的，因为 $H_{\rm pad}$ 被 $L$ 整除。
把一个包的位置从 $0$ 到 $L-1$ 编号。
其第 $i$ 位在第 $i$ 对不等时指定取向，
相等时指定副本。每个包内使用同一取向约定。

包数为 $H_{\rm pad}D_{\rm arr}/L\le3M/L$。
在固定坐标中，给定符号在两侧的全部垫后对标号中恰占
$1/r$。其实端口数因此至多
$H_{\rm pad}D_{\rm arr}/r\le3M/r$。
\end{proof}

\subsubsection{Cauchy 变形与整的局部坐标}
\label{subsec:tensor-deformation}

固定一个活动调用、其包与其逐分量提升的张量。
使用
\[
 S=\{-L,-L+1,\ldots,L-1\},\qquad
 \mathcal C(T)=\prod_{t\in S}(T-t),\qquad \alpha=2^{a+1}.
\]
在一个包内，对其行槽位 $x$ 与列槽位 $y$，置
\begin{equation}\label{eq:tensor-cauchy}
 W_{xy}(T)=\frac{\mathcal C(T)}
                  {\mathcal C'(0)(T-(y-x))}.
\end{equation}
分母中的因子消去分子中的一个因子，因为 $y-x\in S$。
故这是 $\mathbb Q$ 上次数 $2L-1$ 的多项式。
把不同包之间的分量置零，并限制到实端口。
在原来的边等同关系下，$W(0)=I$：
对角分量为二，其他每个分量都有一个未消去的零因子。

设 $A(Y)$ 与 $B(X)$ 分别是提升后的行张量与列张量
在独立端口赋值 $Y,X\in\{0,1\}^M$ 上的乘积。
写 $l_0(z)=1-z$、$l_1(z)=z$。定义
\begin{equation}\label{eq:tensor-deformed-contraction}
 F(T)=\sum_{X,Y\in\{0,1\}^M} A(Y)B(X)
                  \prod_{\text{行端口 }x}l_{Y_x}((W(T)X)_x).
\end{equation}
其系数落在 $\mathcal K[V]/(V^\ell)$ 中，
次数至多 $M(2L-1)$，而 $F(0)$
是提升后张量的收缩。特别地，
把 $F(0)$ 模该调用的精度化简即得所需答案。

\begin{lemma}\label{lem:tensor-integrality}
对每个 $t\in S$，$W(t+\alpha z)$ 的分量属于
$\mathcal O[z]$，且 $F(t+\alpha z)$ 属于
$\mathcal A[z]$。
给定 $\zeta$ 模 $2^p$ 的 $e$ 个坐标，
每个分量 $W_{xy}(t+\alpha\zeta)\bmod 2^p$
可以用 $s,p$ 上多项式次位运算计算。
\end{lemma}
\begin{proof}
网格 $S$ 是模 $2^{a+1}$ 的完全剩余系。
对 $S$ 中 $t\ne q$，$v_2(t-q)\le a$，
且每个赋值处非零差的个数与 $t$ 无关。
故 $v_2(\mathcal C'(t))=v_2(\mathcal C'(0))$。
在
\[
 \frac{\mathcal C(t+\alpha z)}{\mathcal C'(0)}
  =\frac{\alpha z\prod_{q\in S\setminus\{t\}}
                         (t-q+\alpha z)}{\mathcal C'(0)}
\]
中，每个以 $q\ne t$ 为下标的因子逐系数被
$2^{v_2(t-q)}$ 整除。对这些幂与分母相消后，
所示多项式被 $2^{a+1}$ 整除。
为得 \eqref{eq:tensor-cauchy}，只省略以 $y-x$
为下标的因子。若 $y-x=t$，这去掉 $\alpha z$
并留下整系数多项式。否则 $\alpha z$
保留，而被省略因子有二进容量
$2^{v_2(t-y+x)}$（$v_2(t-y+x)\le a$），
故结果仍被二整除。代入
\eqref{eq:tensor-deformed-contraction}
即证明关于 $F$ 的断言。

对求值，先把每个以 $q\ne t$ 为下标的因子
$t-q+\alpha z$ 除以已知整数 $2^{v_2(t-q)}$。
所得一次多项式的两个系数都是整数。
如上对分母消去其余二的显式幂，
并把其奇部模 $2^p$ 求逆。
现在可以直接在所给 $\zeta$
的剩余处求值该乘积（含 $\zeta=0$）。
输入有 $ep$ 位，乘积有 $O(L)$ 个因子。
全部坐标运算在 $s,p$ 上花费多项式时间。
\end{proof}

\subsubsection{锥形 Hermite 恢复}
\label{subsec:tensor-hermite}

下一条引理单列出精度估计。
其陈述中全部幂级数截断都是 $z$
中的通常截断。

\begin{lemma}[锥形 Hermite 恢复]\label{lem:tapered-hermite}
设整数 $u,M$ 满足 $1\le u\le a$、$M\ge u$，
并置 $h=\lceil3M/u\rceil$。
对 $t\in S\setminus\{0\}$ 置
\[
 i=v_2(t),\qquad q_t=q(i)=h\max(0,u-i),\qquad D_t=L/2^i.
\]
设 $F\in(\mathcal K[V]/(V^\ell))[T]$
次数至多 $M(2L-1)$，$F(0)\in\mathcal A$，
且对每个 $t\in S$ 有
$F(t+\alpha z)\in\mathcal A[z]$。
对 $p\ge1$，其值 $F(0)$ 模 $2^p$
可以从
\[
 F(t+\alpha z)\pmod{2^{p_t}},
 \qquad
 p_t=p+8hD_t+8M(a+1),\qquad q_t>0
\]
恢复。只需每个多项式的前 $q_t$ 个系数。
恢复在 $M,L,\ell,e$ 与各精度上花费多项式时间与空间。
\end{lemma}
\begin{proof}
$S$ 中有 $L$ 个奇点，故
$\sum_{t\ne0}q_t\ge Lhu\ge3ML>M(2L-1)$。
且 $h\le4M/u$，故 $q_t\le4M$。
对每个正重数的 $t$ 置
\[
 E_t(T)=\prod_{q\in S\setminus\{0,t\}}
                    \left(\frac{T-q}{t-q}\right)^{q_q}.
\]
Hermite 插值给出
\begin{equation}\label{eq:tensor-hermite-origin}
 F(0)=\sum_{t:q_t>0}E_t(0)
       \left(\left[\frac{F(t+\alpha z)}
                            {E_t(t+\alpha z)}\right]_{<q_t}\right)
                            (-t/\alpha).
\end{equation}
为完整起见，把右端的零换成未定元 $T$，
把 $-t/\alpha$ 换成 $(T-t)/\alpha$。
每个被加项次数小于 $\sum q_t$；
在其自身节点处有指定的 jet，
在其他每个节点处按所需阶为零。
它与 $F$ 之差被
$\prod_{t\ne0}(T-t)^{q_t}$ 整除且次数更小，
故为零。这是 Hermite 插值的通常中国剩余构造
~\cite[Section~5.6]{vzGG13}。

$E_t(t+\alpha z)$ 的每个因子是
$1+\alpha z/(t-q)$ 的幂。它是整的且常项为一，
其截断到 $q_t$ 的倒数也是整的，
故作成 \eqref{eq:tensor-hermite-origin}
中的商不损失精度。

固定 $i=v_2(t)$ 与 $D=D_t$。直接从乘积，
\begin{align}
 v_2(E_t(0))
 &=\sum_{q\ne0,t}q_q\bigl(v_2(q)-v_2(t-q)\bigr)\notag\\
 &=D\sum_{j=1}^{a-i}\frac{j}{2^j}
                           \bigl(q(i+j)-q(i)\bigr)
 \ \ge -6Dh. \label{eq:tensor-hermite-loss}
\end{align}
第二等式的理由：赋值小于 $i$ 的点贡献零。
对 $1\le j\le a-i$，有 $D/2^j$
个 $v_2(q)=i+j$ 的点，同样个数的
$v_2(t-q)=i+j$ 的点。后者有 $v_2(q)=i$，
而这些层穷尽了该赋值处除 $t$
外的全部点。这些计数来自完全剩余网格：
其非零整数差的赋值至多 $a$。
最后，$q(i+j)-q(i)\ge-hj$ 且
$\sum_{j\ge1}j^2/2^j=6$。

次数低于 $q_t$ 的多项式在 $-t/\alpha$
处求值至多再损失 $4M(a+1)$ 个二的幂。
于是把每个商系数换成其模 $2^{p_t}$
剩余的任意提升，只会把其在
\eqref{eq:tensor-hermite-origin}
中的贡献改变一个 $2^p\mathcal A$
中的元素。对这些误差求和不损失更多精度。
重构的和是整的，因为它与整值 $F(0)$
相差一个 $2^p\mathcal A$ 中的元素。

最后这一组合可以用所指定基下的有理算术执行。
乘子 $E_t(0)(-t/\alpha)^j$（$j<q_t$）
是计重数的 $O(a)$ 位有理因子的多项式个乘积。
其分子与分母长度因此在 $M,L,a$
上是多项式的。这样的系数只有多项式个，
故即使其分母的乘积也是足够的多项式位公分母。
相加之后，消去其二进幂，
并把所得整坐标模 $2^p$ 化简。
这证明了算法断言。
\end{proof}

在输入精度为 $p$ 的调用处，
以 \eqref{eq:tensor-parameters} 的参数应用该引理。
对每个活动的 $t$，
通过在剩余域元素的 $M(2L-1)+1$
个互异的二值坐标提升处求值再做 Lagrange 插值，
得到整个多项式 $F(t+\alpha z)$ 模 $2^{p_t}$。
引理~\ref{lem:tensor-integrality}
容许该插值在 $\mathcal A_{p_t}$ 上进行。
由 \eqref{eq:tensor-parameter-bounds}，
点数小于 $s^4$。
下文的求值算法把精度 $p_t$ 传给其孩子。

输入精度与孩子精度满足
\begin{equation}\label{eq:tensor-precision-bound}
 p\le s\lambda^{O(1)}2^{a_{d-1}},
 \qquad
 p_t=O_C(MD_t/u).
\end{equation}
第一个断言在根处成立，因为 $P_0\le Cs$。
沿一条路径，精度增量是 $8hD_t+8M(a+1)$，
其中 $h\le4M/u$、$M\le Cs$、$D_t\le2^{a_d}$。
在 $O(\log\lambda)$ 层深度上求和，
即对每个调用（含叶）证明第一个断言。
对第二个：$q_t>0$ 蕴含 $i<u$，故
\[
 D_t\ge2^{a_d-u}=2^{a_{d-1}+u}.
\]
由 $M>H$ 与 $u\ge\lambda^{4/5}$，
这一额外的因子 $2^u$
支配 $\lambda$ 的每个固定幂。
于是 $p$ 的第一个界以及 $8M(a+1)$
都是 $O_C(MD_t/u)$。
余下的增量 $8hD_t$ 直接具有该界。
由 $a_d\le4\varepsilon\lambda$，
全部精度都保持在 $s$ 的多项式内。

锥形化解释了该估计中的因子 $1/u$。
作为对比，给 $u$ 以下的每一层指派同一重数
$q_*=\Theta(M)$，其余层指派零。
在最后保留的层 $i=u-1$ 处，
恒等式 \eqref{eq:tensor-hermite-loss}
将给出损失
\[
 -q_*D_t\sum_{j=1}^{a-i}\frac{j}{2^j}=-\Theta(MD_t).
\]
而我们的重数每层只变化 $h=\Theta(M/u)$，
给出损失 $O(MD_t/u)$。
$a_d$ 的累积选取提供了吸收继承精度与求值损失所需的
额外因子 $2^u$。
因此在每个保留层、每个递归深度处都成立较小的界。

\subsubsection{组间相互作用的模}
\label{subsec:tensor-modes}

固定活动 $t$ 与 $z$ 的插值取值
$\zeta\in\mathcal O$。
本小节中 $T=t+\alpha\zeta$，
全部计算按模 $2^{p_t}$ 进行。
置 $i=v_2(t)<u$、$D=L/2^i$。
两侧顶点按其在坐标 $i$ 中的符号分组。
$W(T)$ 中两端口同组的分量称为局部的，
保持不变。

这一坐标选取使 $W(t)$ 是局部的。
事实上，除非 $y-x=t$，否则 $W_{xy}(t)=0$。
加一个赋值为 $i$ 的整数会翻转第 $i$ 位，
故存留的对有不同的第 $i$ 位。
当包在该坐标中的两个符号不同时，
这意味着行符号与列符号一致。
符号相等的对本来就同组。
从 $t$ 走到 $t+\alpha\zeta$
引入我们现在展开的交叉项。

在其坐标 $i$ 中两符号不同的包里，
行位置 $x$ 与列位置 $y$ 不同组
当且仅当它们的第 $i$ 位相等。
把列分划成瓦片（tile）
\[
 y=y_0+2^{i+1}j,\qquad 0\le y_0<2^{i+1}.
\]
列组在瓦片上为常数。对跨越到该瓦片的行，
\[
 v_2(t-y_0+x)\le i.
\]
否则 $y_0$ 将与 $x+2^i$ 模 $2^{i+1}$
一致，与其第 $i$ 位相等矛盾。
加 $\alpha\zeta$ 不改变此赋值。
几何展开因此给出
\begin{equation}\label{eq:tensor-mode-expansion}
 W_{xy}(T)=
 \sum_{\substack{b\ge0\\b+1<p_t}}
 \frac{\mathcal C(T)}
      {\mathcal C'(0)(T-y_0+x)}
 \left(\frac{2^{i+1}}{T-y_0+x}\right)^b j^b
 \pmod {2^{p_t}}.
\end{equation}
事实上，第一个分式的赋值至少 $a+1-i\ge1$，
而该比值取 $b$ 次幂前赋值至少为一。
几何余项被 $2^{p_t}$ 整除。
这些表达式可以经分解已知的二的幂、
只用单位逆而计算。

对每个这样的包、瓦片与 $b$，
引入权重 $w_Z=b+1$ 的模变量 $Z$，
其将替换为
\[
 L_Z(X)=\sum_{\text{瓦片中的实列 }y}j^bX_y.
\]
照常 $j^0=1$（含 $j=0$）。
其源是该瓦片的列组，其目标是该包坐标对中的另一组。
公式 \eqref{eq:tensor-mode-expansion}
把每个行型表示为其局部部分加上其目标模的线性组合。
$Z$ 的系数被 $2^{w_Z}$ 整除。
任一权重的变量至多
\begin{equation}\label{eq:tensor-mode-multiplicity}
 (3M/L)2^{i+1}=6M/D
\end{equation}
个。

在仿射因子 $l_{Y_x}((W(T)X)_x)$
的乘积中，模单项 $Z^\beta$
的每个系数都被
$2^{\sum_Zw_Z\beta_Z}$ 整除。
因此只需考虑满足
\begin{equation}\label{eq:tensor-mode-budget}
 \sum_Zw_Z\beta_Z<p_t
\end{equation}
的指数元组。这些元组的个数是
\begin{equation}\label{eq:tensor-partition-count}
 \exp\bigl(O(\sqrt{Mp_t/D})\bigr)
       =\exp\bigl(O_C(M/\sqrt u)\bigr).
\end{equation}
以下是完整的计数估计。对任何 $\theta>0$，
用生成乘积与 \eqref{eq:tensor-mode-multiplicity}
把个数界定为
\[
 \exp\left(\theta p_t+\frac{6M}{D}
                   \sum_{j\ge1}-\ln(1-e^{-\theta j})\right).
\]
展开对数并对几何级数求和，
\[
 \sum_{j\ge1}-\ln(1-e^{-\theta j})
  =\sum_{h\ge1}\frac1{h(e^{\theta h}-1)}
  \le\frac1\theta\sum_{h\ge1}\frac1{h^2}.
\]
取 $\theta=\sqrt{M/(Dp_t)}$ 并用
\eqref{eq:tensor-precision-bound}。
全部模变量及其指数都有多项式规模描述。
满足 \eqref{eq:tensor-mode-budget}
的元组可以按多项式延迟与空间枚举：
对变量排序，逐个在剩余预算内选取每个指数。
每个可行的部分选取都有全零的延续，
故被访问的部分选取数至多为完整元组数乘以变量数
（差一个可加常数）。

\subsubsection{孩子网络及其描述}
\label{subsec:tensor-children}

对 \eqref{eq:tensor-mode-budget} 中的一个元组
$\beta$，为每个组 $g\in[r]$
构造一个孩子网络。把当前网络两侧属于组 $g$
的全部当前顶点放到孩子的同一侧。
保留其端口与所选提升分量（现按模 $2^{p_t}$
化简）。另一侧为每个包放一个顶点，
连边到该包的每个保留端口。
其下标恰为该包中组 $g$ 的行位 $Y$
与列位 $X$。

为定义这个新的包张量，把包中组 $g$
的行因子 $l_{Y_x}$ 相乘，
其自变量用局部列位与形式模表示。
提取以该包中所有目标为 $g$ 的模的指数 $\beta_Z$
的系数。把它乘以
\[
 \prod_{\substack{Z\text{ 在此包中}\\
                  \text{source}(Z)=g}}L_Z(X)^{\beta_Z}.
\]
空乘积为一。若无目标行，则只要某目标指数为正，
提取的系数为零。若 $Z$
的源瓦片没有实列，则 $L_Z(X)=0$，
其幂在 $\beta_Z=0$ 时贡献一，否则贡献零。

每个模恰有一个目标与一个源。
因此系数提取在「包—组」行乘积上可分解，
且每个模的代入只依赖其源组中的位。
在 \eqref{eq:tensor-deformed-contraction}
中展开乘积，随后在组内对独立位求和，
证明精确恒等式
\begin{equation}\label{eq:tensor-recursion-identity}
 F(t+\alpha\zeta)
   =\sum_{\beta:\,\sum w_Z\beta_Z<p_t}
                       \prod_{g=1}^r Z_{g,\beta}
                  \pmod {2^{p_t}},
\end{equation}
其中 $Z_{g,\beta}$ 是刚才定义的孩子的收缩。
被省略的单项由其系数可除性而为零。

由引理~\ref{lem:tensor-packets}，
孩子至多 $6M/r=3M/32<M/8$ 条边。
旧顶点的度不变，新的包顶点至多 $2L\le\Delta$
个端口。于是 \eqref{eq:tensor-parameter-bounds}
所用的不变量归纳成立。

全部新张量都是隐式给出的。
为求值一个包张量分量，展开其至多 $L$
个仿射行因子的乘积。每个因子在 $s$
上有多项式个项（含其局部项）。
枚举这些选取、检验所需指数并对所得项求和花费
\[
 \exp(O(L\log s))
\]
时间与多项式空间。因子 $L_Z(X)^{\beta_Z}$
用 $\mathcal A_{p_t}$ 中通常的反复平方计算。
由 $L\le s^{1/1000}$，该分量成本为
$\exp(O(s^{9/10}))$。

新张量的描述由其包数据、元组 $\beta$、
插值取值与创建精度组成。
这些数据与局部核足以直接计算其分量。
保留的张量从其原描述求值，
或从其自己的创建时描述求值（若它在递归中较早创建）。
其在那个精度处的剩余随后被典范提升到当前精度。
旧描述不可变并按引用共享。
因此分量例程的复合不随深度增长，
全部分量例程保持所需的时间与空间界。
沿一条路径只添加多项式个顶点与多项式多的新描述数据。

\subsubsection{时间与空间界}
\label{subsec:tensor-cost}

\begin{proof}[定理~\ref{thm:tensor-contraction} 的证明]
把原始分量嵌入 $\mathcal A_{P_0}$，
执行上述递归算法。在活动调用处，
对每个正重数的节点与每个所需的、剩余互异的插值取值，
用其孩子收缩计算
\eqref{eq:tensor-recursion-identity}。
插值 $F(t+\alpha z)$ 并应用
引理~\ref{lem:tapered-hermite}。
在叶处，直接对全部布尔边赋值求和。
引理~\ref{lem:tensor-integrality}、
递归恒等式与恢复引理由叶向上归纳证明正确性。

一次调用处的插值节点数、取值数与组数是 $s$
的多项式，与 $\mathcal L$ 无关。
由 \eqref{eq:tensor-partition-count}，
深度 $d$ 的活动调用的孩子数至多
\[
 s^{O(1)}\exp\left(
       O_C\left(\frac{s}{\sqrt\lambda}
                    \left(\frac{\sqrt2}{8}\right)^d\right)\right).
\]
事实上 $M\le Cs\,8^{-d}$
且 $u_d=\Theta_C(\lambda/2^d)$。
把这些界在 $O_C(\log\lambda)$ 层深度上相乘，
或以此方式界定每层节点数，得总调用数至多
\[
 \exp\left(O_C(s/\sqrt\lambda)
                   +O_C(\lambda\log\lambda)\right)
       =\exp(O_C(s/\sqrt\lambda)).
\]
在叶处，直接求和至多用 $2^H$
个赋值，每赋值多项式次分量求值。
其成本被 $s+\mathcal L$
的一个多项式乘以
\[
 \exp(O_C(H)+O_{C,A}(s^{9/10}))
   =\exp(O_{C,A}(s/\sqrt\lambda))
\]
界定。局部算术、装箱、描述与恢复对每个生成的项
花费多项式时间（已算界的分量求值除外）。
依赖描述的因子界定每次调用的工作。
由于分支界与 $\mathcal L$ 无关，
总工作保持 $s+\mathcal L$
的固定多项式因子。这证明时间界。

对空间，深度优先遍历递归，
一次累加一个关于 $\beta$ 的和与一个关于组的积。
存储一次调用的恢复所需的多项式个 Taylor 数据、
其当前插值与模下标、其图与新张量描述。
其规模在 $s+\mathcal L$ 上是多项式的
（精度由 \eqref{eq:tensor-precision-bound}
计入）。栈层只有 $O_C(\log\lambda)$ 层。
引理~\ref{lem:tapered-hermite}
中的公分母计算也只用多项式位整数。
每次调用按其自身的固定精度累加其收缩项。
原始分量例程按假设使用多项式空间，
新例程亦然。这证明空间界。
最后提取 $\mathcal A_{P_0}$
中的基坐标，它们给出 $\mathcal B_{P_0,\ell}$
中的答案。
\end{proof}


\begin{thebibliography}{CHKX06}

\bibitem[ABHS19]{ABHS19}
Amir Abboud, Karl Bringmann, Danny Hermelin, and Dvir Shabtay.
\href{https://doi.org/10.1137/1.9781611975482.3}
{\emph{SETH-Based Lower Bounds for Subset Sum and Bicriteria Path}}.
30th ACM--SIAM Symposium on Discrete Algorithms, 2019, 41--57.

\bibitem[AW14]{AW14}
Amir Abboud and Virginia Vassilevska Williams.
\href{https://doi.org/10.1109/FOCS.2014.53}
{\emph{Popular Conjectures Imply Strong Lower Bounds for Dynamic Problems}}.
55th IEEE Symposium on Foundations of Computer Science, 2014, 434--443.

\bibitem[AGP23]{AGP23}
Josh Alman, Yunfeng Guan, and Ashwin Padaki.
\href{https://arxiv.org/abs/2211.05217}
{\emph{Smaller Low-Depth Circuits for Kronecker Powers}}.
34th ACM--SIAM Symposium on Discrete Algorithms, 2023, 4159--4187.

\bibitem[AW15]{AW15}
Josh Alman and Ryan Williams.
\href{https://arxiv.org/abs/1507.05106}
{\emph{Probabilistic Polynomials and Hamming Nearest Neighbors}}.
56th IEEE Symposium on Foundations of Computer Science, 2015, 136--150.

\bibitem[AW17]{AW17}
Josh Alman and Ryan Williams.
\href{https://arxiv.org/abs/1611.05558}
{\emph{Probabilistic Rank and Matrix Rigidity}}.
49th ACM Symposium on Theory of Computing, 2017, 641--652.

\bibitem[AR94]{AR94}
Noga Alon and Yuval Roichman.
\href{https://doi.org/10.1002/rsa.3240050203}
{\emph{Random Cayley Graphs and Expanders}}.
Random Structures \& Algorithms \textbf{5}(2), 1994, 271--284.

\bibitem[BI16]{BI16}
Arturs Backurs and Piotr Indyk.
\href{https://doi.org/10.1109/FOCS.2016.56}
{\emph{Which Regular Expression Patterns Are Hard to Match?}}
57th IEEE Symposium on Foundations of Computer Science, 2016, 457--466.

\bibitem[BI18]{BackursIndyk18}
Arturs Backurs and Piotr Indyk.
\href{https://doi.org/10.1137/15M1053128}
{\emph{Edit Distance Cannot Be Computed in Strongly Subquadratic Time
(unless SETH is false)}}.
SIAM Journal on Computing \textbf{47}(3), 2018, 1087--1097.
Conference version in STOC 2015.

\bibitem[Bar23]{Bar23}
Boaz Barak.
\href{https://introtcs.org/public/lec_13_Cook_Levin.html}
{\emph{Introduction to Theoretical Computer Science}}.
Online textbook, 2023. Chapter~15, Lemmas~15.9--15.10.
\href{https://github.com/boazbk/tcs/releases/tag/v0.95}{Archived 2023 edition (v0.95)}.

\bibitem[Bar68]{Bareiss68}
Erwin H. Bareiss.
\href{https://doi.org/10.1090/S0025-5718-1968-0226829-0}
{\emph{Sylvester's Identity and Multistep Integer-Preserving Gaussian Elimination}}.
Mathematics of Computation \textbf{22}(103), 1968, 565--578.

\bibitem[BI13]{BeckImp13}
Chris Beck and Russell Impagliazzo.
\href{https://doi.org/10.1145/2488608.2488669}
{\emph{Strong ETH Holds for Regular Resolution}}.
45th ACM Symposium on Theory of Computing, 2013, 487--494.

\bibitem[BT94]{BT94}
Richard Beigel and Jun Tarui.
\href{https://doi.org/10.1007/BF01263423}
{\emph{On ACC}}.
Computational Complexity \textbf{4}(4), 1994, 350--366.

\bibitem[BGKM23]{BGKM23}
Vishwas Bhargava, Sumanta Ghosh, Mrinal Kumar, and Chandra Kanta Mohapatra.
\href{https://doi.org/10.1145/3625226}
{\emph{Fast, Algebraic Multivariate Multipoint Evaluation in Small
Characteristic and Applications}}.
Journal of the ACM \textbf{70}(6), 2023, 42:1--42:46.
Conference version in STOC 2022.

\bibitem[Bj\"o16]{Bjo16}
Andreas Bj\"orklund.
\href{https://doi.org/10.4230/LIPIcs.SWAT.2016.17}
{\emph{Below All Subsets for Some Permutational Counting Problems}}.
15th Scandinavian Symposium and Workshops on Algorithm Theory,
LIPIcs \textbf{53}, 2016, 17:1--17:11.

\bibitem[BKW19]{BKW19}
Andreas Bj\"orklund, Petteri Kaski, and Ryan Williams.
\href{https://doi.org/10.1007/s00453-018-0513-7}
{\emph{Generalized Kakeya Sets for Polynomial Evaluation and Faster
Computation of Fermionants}}.
Algorithmica \textbf{81}, 2019, 4010--4028.
Conference version in IPEC 2017.

\bibitem[Bri14]{Bri14}
Karl Bringmann.
\href{https://doi.org/10.1109/FOCS.2014.76}
{\emph{Why Walking the Dog Takes Time: Fr\'echet Distance Has No Strongly
Subquadratic Algorithms Unless SETH Fails}}.
55th IEEE Symposium on Foundations of Computer Science, 2014, 661--670.

\bibitem[BK15]{BK15}
Karl Bringmann and Marvin K\"unnemann.
\href{https://doi.org/10.1109/FOCS.2015.15}
{\emph{Quadratic Conditional Lower Bounds for String Problems and
Dynamic Time Warping}}.
56th IEEE Symposium on Foundations of Computer Science, 2015, 79--97.

\bibitem[Cal09]{CalabroThesis}
Chris Calabro.
\href{https://cseweb.ucsd.edu/~ccalabro/dissertation_single_spaced.pdf}
{\emph{The Exponential Complexity of Satisfiability Problems}}.
Ph.D. dissertation, University of California, San Diego, 2009.
Chapter~5.

\bibitem[CIP06]{CIP}
Chris Calabro, Russell Impagliazzo, and Ramamohan Paturi.
\href{https://cseweb.ucsd.edu/~paturi/myPapers/pubs/CalabroImpagliazzoPaturi_2006_ccc.pdf}
{\emph{A Duality between Clause Width and Clause Density for SAT}}.
21st IEEE Conference on Computational Complexity, 2006, 252--260.

\bibitem[CGI$^{+}$16]{CGIMPS16}
Marco L. Carmosino, Jiawei Gao, Russell Impagliazzo, Ivan Mihajlin,
Ramamohan Paturi, and Stefan Schneider.
\href{https://doi.org/10.1145/2840728.2840746}
{\emph{Nondeterministic Extensions of the Strong Exponential Time
Hypothesis and Consequences for Non-reducibility}}.
7th Innovations in Theoretical Computer Science Conference, 2016, 261--270.

\bibitem[CCK$^{+}$20]{CCK20}
Parinya Chalermsook, Marek Cygan, Guy Kortsarz, Bundit Laekhanukit,
Pasin Manurangsi, Danupon Nanongkai, and Luca Trevisan.
\href{https://doi.org/10.1137/18M1166869}
{\emph{From Gap-Exponential Time Hypothesis to Fixed Parameter Tractable
Inapproximability: Clique, Dominating Set, and More}}.
SIAM Journal on Computing \textbf{49}(4), 2020, 772--810.
Conference version in FOCS 2017.

\bibitem[CW21]{CW21}
Timothy M. Chan and R. Ryan Williams.
\href{https://doi.org/10.1145/3402926}
{\emph{Deterministic APSP, Orthogonal Vectors, and More:
Quickly Derandomizing Razborov--Smolensky}}.
ACM Transactions on Algorithms \textbf{17}(1), Article~2, 2021, 1--14.
\href{https://people.csail.mit.edu/rrw/derand-apsp-ov-cr-talg.pdf}{Authors' manuscript}, Theorems~1.6 and~2.5.

\bibitem[CHKX06]{CHKX06}
Jianer Chen, Xiuzhen Huang, Iyad A. Kanj, and Ge Xia.
\href{https://doi.org/10.1016/j.jcss.2006.04.007}
{\emph{Strong Computational Lower Bounds via Parameterized Complexity}}.
Journal of Computer and System Sciences \textbf{72}(8), 2006, 1346--1367.

\bibitem[CLWYY25]{CLWYY25}
Zongchen Chen, Aditya Lonkar, Chunyang Wang, Kuan Yang, and Yitong Yin.
\href{https://doi.org/10.1145/3717823.3718163}
{\emph{Counting Random $k$-SAT near the Satisfiability Threshold}}.
57th ACM Symposium on Theory of Computing, 2025, 867--878.
Theorem~1.2 is numbered as in the \href{https://wcysai.com/publication/randomksat/randomkSAT.pdf}{authors' full version}.

\bibitem[Coo71]{Cook71}
Stephen A. Cook.
\href{https://www.cs.utoronto.ca/~sacook/homepage/1971.pdf}
{\emph{The Complexity of Theorem-Proving Procedures}}.
3rd ACM Symposium on Theory of Computing, 1971, 151--158.

\bibitem[Cur18]{Cur18}
Radu Curticapean.
\href{https://doi.org/10.1016/j.ic.2018.02.008}
{\emph{Block Interpolation: A Framework for Tight Exponential-Time
Counting Complexity}}.
Information and Computation \textbf{261}, Part~2, 2018, 265--280.
Conference version in ICALP 2015.

\bibitem[CDL$^{+}$16]{CDL16}
Marek Cygan, Holger Dell, Daniel Lokshtanov, D\'aniel Marx,
Jesper Nederlof, Yoshio Okamoto, Ramamohan Paturi, Saket Saurabh,
and Magnus Wahlstr\"om.
\href{https://doi.org/10.1145/2925416}
{\emph{On Problems as Hard as CNF-SAT}}.
ACM Transactions on Algorithms \textbf{12}(3), Article~41, 2016, 1--24.

\bibitem[DHM$^{+}$14]{DHMTW}
Holger Dell, Thore Husfeldt, D\'aniel Marx, Nina Taslaman, and Martin Wahl\'en.
\href{https://arxiv.org/abs/1206.1775}
{\emph{Exponential Time Complexity of the Permanent and the Tutte Polynomial}}.
ACM Transactions on Algorithms \textbf{10}(4), Article~21, 2014, 1--32.

\bibitem[Din16]{Din16}
Irit Dinur.
\href{https://eccc.weizmann.ac.il/report/2016/128/}
{\emph{Mildly Exponential Reduction from Gap 3SAT to Polynomial-Gap
Label-Cover}}.
Electronic Colloquium on Computational Complexity, Report TR16-128, 2016.

\bibitem[DL20]{DL20}
Zeev Dvir and Allen Liu.
\href{https://theoryofcomputing.org/articles/v016a020/}
{\emph{Fourier and Circulant Matrices are Not Rigid}}.
Theory of Computing \textbf{16}(20), 2020, 1--48.

\bibitem[vzGG13]{vzGG13}
Joachim von zur Gathen and J\"urgen Gerhard.
\href{https://doi.org/10.1017/CBO9781139856065}{\emph{Modern Computer Algebra}}.
3rd edition, Cambridge University Press, 2013.

\bibitem[Gol00]{Gol00}
Oded Goldreich.
\href{https://eccc.weizmann.ac.il/eccc-reports/2000/TR00-090/index.html}
{\emph{Candidate One-Way Functions Based on Expander Graphs}}.
Electronic Colloquium on Computational Complexity, Report TR00-090, 2000.

\bibitem[Gol08]{Gol08}
Oded Goldreich.
\href{https://www.wisdom.weizmann.ac.il/~oded/CC/bc-1.pdf}
{\emph{P, NP, and NP-Completeness: The Basics of Computational Complexity}}.
Author's draft, June 14, 2008, Proposition~4.7.
The Cambridge University Press edition appeared in 2010.

\bibitem[HKZZ19]{HKZZ19}
Thomas Dueholm Hansen, Haim Kaplan, Or Zamir, and Uri Zwick.
\href{https://doi.org/10.1145/3313276.3316359}
{\emph{Faster $k$-SAT Algorithms Using Biased-PPSZ}}.
51st ACM Symposium on Theory of Computing, 2019, 578--589.

\bibitem[Her14]{Hertli14}
Timon Hertli.
\href{https://doi.org/10.1137/120868177}
{\emph{3-SAT Faster and Simpler---Unique-SAT Bounds for PPSZ Hold in General}}.
SIAM Journal on Computing \textbf{43}(2), 2014, 718--729.
Conference version in FOCS 2011.

\bibitem[Hus11]{Hus11}
Thore Husfeldt.
\href{https://doi.org/10.1007/978-3-642-22012-8_3}
{\emph{Invitation to Algorithmic Uses of Inclusion--Exclusion}}.
38th International Colloquium on Automata, Languages, and Programming,
Part II, Lecture Notes in Computer Science \textbf{6756}, 2011, 42--59.

\bibitem[IMP12]{IMP12}
Russell Impagliazzo, William Matthews, and Ramamohan Paturi.
\href{https://cseweb.ucsd.edu/~paturi/myPapers/pubs/ImpagliazzoMatthewsPaturi_2012_soda.pdf}
{\emph{A Satisfiability Algorithm for $AC^0$}}.
23rd ACM--SIAM Symposium on Discrete Algorithms, 2012, 961--972.
\href{https://doi.org/10.1137/1.9781611973099.77}{Publisher's version}.

\bibitem[IP01]{IP01}
Russell Impagliazzo and Ramamohan Paturi.
\href{https://doi.org/10.1006/jcss.2000.1727}
{\emph{On the Complexity of $k$-SAT}}.
Journal of Computer and System Sciences \textbf{62}(2), 2001, 367--375.

\bibitem[IPZ01]{IPZ}
Russell Impagliazzo, Ramamohan Paturi, and Francis Zane.
\href{https://doi.org/10.1006/jcss.2001.1774}
{\emph{Which Problems Have Strongly Exponential Complexity?}}
Journal of Computer and System Sciences \textbf{63}(4), 2001, 512--530.

\bibitem[JMV13]{JMV13}
Hamid Jahanjou, Eric Miles, and Emanuele Viola.
\href{https://eccc.weizmann.ac.il/report/2013/099/revision/2/download/}
{\emph{Local Reductions}}.
Electronic Colloquium on Computational Complexity, Report TR13-099,
revision 2, 2013. Conference version in ICALP 2015, 749--760.

\bibitem[JVV86]{JVV86}
Mark R. Jerrum, Leslie G. Valiant, and Vijay V. Vazirani.
\href{https://doi.org/10.1016/0304-3975(86)90174-X}
{\emph{Random Generation of Combinatorial Structures from a Uniform Distribution}}.
Theoretical Computer Science \textbf{43}(2--3), 1986, 169--188.

\bibitem[JC26]{JC26}
Tao Jiang and Shaowei Cai.
\href{https://arxiv.org/abs/2607.10697}
{\emph{A Better Analysis For PPSZ For 3-SAT}}.
arXiv:2607.10697, 2026.

\bibitem[Koi06]{Koi06}
Mikko Koivisto.
\href{https://doi.org/10.1016/j.ipl.2005.11.013}
{\emph{Optimal 2-Constraint Satisfaction via Sum--Product Algorithms}}.
Information Processing Letters \textbf{98}(1), 2006, 24--28.

\bibitem[KSY14]{KSY14}
Swastik Kopparty, Shubhangi Saraf, and Sergey Yekhanin.
\href{https://doi.org/10.1145/2629416}
{\emph{High-Rate Codes with Sublinear-Time Decoding}}.
Journal of the ACM \textbf{61}(5), 2014, 28:1--28:20.

\bibitem[KW16]{KW16}
Ioannis Koutis and Ryan Williams.
\href{https://doi.org/10.1145/2742544}
{\emph{Algebraic Fingerprints for Faster Algorithms}}.
Communications of the ACM \textbf{59}(1), 2016, 98--105.

\bibitem[Kut26]{Kutzkov26}
Konstantin Kutzkov.
\href{https://doi.org/10.24963/ijcai.2026/254}
{\emph{Complementary Branching for \#SAT: Bounded-Degree Bounds and Faster \#3-SAT}}.
35th International Joint Conference on Artificial Intelligence, 2026, 2282--2290.

\bibitem[Li21]{Li21}
Ray Li.
\href{https://arxiv.org/abs/2008.05106}
{\emph{Settling SETH vs. Approximate Sparse Directed Unweighted Diameter
(up to (NU)NSETH)}}.
53rd ACM Symposium on Theory of Computing, 2021, 1684--1696.

\bibitem[Liu18]{Liu18}
Sixue Liu.
\href{https://doi.org/10.4230/LIPIcs.ICALP.2018.88}
{\emph{Chain, Generalization of Covering Code, and Deterministic Algorithm for $k$-SAT}}.
45th International Colloquium on Automata, Languages, and Programming,
LIPIcs \textbf{107}, 2018, 88:1--88:13.

\bibitem[LMS18]{LMS18}
Daniel Lokshtanov, D\'aniel Marx, and Saket Saurabh.
\href{https://doi.org/10.1145/3170442}
{\emph{Known Algorithms on Graphs of Bounded Treewidth Are Probably Optimal}}.
ACM Transactions on Algorithms \textbf{14}(2), Article~13, 2018, 1--30.
Conference version in SODA 2011.

\bibitem[MR17]{MR17}
Pasin Manurangsi and Prasad Raghavendra.
\href{https://doi.org/10.4230/LIPIcs.ICALP.2017.78}
{\emph{A Birthday Repetition Theorem and Complexity of Approximating
Dense CSPs}}.
44th International Colloquium on Automata, Languages, and Programming,
LIPIcs \textbf{80}, 2017, 78:1--78:15.

\bibitem[Mar10]{Marx10}
D\'aniel Marx.
\href{https://doi.org/10.4086/toc.2010.v006a005}
{\emph{Can You Beat Treewidth?}}
Theory of Computing \textbf{6}, 2010, 85--112.

\bibitem[NN93]{NN93}
Joseph Naor and Moni Naor.
\href{https://doi.org/10.1137/0222053}
{\emph{Small-Bias Probability Spaces: Efficient Constructions and Applications}}.
SIAM Journal on Computing \textbf{22}(4), 1993, 838--856.

\bibitem[Ope26]{OpenAIUGC26}
OpenAI, \emph{The Unique Games Theorem},
OpenAI Math Release preprint
\href{https://github.com/openai/math/blob/main/preprints/The-Unique-Games-Theorem-September-23-2026/paper.pdf}{OAI:The-Unique-Games-Theorem-September-23-2026},
2026.

\bibitem[PPSZ05]{PPSZ05}
Ramamohan Paturi, Pavel Pudl\'ak, Michael E. Saks, and Francis Zane.
\href{https://doi.org/10.1145/1066100.1066101}
{\emph{An Improved Exponential-Time Algorithm for $k$-SAT}}.
Journal of the ACM \textbf{52}(3), 2005, 337--364.
Conference version in FOCS 1998.

\bibitem[Rag88]{Raghavan88}
Prabhakar Raghavan.
\href{https://doi.org/10.1016/0022-0000(88)90003-7}
{\emph{Probabilistic Construction of Deterministic Algorithms:
Approximating Packing Integer Programs}}.
Journal of Computer and System Sciences \textbf{37}(2), 1988, 130--143.

\bibitem[Rag08]{Rag08}
Prasad Raghavendra.
\href{https://doi.org/10.1145/1374376.1374414}
{\emph{Optimal Algorithms and Inapproximability Results for Every CSP?}}
40th ACM Symposium on Theory of Computing, 2008, 245--254.

\bibitem[RV13]{RV13}
Liam Roditty and Virginia Vassilevska Williams.
\href{https://doi.org/10.1145/2488608.2488673}
{\emph{Fast Approximation Algorithms for the Diameter and Radius of Sparse Graphs}}.
45th ACM Symposium on Theory of Computing, 2013, 515--524.

\bibitem[Sch24]{Scheder24}
Dominik Scheder.
\href{https://doi.org/10.46298/theoretics.24.5}
{\emph{PPSZ is better than you think}}.
TheoretiCS \textbf{3}, Article~5, 2024, 1--37.

\bibitem[Sch99]{Schon99}
Uwe Sch\"oning.
\href{https://doi.org/10.1109/SFFCS.1999.814612}
{\emph{A Probabilistic Algorithm for $k$-SAT and Constraint Satisfaction
Problems}}.
40th IEEE Symposium on Foundations of Computer Science, 1999, 410--414.

\bibitem[SS09]{SS09}
Alexander D. Scott and Gregory B. Sorkin.
\href{https://doi.org/10.1145/1597036.1597049}
{\emph{Polynomial Constraint Satisfaction Problems, Graph Bisection,
and the Ising Partition Function}}.
ACM Transactions on Algorithms \textbf{5}(4), 2009, 45:1--45:27.

\bibitem[SKC94]{SKC94}
Bart Selman, Henry A. Kautz, and Bram Cohen.
\href{https://cdn.aaai.org/AAAI/1994/AAAI94-051.pdf}
{\emph{Noise Strategies for Improving Local Search}}.
12th National Conference on Artificial Intelligence, 1994, 337--343.

\bibitem[SS96]{SS96}
Michael Sipser and Daniel A. Spielman.
\href{https://doi.org/10.1109/18.556667}
{\emph{Expander Codes}}.
IEEE Transactions on Information Theory \textbf{42}(6), 1996, 1710--1722.

\bibitem[Tan81]{Tanner81}
R. Michael Tanner.
\href{https://doi.org/10.1109/TIT.1981.1056404}
{\emph{A Recursive Approach to Low Complexity Codes}}.
IEEE Transactions on Information Theory \textbf{27}(5), 1981, 533--547.

\bibitem[Tov84]{Tov84}
Craig A. Tovey.
\href{https://doi.org/10.1016/0166-218X(84)90081-7}
{\emph{A Simplified NP-Complete Satisfiability Problem}}.
Discrete Applied Mathematics \textbf{8}(1), 1984, 85--89.

\bibitem[Val77]{Val77}
Leslie G. Valiant.
\href{https://doi.org/10.1007/3-540-08353-7_135}
{\emph{Graph-Theoretic Arguments in Low-Level Complexity}}.
Mathematical Foundations of Computer Science 1977,
Lecture Notes in Computer Science \textbf{53}, 1977, 162--176.

\bibitem[Val79]{Val79}
Leslie G. Valiant.
\href{https://doi.org/10.1137/0208032}
{\emph{The Complexity of Enumeration and Reliability Problems}}.
SIAM Journal on Computing \textbf{8}(3), 1979, 410--421.

\bibitem[VW19]{VW19}
Virginia Vassilevska Williams.
\href{https://doi.org/10.1142/9789813272880_0188}
{\emph{On Some Fine-Grained Questions in Algorithms and Complexity}}.
Proceedings of the International Congress of Mathematicians (ICM 2018),
World Scientific, 2019, 3447--3487.

\bibitem[VW21]{VW21}
Nikhil Vyas and Ryan Williams.
\href{https://doi.org/10.1613/jair.1.11859}
{\emph{On Super Strong ETH}}.
Journal of Artificial Intelligence Research \textbf{70}, 2021, 473--495.

\bibitem[Wah08]{Wahlstrom08}
Magnus Wahlstr\"om.
\href{https://doi.org/10.1007/978-3-540-79723-4_19}
{\emph{A Tighter Bound for Counting Max-Weight Solutions to 2SAT Instances}}.
3rd International Workshop on Parameterized and Exact Computation,
Lecture Notes in Computer Science \textbf{5018}, 2008, 202--213.

\bibitem[Wil05]{Wil05}
Ryan Williams.
\href{https://doi.org/10.1016/j.tcs.2005.09.023}
{\emph{A New Algorithm for Optimal 2-Constraint Satisfaction and Its Implications}}.
Theoretical Computer Science \textbf{348}(2--3), 2005, 357--365.

\bibitem[Wil13]{Wil13}
Ryan Williams.
\href{https://doi.org/10.1137/10080703X}
{\emph{Improving Exhaustive Search Implies Superpolynomial Lower Bounds}}.
SIAM Journal on Computing \textbf{42}(3), 2013, 1218--1244.
Conference version in STOC 2010. The numbering of Lemma~3.2 and Theorem~6.1 follows the
\href{https://people.csail.mit.edu/rrw/improved-algs-lbs2.pdf}{author's full version}.

\bibitem[Wil14]{Wil14}
Richard Ryan Williams.
\href{https://doi.org/10.4230/LIPIcs.FSTTCS.2014.47}
{\emph{The Polynomial Method in Circuit Complexity Applied to Algorithm
Design (Invited Talk)}}.
34th International Conference on Foundation of Software Technology and
Theoretical Computer Science, LIPIcs \textbf{29}, 2014, 47--60.

\bibitem[Wil16]{Wil16}
Richard Ryan Williams.
\href{https://doi.org/10.4230/LIPIcs.CCC.2016.2}
{\emph{Strong ETH Breaks With Merlin and Arthur: Short Non-Interactive
Proofs of Batch Evaluation}}.
31st Conference on Computational Complexity, LIPIcs \textbf{50}, 2016,
2:1--2:17.

\bibitem[WY07]{WY07}
David Woodruff and Sergey Yekhanin.
\href{https://doi.org/10.1137/06065773X}
{\emph{A Geometric Approach to Information-Theoretic Private Information
Retrieval}}.
SIAM Journal on Computing \textbf{37}(4), 2007, 1046--1056.

\bibitem[\v{Z}\'ak83]{Zak83}
Stanislav \v{Z}\'ak.
\href{https://doi.org/10.1016/0304-3975(83)90015-4}{\emph{A Turing
Machine Time Hierarchy}}.
Theoretical Computer Science \textbf{26}(3), 1983, 327--333.

\end{thebibliography}
\end{document}
